% Springer-compatible LaTeX template for APJM / MIR submission
% Compiled from Markdown via pandoc.
% Document class: `article` (Springer Nature's sn-jnl is not bundled; this template
% follows Springer's general manuscript guidelines and is compatible with both
% Springer Nature LaTeX template and Word-track submission portals).

\documentclass[12pt,a4paper]{article}

% --- Geometry & spacing ---
\usepackage[margin=2.5cm]{geometry}
\usepackage{setspace}
\onehalfspacing

% --- Encoding & fonts ---
% Conditional setup: under pdflatex use legacy inputenc; under xelatex/lualatex
% use fontspec for native Unicode support (needed for β, −, × symbols).
\usepackage{iftex}
\ifPDFTeX
  \usepackage[utf8]{inputenc}
  \usepackage[T1]{fontenc}
  \usepackage{lmodern}
\else
  \usepackage{fontspec}
  \setmainfont{Liberation Serif}[
    Scale=1.0,
    BoldFont={Liberation Serif Bold},
    ItalicFont={Liberation Serif Italic},
    BoldItalicFont={Liberation Serif Bold Italic}
  ]
  % Liberation Mono carries the full Vietnamese diacritic set; DejaVu Sans
  % Mono does not (renders Ấ/ầ/ế/ố as missing characters in code blocks).
  \setmonofont{Liberation Mono}[Scale=0.85]
  % Map common Unicode math/typographic symbols that may appear in markdown
  % body text. Liberation Serif lacks them, so route via amssymb/math mode.
  % (We deliberately avoid unicode-math, which conflicts with the amssymb
  % package that pandoc auto-loads in the preamble.)
  \usepackage{newunicodechar}
  \newunicodechar{∈}{\ensuremath{\in}}
  \newunicodechar{∉}{\ensuremath{\notin}}
  \newunicodechar{≤}{\ensuremath{\leq}}
  \newunicodechar{≥}{\ensuremath{\geq}}
  \newunicodechar{≠}{\ensuremath{\neq}}
  \newunicodechar{≈}{\ensuremath{\approx}}
  \newunicodechar{∞}{\ensuremath{\infty}}
  \newunicodechar{∑}{\ensuremath{\sum}}
  \newunicodechar{∏}{\ensuremath{\prod}}
  \newunicodechar{∗}{\ensuremath{\ast}}
  \newunicodechar{·}{\ensuremath{\cdot}}
  \newunicodechar{−}{\ensuremath{-}}
  \newunicodechar{×}{\ensuremath{\times}}
  \newunicodechar{÷}{\ensuremath{\div}}
  \newunicodechar{±}{\ensuremath{\pm}}
  % Typographic checkmarks/crosses (used as bullet markers in markdown).
  \newunicodechar{✓}{\ensuremath{\checkmark}}
  \newunicodechar{✗}{\ensuremath{\times}}
  \newunicodechar{⚠}{\textbf{!}}
  \newunicodechar{→}{\ensuremath{\rightarrow}}
  \newunicodechar{←}{\ensuremath{\leftarrow}}
  \newunicodechar{⇒}{\ensuremath{\Rightarrow}}
  % Allow URL line breaks at slashes (fix Overfull \hbox in long URLs)
  \usepackage{xurl}
\fi

% --- Math ---
\usepackage{amsmath,amssymb,amsfonts,mathtools}
\usepackage{bm}

% --- Tables ---
\usepackage{booktabs}
\usepackage{longtable}
\usepackage{multirow}
\usepackage{array}
\usepackage{tabularx}

% --- Figures ---
\usepackage{graphicx}
\usepackage{float}
\usepackage{caption}
\captionsetup{font=small,labelfont=bf}

% --- Misc text ---
\usepackage{microtype}
\usepackage{enumitem}
\setlist{nosep,leftmargin=1.5em}
\usepackage{textcomp}
\usepackage{xcolor}
\usepackage{ulem}\normalem

% --- Code block support (pandoc Shaded / Highlighting environments) ---
% Required when source markdown contains fenced code blocks (```...```).
\usepackage{fancyvrb}
\usepackage{framed}
\definecolor{shadecolor}{RGB}{248,248,248}
\newenvironment{Shaded}{\begin{snugshade}}{\end{snugshade}}
% Highlighting environment for syntax-highlighted code
\providecommand{\BuiltInTok}[1]{#1}
\providecommand{\CommentTok}[1]{\textit{#1}}
\providecommand{\KeywordTok}[1]{\textbf{#1}}
\providecommand{\NormalTok}[1]{#1}
\providecommand{\StringTok}[1]{\textit{#1}}
\providecommand{\OperatorTok}[1]{#1}
\providecommand{\DataTypeTok}[1]{#1}
\providecommand{\FunctionTok}[1]{#1}

% --- Links & references ---
\usepackage[hidelinks,breaklinks=true]{hyperref}
\usepackage{cleveref}

% --- Bibliography (APA7 via natbib + apalike) ---
% APJM/MIR accept APA7. natbib's apalike close approximation; for full APA7
% compliance use biblatex with apa style.
\usepackage[round,sort]{natbib}
\bibliographystyle{apalike}

% --- Pandoc compatibility shims ---
\providecommand{\tightlist}{\setlength{\itemsep}{0pt}\setlength{\parskip}{0pt}}
\providecommand{\passthrough}[1]{\lstset{mathescape=false}#1\lstset{mathescape=true}}

% --- Title page (filled in per paper) ---
\title{Chuyên đề tiến sĩ số 1 --- Thực trạng hiệu quả hoạt động kinh
doanh ở Châu Á}
\author{Author details withheld for blind review}
\date{2026-05-21}

\begin{document}

\maketitle

\hypertarget{chuyuxean-ux111ux1ec1-tiux1ebfn-sux129-sux1ed1-1--bux1ea3n-nhuxe1p-ux111ux1ea7y-ux111ux1ee7-phux1ea7n-1-buxeca-tuxf3m-tux1eaft-chux1b0ux1a1ng-13}{%
\section{CHUYÊN ĐỀ TIẾN SĨ SỐ 1 --- BẢN NHÁP ĐẦY ĐỦ (PHẦN 1: BÌA, TÓM
TẮT, CHƯƠNG
1--3)}\label{chuyuxean-ux111ux1ec1-tiux1ebfn-sux129-sux1ed1-1--bux1ea3n-nhuxe1p-ux111ux1ea7y-ux111ux1ee7-phux1ea7n-1-buxeca-tuxf3m-tux1eaft-chux1b0ux1a1ng-13}}

\begin{quote}
Căn cứ Quyết định số 4768/QĐ-ĐHCT ngày 15/10/2024 (giao chuyên đề);
Quyết định số 4769/QĐ-ĐHCT ngày 15/10/2024 (tên luận án --- "Châu Á").
Outline tham chiếu: \texttt{thesis/12\_chuyen\_de\_1\_outline\_vi.md}.
Phần 2: \texttt{thesis/15\_cd1\_part2\_findings\_vi.md} · Phần 3:
\texttt{thesis/16\_cd1\_part3\_cases\_conclusion\_vi.md}. Bảng thuật ngữ
Anh-Việt: \texttt{thesis/09b\_vn\_term\_glossary.md} · Chuẩn văn phong:
\texttt{thesis/09\_academic\_writing\_standards\_vi.md}. Hình minh họa:
\texttt{thesis/figures/} (11 hình).
\end{quote}

\begin{center}\rule{0.5\linewidth}{0.5pt}\end{center}

\hypertarget{trang-buxeca}{%
\subsection{TRANG BÌA}\label{trang-buxeca}}

\begin{verbatim}
BỘ GIÁO DỤC VÀ ĐÀO TẠO
TRƯỜNG ĐẠI HỌC CẦN THƠ
TRƯỜNG KINH TẾ

─────────────

ĐỖ THÙY HƯƠNG

CHUYÊN ĐỀ TIẾN SĨ SỐ 1

THỰC TRẠNG VỀ HIỆU QUẢ HOẠT ĐỘNG KINH DOANH
CỦA CÁC DOANH NGHIỆP Ở CHÂU Á

Ngành: Quản trị kinh doanh
Mã ngành: 9340101
Mã nghiên cứu sinh: P1323001

NGƯỜI HƯỚNG DẪN KHOA HỌC
TS. NGUYỄN MINH CẢNH

Cần Thơ, năm 2026
\end{verbatim}

\begin{center}\rule{0.5\linewidth}{0.5pt}\end{center}

\hypertarget{lux1eddi-cam-ux111oan}{%
\subsection{LỜI CAM ĐOAN}\label{lux1eddi-cam-ux111oan}}

Tôi xin cam đoan chuyên đề tiến sĩ số 1 với tiêu đề "Thực trạng về hiệu
quả hoạt động kinh doanh của các doanh nghiệp ở Châu Á" là công trình
nghiên cứu của riêng tôi dưới sự hướng dẫn khoa học của TS. Nguyễn Minh
Cảnh, Trường Đại học Cần Thơ. Các kết quả phân tích, số liệu, bảng biểu,
hình vẽ trình bày trong chuyên đề đều trung thực, có nguồn gốc rõ ràng.
Tài liệu tham khảo được trích dẫn theo chuẩn APA 7th. Những trích dẫn từ
các nghiên cứu khác đều được nêu rõ trong danh mục tài liệu tham khảo.
Tôi xin chịu trách nhiệm hoàn toàn về nội dung khoa học của chuyên đề
này.

\emph{Cần Thơ, ngày \ldots{} tháng \ldots{} năm 2026}

\emph{Nghiên cứu sinh}

\textbf{Đỗ Thùy Hương}

\begin{center}\rule{0.5\linewidth}{0.5pt}\end{center}

\hypertarget{tuxf3m-tux1eaft}{%
\subsection{TÓM TẮT}\label{tuxf3m-tux1eaft}}

\textbf{Mục đích.} Chuyên đề thực hiện phân tích mô tả-chẩn đoán thực
trạng hiệu quả hoạt động kinh doanh của doanh nghiệp ở các nền kinh tế
châu Á trong giai đoạn 2009-2025, đồng thời mở rộng so sánh với các Quốc
đảo nhỏ Thái Bình Dương như trường hợp biên. Mục tiêu là nhận diện các
mẫu hình điển hình về năng suất, lợi nhuận, tăng trưởng và đổi mới qua
sáu phân nhóm chế độ thể chế, từ đó định hướng giả thuyết cho mô hình
thực nghiệm trong Chuyên đề tiến sĩ số 2.

\textbf{Thiết kế / phương pháp luận / cách tiếp cận.} Chuyên đề sử dụng
dữ liệu vi mô gộp gồm 101.185 doanh nghiệp, 108 cặp quốc gia-năm và 14
mốc khảo sát từ Khảo sát Doanh nghiệp của Ngân hàng Thế giới, bao trùm
47 nền kinh tế phân loại thành sáu phân nhóm theo khung Biến thiên Chế
độ Thể chế (Institutional Context Regime Variation, ICRV): Tiên tiến đổi
mới sáng tạo dẫn dắt, Tiên tiến tài nguyên dẫn dắt, Thu nhập trung bình
cao, Mới nổi, Cận biên, và Quốc đảo nhỏ Thái Bình Dương. Phương pháp chủ
đạo là thống kê mô tả đa chiều kết hợp đồ thị, bảng so sánh xuyên quốc
gia-xuyên ngành và phân tích hai biến trên bốn chiều hiệu quả: năng suất
lao động, lợi nhuận và biên lợi nhuận, tăng trưởng, đổi mới và cấu trúc
doanh nghiệp.

\textbf{Kết quả nghiên cứu.} Hiệu quả doanh nghiệp châu Á có phân tán
lớn và không hội tụ qua giai đoạn 2009-2025. Bảy tiểu cảnh điển hình tồn
tại đồng thời, gồm Singapore, Vùng Vịnh (Saudi Arabia, Qatar, Kuwait),
Việt Nam, Trung Quốc, Đông Á-Đông Nam Á mới nổi rộng hơn, Mông Cổ trong
phạm vi chính, cùng các Quốc đảo nhỏ Thái Bình Dương như trường hợp
biên, với các điểm mạnh hiệu quả khác biệt. Tương quan giữa các yếu tố
giải thích sơ bộ và hiệu quả thay đổi dấu và cường độ giữa các phân nhóm
thể chế, đặc biệt bốn lần đảo dấu của hệ số FDI/FSTS/TCI/DAI xuyên năm
phân nhóm chính, gợi ý sự cần thiết của các mô hình phi tuyến và điều
tiết đa tầng. Phát hiện hai phân nhóm con trong chế độ Tiên tiến (đổi
mới sáng tạo so với tài nguyên dẫn dắt) và xác nhận mẫu hình chi phí
buộc phải quốc tế hoá ở các Quốc đảo Thái Bình Dương.

\textbf{Tính mới / đóng góp.} Chuyên đề đóng góp ba điểm cho tài liệu
kinh doanh quốc tế ở thị trường mới nổi: thứ nhất, mở rộng phạm vi mẫu
so với các tổng hợp châu Á trước đây lên 47 nền kinh tế và 14 đợt khảo
sát; thứ hai, phân tách phân nhóm Tiên tiến theo cơ chế tăng trưởng, một
đóng góp lý thuyết kiểu Varieties of Capitalism cho bối cảnh châu Á; thứ
ba, thiết lập trường hợp biên Thái Bình Dương cho lý thuyết chữ U ngược
Lu-Beamish, định hướng giả thuyết H6 cho Chuyên đề tiến sĩ số 2.

\textbf{Từ khoá:} hiệu quả doanh nghiệp; châu Á; Khảo sát Doanh nghiệp
của Ngân hàng Thế giới; năng suất lao động; phân nhóm con Tiên tiến;
trường hợp biên Thái Bình Dương; Biến thiên Chế độ Thể chế.

\textbf{Phân loại JEL:} F23; O53; L25; D22; O19.

\textbf{Loại bài:} Chuyên đề tiến sĩ.

\begin{center}\rule{0.5\linewidth}{0.5pt}\end{center}

\hypertarget{abstract}{%
\subsection{ABSTRACT}\label{abstract}}

\textbf{Purpose.} This specialty essay conducts a descriptive-diagnostic
analysis of the firm-performance landscape across Asian economies during
2009-2025, with comparative extension to Pacific Small Island Developing
States as a boundary case. The aim is to identify stylised patterns in
productivity, profitability, growth, and innovation across six
institutional sub-regimes, thereby orienting hypothesis development for
the empirical model in Specialty Essay 2.

\textbf{Design/methodology/approach.} The study uses a pooled microdata
set of 101,185 firms, 108 country-year pairs, and 14 survey waves from
the World Bank Enterprise Surveys covering 47 economies, classified by
the Institutional Context Regime Variation (ICRV) framework into six
sub-regimes: Advanced innovation-driven, Advanced resource-driven,
Upper-middle, Emerging, Frontier, and Pacific SIDS. The primary method
combines multidimensional descriptive statistics with charts,
cross-country and cross-industry comparison tables, and bivariate
analyses on four performance dimensions: labour productivity,
profitability and margins, growth, and innovation and structural
composition.

\textbf{Findings.} Firm performance across Asian economies exhibits
substantial and persistent dispersion with no convergence pattern over
2009-2025. Seven salient archetypes coexist with distinct performance
strengths: Singapore, the Gulf trio of Saudi Arabia, Qatar, and Kuwait,
Vietnam, China, broader Emerging East and Southeast Asia, and Mongolia
within the primary scope, together with Pacific SIDS as boundary case.
Bivariate correlations between candidate explanatory factors and
performance vary in sign and magnitude across regimes; in particular,
four sign reversals in the foreign direct investment, export intensity,
technological capability, and digital adoption coefficients across the
five primary sub-regimes suggest the need for nonlinear and multi-level
moderation models. The essay identifies a sub-grouping of the Advanced
regime into innovation-driven and resource-driven types and confirms the
forced-internationalisation-penalty pattern in Pacific SIDS.

\textbf{Originality/value.} The essay contributes to emerging-market
international business research in three ways: it extends the sample
frame of prior Asian pooled studies to 47 economies and 14 survey waves;
it splits the Advanced regime by growth mechanism, a theoretical
contribution in the Varieties-of-Capitalism tradition adapted to Asia;
and it establishes the Pacific boundary case as a stress-test for the
Lu-Beamish inverted-U logic, orienting hypothesis H6 for Specialty Essay
2.

\textbf{Keywords:} firm performance; Asia; World Bank Enterprise
Surveys; labour productivity; Advanced sub-grouping; Pacific boundary
case; Institutional Context Regime Variation.

\textbf{JEL Classification:} F23; O53; L25; D22; O19.

\textbf{Paper type:} Doctoral specialty essay.

\begin{center}\rule{0.5\linewidth}{0.5pt}\end{center}

\hypertarget{mux1ee5c-lux1ee5c}{%
\subsection{MỤC LỤC}\label{mux1ee5c-lux1ee5c}}

(Tự động tạo khi xuất Word/PDF từ markdown, căn cứ heading H2/H3.)

\begin{center}\rule{0.5\linewidth}{0.5pt}\end{center}

\hypertarget{danh-mux1ee5c-bux1ea3ng}{%
\subsection{DANH MỤC BẢNG}\label{danh-mux1ee5c-bux1ea3ng}}

\begin{longtable}[]{@{}lll@{}}
\toprule\noalign{}
Bảng & Tiêu đề & Trang \\
\midrule\noalign{}
\endhead
\bottomrule\noalign{}
\endlastfoot
Bảng 2.6.1 & Ma trận đối chiếu 3 trạng thái U-curve theo 3 điều kiện cấu
trúc & Ch.2 \\
Bảng 3.1 & 47 nền kinh tế phân theo 6 nhóm ICRV & Ch.3 \\
Bảng 3.2 & Đặc điểm các nhóm ICRV (tiêu chí + đại diện + pattern dự
kiến) & Ch.3 \\
Bảng 4.1 & Phân tán năng suất lao động theo phân nhóm con thể chế &
Ch.4 \\
Bảng 4.3 & FSTS, doanh nghiệp xuất khẩu và CAGR việc làm theo phân nhóm
con & Ch.4 \\
Bảng 4.4 & Đổi mới sáng tạo và áp dụng số (\%) & Ch.4 \\
Bảng 4.4.0 & Khung 5 lĩnh vực chỉ số WBES & Ch.4 \\
Bảng 4.5 & Cấu trúc doanh nghiệp theo phân nhóm con (\%) & Ch.4 \\
Bảng 4.5.5 & Bốn rào cản hàng đầu Asia-Pacific WBES & Ch.4 \\
Bảng 4.6 & Thay đổi theo thời gian --- Δ điểm phần trăm 2018--2025 so
với 2009--2012 & Ch.4 \\
Bảng 4.7.5 & Phân biệt 3 chỉ số tham nhũng WBES --- Bribery vs Graft &
Ch.4 \\
Bảng 4.8.1 & Khung phân loại 9 ngành ISIC Rev. 4 cho nhóm dữ liệu WBES &
Ch.4 \\
Bảng 4.9 & Phân nhóm con Emerging, số liệu tổng hợp 2009--2025 & Ch.4 \\
Bảng 4.10 & 14 quốc gia trong đợt khảo sát 2025 & Ch.4 \\
Bảng 5.1 & Ma trận so sánh đa chiều bảy tiểu cảnh điển hình & Ch.5 \\
Bảng 6.1 & Hệ số tương quan Pearson (n=101.185) & Ch.6 \\
Bảng 6.1.1 & Pattern giới tính trong TMT + ownership, phân tầng theo khu
vực & Ch.6 \\
\end{longtable}

\begin{center}\rule{0.5\linewidth}{0.5pt}\end{center}

\hypertarget{danh-mux1ee5c-huxecnh}{%
\subsection{DANH MỤC HÌNH}\label{danh-mux1ee5c-huxecnh}}

\begin{longtable}[]{@{}lll@{}}
\toprule\noalign{}
Hình & Tiêu đề & Trang \\
\midrule\noalign{}
\endhead
\bottomrule\noalign{}
\endlastfoot
Hình 4.0 & Phân bố mẫu 101.185 doanh nghiệp xuyên 17 năm (2009--2025)
theo 3 thế hệ schema WBES & Ch.4 \\
Hình 4.1.1 & Phân bố pool 5 phân nhóm con thể chế (n=101.185 doanh
nghiệp) & Ch.4 \\
Hình 4.2 & Phân phối năng suất theo 6 phân nhóm con (kernel density,
n=101.185) & Ch.4 \\
Hình 4.4 & Slope chart 2009--2012 vs 2018--2025 (Δ điểm phần trăm 5 chỉ
số) & Ch.4 \\
Hình 4.6 & Phân phối quy mô doanh nghiệp theo phân nhóm con (n=101.185)
& Ch.4 \\
Hình 4.7 & Spider chart 5 chiều × 2 mốc thời gian (Advanced, Emerging,
SIDS) & Ch.4 \\
Hình 5.6.1 & Mongolia: Evolution 2009--2025 (4 đợt khảo sát WBES) &
Ch.5 \\
\end{longtable}

\begin{center}\rule{0.5\linewidth}{0.5pt}\end{center}

\hypertarget{danh-mux1ee5c-tux1eeb-viux1ebft-tux1eaft}{%
\subsection{DANH MỤC TỪ VIẾT
TẮT}\label{danh-mux1ee5c-tux1eeb-viux1ebft-tux1eaft}}

\begin{longtable}[]{@{}lll@{}}
\toprule\noalign{}
Từ viết tắt & Tiếng Anh & Tiếng Việt \\
\midrule\noalign{}
\endhead
\bottomrule\noalign{}
\endlastfoot
ADB & Asian Development Bank & Ngân hàng Phát triển Châu Á \\
AIPI & AI Productivity Impact & Tác động năng suất của AI \\
ASEAN & Association of Southeast Asian Nations & Hiệp hội các quốc gia
Đông Nam Á \\
BEE & Business Environment and Enterprise & --- \\
BREADY & WBES 2018+ schema & Thế hệ bảng hỏi WBES mới từ 2018 \\
CAGR & Compound Annual Growth Rate & Tốc độ tăng trưởng hàng năm kép \\
CDCM & Context-Contingent Digital Capability Model & Mô hình năng lực số
phụ thuộc bối cảnh \\
CĐ1 & Chuyên đề tiến sĩ số 1 & Specialty Essay 1 \\
CĐ2 & Chuyên đề tiến sĩ số 2 & Specialty Essay 2 \\
CIEM & Central Institute for Economic Management & Viện Nghiên cứu Quản
lý Kinh tế Trung ương \\
CRM & Customer Relationship Management & Quản lý quan hệ khách hàng \\
DAI & Digital Adoption Index & Chỉ số áp dụng số \\
DAI (Tier-1) & Website binary & Sự hiện diện số cơ bản \\
DAI (Tier-2) & Website + e-payment composite & Chuyển đổi số cơ bản \\
DPL & Digital Paradox Lifecycle & Vòng đời nghịch lý số \\
EAP & East Asia and Pacific & Đông Á và Thái Bình Dương \\
ECA & Europe and Central Asia & Châu Âu và Trung Á \\
ERP & Enterprise Resource Planning & Hoạch định nguồn lực doanh
nghiệp \\
FDA & Foundational Digital Adoption & Áp dụng số nền tảng \\
FDI & Foreign Direct Investment & Đầu tư trực tiếp nước ngoài \\
FSTS & Foreign Trade Sales to Total Sales & Tỷ lệ doanh thu xuất khẩu
trên tổng doanh thu \\
GCC & Gulf Cooperation Council & Hội đồng Hợp tác vùng Vịnh \\
GDP & Gross Domestic Product & Tổng sản phẩm quốc nội \\
GNI & Gross National Income & Tổng thu nhập quốc gia \\
GVC & Global Value Chain & Chuỗi giá trị toàn cầu \\
HC1 & HC1 Robust Standard Errors & Sai số chuẩn vững HC1 \\
ICRV & Institutional Context Regime Variation & Biến thiên chế độ thể
chế \\
IFC & International Finance Corporation & Tổ chức Tài chính Quốc tế \\
IMF & International Monetary Fund & Quỹ Tiền tệ Quốc tế \\
IPR & Intellectual Property Rights & Quyền sở hữu trí tuệ \\
ISIC & International Standard Industrial Classification & Phân loại
ngành kinh tế quốc tế \\
ISO & International Organization for Standardization & Tổ chức Tiêu
chuẩn Quốc tế \\
IV & Instrumental Variable & Biến công cụ \\
LP & Labor Productivity & Năng suất lao động \\
MENA & Middle East and North Africa & Trung Đông và Bắc Phi \\
MIRAB & Migration, Remittances, Aid, Bureaucracy & Mô hình kinh tế SIDS
Pacific \\
MNE & Multinational Enterprise & Doanh nghiệp đa quốc gia \\
ODA & Official Development Assistance & Hỗ trợ phát triển chính thức \\
OLI & Ownership-Location-Internalization & Mô hình chiết trung OLI \\
PICS3 & WBES 2009--2012 schema & Thế hệ bảng hỏi WBES 2009--2012 \\
PICs & Pacific Island Countries & Các quốc đảo Thái Bình Dương \\
PLMS & Pacific Labour Mobility Scheme & Chương trình lao động di chuyển
Thái Bình Dương \\
PPP & Purchasing Power Parity & Ngang giá sức mua \\
PRC & People\textquotesingle s Republic of China & Cộng hòa Nhân dân
Trung Hoa \\
R\&D & Research and Development & Nghiên cứu và Phát triển \\
RBV & Resource-Based View & Quan điểm dựa trên nguồn lực \\
RSE & Recognised Seasonal Employer & --- \\
SIDS & Small Island Developing States & Quốc đảo nhỏ đang phát triển \\
SME & Small and Medium Enterprise & Doanh nghiệp vừa và nhỏ \\
TCE & Transaction Cost Economics & Lý thuyết chi phí giao dịch \\
TCI & Technological Capability Index & Chỉ số năng lực công nghệ \\
TMT & Top Management Team & Nhóm quản trị cấp cao \\
UNCTAD & United Nations Conference on Trade and Development & --- \\
WBES & World Bank Enterprise Surveys & Điều tra doanh nghiệp của Ngân
hàng Thế giới \\
WGI & World Governance Indicators & Chỉ số quản trị quốc gia \\
\end{longtable}

\begin{center}\rule{0.5\linewidth}{0.5pt}\end{center}

\begin{center}\rule{0.5\linewidth}{0.5pt}\end{center}

\hypertarget{chux1b0ux1a1ng-1--giux1edbi-thiux1ec7u}{%
\subsection{CHƯƠNG 1 --- GIỚI
THIỆU}\label{chux1b0ux1a1ng-1--giux1edbi-thiux1ec7u}}

\hypertarget{11-ux111ux1eb7t-vux1ea5n-ux111ux1ec1}{%
\subsubsection{1.1 Đặt vấn
đề}\label{11-ux111ux1eb7t-vux1ea5n-ux111ux1ec1}}

Khu vực châu Á đã trở thành \textbf{động lực} tăng trưởng của kinh tế
thế giới trong gần hai thập niên qua. Theo Asian Development Bank (ADB,
2024), châu Á đóng góp khoảng 60\% tăng trưởng GDP toàn cầu giai đoạn
2010--2023, với tổng quy mô kinh tế chiếm xấp xỉ 40\% GDP toàn cầu
(World Bank, 2024). Bên trong khu vực, hiệu quả hoạt động của doanh
nghiệp -- đơn vị tế bào của nền kinh tế -- là yếu tố quyết định mức độ
chuyển hoá tăng trưởng quốc gia thành \textbf{thịnh vượng} và năng lực
cạnh tranh dài hạn (Cusolito \& Maloney, 2018). Vì vậy, đánh giá thực
trạng hiệu quả hoạt động kinh doanh của doanh nghiệp ở châu Á có ý nghĩa
lý luận lẫn chính sách.

Bức tranh hiệu quả doanh nghiệp châu Á 2009--2025 được định hình bởi
\textbf{bảy lớp bối cảnh đan xen}. Bốn lớp đầu gồm: hậu khủng hoảng tài
chính toàn cầu 2008--2009; tái cấu trúc chuỗi giá trị toàn cầu sau chiến
tranh thương mại Mỹ -- Trung 2018 (UNCTAD, 2023); đại dịch COVID-19
2020--2022 (World Bank, 2023); và làn sóng chuyển đổi số tăng tốc từ
2018 (Verhoef et al., 2021; Banalieva \& Dhanaraj, 2019). Lớp thứ năm là
\textbf{giai đoạn AI bùng nổ và củng cố hậu COVID 2023--2025} (Stallkamp
\& Schotter, 2021). \textbf{Lớp thứ sáu là xung đột Trung Đông 2026} ---
IMF (2026, April) điều chỉnh giảm tăng trưởng toàn cầu xuống 3,1\% năm
2026 và giảm dự báo Saudi Arabia mạnh từ \textasciitilde4,5\% xuống
3,1\% do giảm sản lượng dầu; ADB (2026, April) trong báo cáo \emph{Asian
Development Outlook April 2026: The Middle East Conflict Challenges
Resilience in Asia and the Pacific} dự báo Pacific 3,4\%, Đông Nam Á
4,7\%, Việt Nam 7,0\%. \textbf{Riêng Việt Nam, ADB (2026, April, Vietnam
Economic Outlook) ghi nhận GDP lịch sử 2025 đạt 8,0\%}, mức cao nhất khu
vực Đông Nam Á, và cập nhật dự báo \textbf{7,2\% (2026), sau đó 7,0\%
(2027)} (đường giảm 8,0\%, sau đó 7,2\%, sau đó 7,0\%, vượt mức bình
quân ASEAN 4,6\% và Đông Nam Á 4,7\%), với động lực từ FDI commitment
\$2,4 tỷ qua 18 dự án chiến lược (fintech, năng lượng sạch, logistics),
năng suất lao động dự phóng +5,1\% giai đoạn 2026-2027 và FDI
manufacturing tiến vào "higher-value" stage. \textbf{Lớp thứ bảy là bất
định chính sách thuế quan Mỹ 2026} --- Mỹ chiếm xấp xỉ \textbf{30\% tổng
kim ngạch xuất khẩu Việt Nam} (đối tác thương mại lớn nhất); thuế toàn
cầu 10\% tạm thời áp dụng đầu 2026 (nguy cơ tăng lên 15\% trong nửa cuối
năm) tạo áp lực kép lên doanh nghiệp xuất khẩu châu Á: (a)
\textbf{front-loading effect} --- đẩy nhanh giao hàng ngắn hạn nhằm tận
dụng cửa sổ thuế thấp, boost doanh thu 2025 (8,0\%) nhưng làm giảm đà
tăng trưởng 2026-2027; (b) \textbf{trì hoãn đầu tư} do bất định chính
sách (policy uncertainty) dẫn đến cắt giảm CapEx + R\&D; (c)
\textbf{phân hóa giữa tập đoàn lớn và SME} --- tập đoàn lớn có khả năng
hấp thụ tariff cost cao hơn (đa thị trường, hedging tài chính, chuỗi
cung ứng đa lớp); SME chịu áp lực biên lợi nhuận và phải lựa chọn giữa
hấp thụ cost vs chuyển giá khách hàng (ADB Vietnam Economic Outlook
2026). Bảy lớp bối cảnh này tác động không đồng đều đến các nhóm nền
kinh tế châu Á.

Trong bối cảnh đó, hệ thống thông tin về hiệu quả doanh nghiệp châu Á
còn phân mảnh. World Bank Enterprise Surveys (WBES) là cơ sở dữ liệu vi
mô tin cậy và có tính so sánh xuyên quốc gia tốt nhất hiện có (World
Bank Enterprise Surveys, 2025). Đỗ và Phan (2026 --- VEFR) đã mở rộng
phạm vi WBES đến 17 nền kinh tế châu Á mới nổi với \textasciitilde40.633
doanh nghiệp, ghi nhận pattern "\textbf{hiệu ứng lá chắn số (digital
shield effect)}". Tuy nhiên, vẫn còn ít có tổng hợp xuyên 47 nền kinh tế
trên giai đoạn dài 2009--2025 sau khi hài hòa ba thế hệ khung dữ liệu
(PICS3, WBES Standardized, BREADY/BEE).

Vì vậy, chuyên đề này lấp khoảng trống thực tiễn bằng việc cập nhật
\textbf{pattern điển hình (stylized facts)} đa chiều về hiệu quả doanh
nghiệp châu Á. Tổng hợp dữ liệu (pool) \textbf{101.185 doanh nghiệp
xuyên 47 nền kinh tế × 108 cặp quốc gia × năm × 14 mốc khảo sát} --- gấp
2,5 lần phạm vi của Đỗ \& Phan (2026 --- VEFR). Đây là phạm vi rộng nhất
từng có cho nghiên cứu I→P trong văn liệu IB hiện hành. Năm 2025 có
\textbf{14 đợt khảo sát với 16.979 doanh nghiệp} (gồm cả Kiribati 2025).

\textbf{Ẩn dụ Fiji vs Singapore}: Trường hợp Fiji 2025 với tỷ lệ website
74,8\% \textbf{vượt Singapore 66,1\%} (xem Bảng 4.10 file 15) là minh
họa rõ nhất cho ranh giới giữa "lý thuyết Uppsala vận hành ở điểm tối
ưu" và "lý thuyết Uppsala bị biến dạng thành đường cong đảo ngược ở điều
kiện biên". Đối với Fiji và 6 SIDS Pacific khác, \textbf{số hóa không
phải công cụ tối ưu hóa chuỗi cung ứng mà là phương tiện sinh tồn ---
website không phải chiến lược marketing mà là cánh cửa duy nhất kết nối
với thế giới}. Ẩn dụ này dẫn vào Mục 2.6 (3 điều kiện cấu trúc của
U-curve) và Mục 5.7 (SIDS Pacific case file 16) --- nơi điều kiện biên
của lý thuyết Lu \& Beamish (2004) được thiết lập rõ ràng.

\hypertarget{12-mux1ee5c-tiuxeau-chuyuxean-ux111ux1ec1}{%
\subsubsection{1.2 Mục tiêu chuyên
đề}\label{12-mux1ee5c-tiuxeau-chuyuxean-ux111ux1ec1}}

\textbf{Mục tiêu chung}: Hệ thống hóa và phân tích mô tả thực trạng hiệu
quả hoạt động kinh doanh của các doanh nghiệp ở \textbf{47 nền kinh tế
châu Á và Thái Bình Dương} (41 nền kinh tế châu Á + 7 SIDS Thái Bình
Dương boundary extension; pool 101.185 doanh nghiệp) trong giai đoạn
\textbf{2009--2025}.

\textbf{Mục tiêu cụ thể}:

\begin{enumerate}
\def\labelenumi{\arabic{enumi}.}
\tightlist
\item
  Hệ thống hóa khái niệm và đa chiều đo lường hiệu quả hoạt động kinh
  doanh trong văn liệu hiện hành;
\item
  Đề xuất khung phân loại 6 phân nhóm con ICRV cho 47 nền kinh tế (5
  phân nhóm châu Á thuần + 1 phân nhóm boundary SIDS Pacific);
\item
  Phân tích thực trạng đa chiều hiệu quả doanh nghiệp ở \textbf{41 nước
  châu Á (scope chính)} + so sánh với 7 SIDS Pacific (boundary case);
\item
  Nhận diện đặc điểm và yếu tố giải thích sơ bộ cho dị biệt hiệu quả
  giữa các phân nhóm thể chế và giữa các nhóm doanh nghiệp;
\item
  Đề xuất khoảng trống thực tiễn để định hướng nghiên cứu tiếp theo.
\end{enumerate}

\hypertarget{13-phux1ea1m-vi-vuxe0-ux111ux1ed1i-tux1b0ux1ee3ng-nghiuxean-cux1ee9u}{%
\subsubsection{1.3 Phạm vi và đối tượng nghiên
cứu}\label{13-phux1ea1m-vi-vuxe0-ux111ux1ed1i-tux1b0ux1ee3ng-nghiuxean-cux1ee9u}}

\textbf{Đối tượng nghiên cứu}: doanh nghiệp đang hoạt động ở các nền
kinh tế trong pool đáp ứng tiêu chí mẫu của WBES (chủ yếu là doanh
nghiệp khu vực ngoài quốc doanh, có ít nhất 5 lao động, thuộc các ngành
phi nông nghiệp).

\textbf{Phân loại quy mô doanh nghiệp WBES chính thức}: WBES sử dụng 4
nhóm quy mô tiêu chuẩn theo số lao động toàn thời gian + tạm thời:
\textbf{Nhỏ (5-19 lao động)}, \textbf{Vừa (20-99)}, \textbf{Lớn (100+)},
\textbf{Rất lớn (top 1\% nền kinh tế)}. Các kết quả mean/median trong
CĐ1 sử dụng phân tầng này khi cần. \textbf{Bốn ngành loại trừ khỏi vũ
trụ khảo sát WBES}: tiện ích công cộng, dịch vụ chính phủ, chăm sóc sức
khỏe và dịch vụ tài chính --- vì các ngành này có cấu trúc thể chế đặc
thù không phù hợp với khung phân tích doanh nghiệp tư nhân chính thức
(cited World Bank, n.d. --- Tài liệu tóm tắt chỉ số Enterprise Surveys;
file 04 v2.5).

\textbf{Phạm vi không gian}:

\textbf{A. Scope chính --- 41 nước châu Á thuần} (theo tiêu đề luận án
"Châu Á" trong Quyết định 4769/QĐ-ĐHCT 15/10/2024):

\begin{itemize}
\tightlist
\item
  \textbf{(i) Advanced đổi mới sáng tạo dẫn dắt (innovation-driven, 5
  nước)}: Singapore, Hong Kong SAR, Hàn Quốc, Đài Loan, Israel.
\item
  \textbf{(ii) Advanced tài nguyên dẫn dắt (resource-driven, 5 nước)}:
  Saudi Arabia, Qatar, Kuwait, Bahrain, Brunei.
\item
  \textbf{(iii) Upper-middle (6 nước)}: Trung Quốc, Malaysia, Thái Lan,
  Kazakhstan, Armenia, Georgia.
\item
  \textbf{(iv) Emerging (7 nước)}: Việt Nam, Indonesia, Philippines, Ấn
  Độ, Sri Lanka, Jordan, Mongolia.
\item
  \textbf{(v) Frontier (17 nước)}: Bangladesh, Pakistan, Lào, Campuchia,
  Myanmar, Nepal, Bhutan, Maldives, Uzbekistan, Tajikistan, Kyrgyz
  Republic, Turkmenistan, Afghanistan, Timor-Leste, Iraq, Lebanon,
  Yemen.
\item
  Cyprus được bao gồm trong nhóm Advanced (cụm Mediterranean --- biên
  giới Đông -- Tây Á về địa lý).
\end{itemize}

\begin{quote}
\textbf{Lưu ý phân biệt phạm vi ICRV vs WB geographic scope (v3.14)}:
Khung ICRV trong chuyên đề này bao gồm \textbf{47 nền kinh tế} với 5
phân nhóm châu Á thuần (gồm Kazakhstan/Armenia/Georgia trong
Upper-middle; Kyrgyzstan/Tajikistan/Uzbekistan/Turkmenistan trong
Frontier) vì các nền kinh tế này có dữ liệu WBES đầy đủ, mức độ dị biệt
thể chế có giá trị phân tích và được xếp vào "Emerging/Central Asia"
theo khung phân loại ADB + nhiều nghiên cứu IB. Tài liệu tham chiếu
\texttt{manuscripts/wb\_country\_classification\_2025\_2026.md} sử dụng
phạm vi hẹp hơn: WB geographic regions EAP + SAS + MENAAP (không bao gồm
ECA), phục vụ mục đích tra cứu phân loại thu nhập FY2026 WB Atlas Method
--- hai tài liệu phục vụ mục đích khác nhau. Để tra cứu chi tiết tình
trạng WBES data availability cho từng nền kinh tế châu Á (Included vs.
Excluded trong PhD + lý do), xem
\texttt{manuscripts/wbes\_asia\_country\_scope.md}.
\end{quote}

\textbf{B. Boundary case extension --- 7 SIDS Thái Bình Dương} (theo
khung "Asia and Pacific" của ADB Regional Classification và World Bank
East Asia and Pacific Region):

\begin{itemize}
\tightlist
\item
  \textbf{(vi) SIDS Thái Bình Dương (7 nước, n=1.371, 1,4\% pool)}:
  Fiji, Papua New Guinea, Solomon Islands, Tonga (gồm cả đợt 2024),
  Vanuatu, Samoa, \textbf{Kiribati (mới 2025)}.
\end{itemize}

\textbf{Lý do bao gồm SIDS Pacific như boundary case extension} (theo
nguyên tắc \emph{boundary condition explicit} của Bello \& Kostova,
2012, và \emph{phenomenon-based scope} của Meyer et al., 2017):

(a) \textbf{Khung "Asia and Pacific" thống nhất} của ADB Regional
Classification (ADB, 2024) và World Bank EAP Region được áp dụng xuyên
suốt các báo cáo thực tiễn --- đảm bảo so sánh nhất quán với văn liệu IB
hiện hành;

(b) \textbf{WBES coverage} cho SIDS Pacific là phần dữ liệu Asia mà
không có nguồn alternative chuẩn hóa khác --- bảo đảm tính khả dụng dữ
liệu so sánh xuyên quốc gia;

(c) \textbf{Giá trị khoa học (scientific value)}: SIDS Pacific là
\textbf{trường hợp biên cực trị (extreme boundary case)} giúp minh họa
và kiểm chứng patterns châu Á qua tương phản --- đặc biệt H6 (forced
internationalization penalty) trong Chuyên đề 2. Kiribati 2025 với FSTS
1,03\%, FDI 0,7\%, website 18,7\% là extreme nhất pool (xem Mục 5.7);

(d) \textbf{Lý thuyết phenomenon-based} (Meyer et al., 2017): hiện tượng
"isolated SIDS" (Kiribati extreme) làm sáng tỏ điều kiện biên của lý
thuyết Uppsala (Johanson \& Vahlne, 1977) và liability of foreignness
(Zaheer, 1995) ở scale cực nhỏ.

\textbf{Quan hệ SIDS Pacific với scope chính}: SIDS Pacific (1,4\% pool)
\textbf{không phải scope chính} của tiêu đề luận án; được đưa vào với 3
vai trò minh họa:

\begin{itemize}
\tightlist
\item
  (i) \textbf{Trường hợp biên cực trị (extreme boundary case)} cho ICRV
  regime classification (Mục 3.2);
\item
  (ii) \textbf{Đối chứng (counterfactual)} cho pattern châu Á --- đặc
  biệt Mục 5.7 SIDS Thái Bình Dương và Mục 7.3.1 (3) digital theatre;
\item
  (iii) \textbf{Bằng chứng cho H6} (forced internationalization penalty)
  trong Chuyên đề 2.
\end{itemize}

\textbf{Chi tiết phân tích firm-level cho SIDS Pacific} (gồm Kiribati
extreme; Pacific MIRAB model, Bertram, 2006; vulnerability, Briguglio,
1995) được tách thành \textbf{bản thảo P8 riêng} (Đỗ \& Phan, đang chuẩn
bị, target \emph{World Development}) để bảo đảm \textbf{scope coherence}
cho luận án thuần "Châu Á" theo tiêu đề chính thức.

\textbf{Phạm vi thời gian}: \textbf{2009--2025} (16 năm × 14 mốc khảo
sát: 2009--2016, 2018--2025). Bao trùm ba thế hệ khung dữ liệu WBES ---
PICS3 (2009--2012, n=14.335), Standardized (2013--2017, n=25.046),
Standardized 2018+ với BREADY/BEE/EAP Core (2018--2025, n=61.804).

\textbf{Phạm vi nội dung}: hiệu quả hoạt động kinh doanh đa chiều ---
năng suất lao động, lợi nhuận, tăng trưởng, đổi mới sáng tạo, cấu trúc
doanh nghiệp.

\hypertarget{14-phux1b0ux1a1ng-phuxe1p-tiux1ebfp-cux1eadn}{%
\subsubsection{1.4 Phương pháp tiếp
cận}\label{14-phux1b0ux1a1ng-phuxe1p-tiux1ebfp-cux1eadn}}

Chuyên đề tiếp cận theo lối \textbf{mô tả -- chẩn đoán}
(descriptive-diagnostic), không thiết lập mô hình hồi quy đa biến để
kiểm định nhân quả. Phương pháp này cung cấp bức tranh tổng thể đáng tin
cậy làm nền tảng thực tiễn cho Chuyên đề 2 (mô hình lý thuyết và thực
nghiêm) và luận án (kiểm định thực nghiệm).

Cụ thể, chuyên đề kết hợp ba kỹ thuật: (a) thống kê mô tả đa biến (trung
vị, trung bình, độ phân tán, kurtosis) theo phân nhóm thể chế và theo
ngành; (b) trực quan hóa dữ liệu (heatmap, kernel density, spider chart,
scatter plot --- 11 hình); (c) phân tích hai biến (bivariate). Khung
PESTEL bổ trợ ở cấp độ giải thích bối cảnh.

\textbf{Áp dụng nguyên tắc transparent data section của Aguinis et al.
(2019 --- \emph{Journal of Management})}: chuyên đề trình bày rõ (i)
population definition (47 nền kinh tế Asia + Pacific với phân tách 41 vs
7); (ii) sampling procedure (WBES official sampling); (iii) sample size
justification (pool 101.185); (iv) response rate (per WBES Country
Profile each); (v) missing data handling (FSTS = d3b + d3c, winsorize
1/99); (vi) operationalization variables; (vii) reliability/validity
(TCI = R\&D + ISO multi-component aspirational, DAI = website
single-component v3.0, multi-component v3.2); (viii) estimation method
(descriptive only in CĐ1; OLS + IV in CĐ2); (ix) standard errors (n/a
CĐ1); (x) robustness (cross-validate with VEFR baseline 17 nước); (xi)
open science (replication package \texttt{wbes/} + \texttt{papers/} +
\texttt{figures/}); (xii) data sharing (WBES public + processed pool
documented in Phụ lục E).

\textbf{Khung lý giải lý thuyết bổ sung --- Kafouros et al. (2023 ---
\emph{Global Strategy Journal})}: Phương pháp descriptive được dẫn dắt
bởi ba cơ chế (mechanisms) qua đó chất lượng thể chế tác động đến hiệu
quả doanh nghiệp ở các nền kinh tế mới nổi: (a) \textbf{chi phí và mức
độ dễ dàng nhận diện và thiết lập đối tác} (partnership identification)
--- giải thích vì sao DAI và mức độ tham gia website thay đổi mạnh xuyên
các phân nhóm thể chế; (b) \textbf{hiệu lực của khung pháp lý và sở hữu
trí tuệ} (legal framework effectiveness) --- giải thích sự phân tán
R\&D, ISO và TCI; (c) \textbf{trao đổi giữa các doanh nghiệp} (interfirm
market exchange) --- giải thích phân tán năng suất nội bộ. Khung 3 cơ
chế này là cơ sở để CĐ2 thiết lập các mô hình điều tiết (moderation) cho
giả thuyết H7 mới --- tương tác giữa chất lượng thể chế và động học
ngành (industry × institutional dynamism interaction).

\textbf{Khung 3 trụ cột bổ sung, IFC Private Sector Development
Blueprint}: Báo cáo IFC (n.d., PSD Blueprint) phân tích WBES theo ba trụ
cột bổ trợ cho Aguinis 12-point và Kafouros 3 mechanisms: (i)
\textbf{nền tảng thể chế} (institutional foundations --- qui định thuế,
quyền sở hữu tài sản, giải quyết tranh chấp, \textbf{6 essential public
transactions}: cấp phép hoạt động, cấp phép xây dựng, kết nối điện, kết
nối nước, cấp phép khảo sát đất, hoàn thuế VAT); (ii) \textbf{nguồn lực
vận hành} (operational resources, lao động, tài chính, năng lượng,
internet); (iii) \textbf{động lực thị trường} (market dynamics, cạnh
tranh, FDI, xuất nhập khẩu). Đáng chú ý: doanh nghiệp xuất khẩu mất
\textasciitilde1 tuần để hoàn tất thủ tục (xuất) / \textasciitilde2 tuần
(nhập); \textasciitilde40\% doanh nghiệp coi thủ tục thương mại là rào
cản lớn nhất (theo IFC analysis). Khung 3 trụ cột này là tham chiếu cho
Mục 6 yếu tố giải thích trong CĐ1 v3.x --- bổ sung lăng kính phân tích
doanh nghiệp WBES tách 3 chiều cấu trúc.

\textbf{Chi tiết phương pháp luận WBES Global}: Áp dụng phương pháp lấy
mẫu Global từ năm 2006 với ba nguyên tắc kỹ thuật: (a) \textbf{Trọng số
lấy mẫu} (sampling weights): mọi giá trị trung bình và trung vị trong
CĐ1 đều được tính có trọng số để phản ánh đúng tổng thể doanh nghiệp
thực tế; (b) \textbf{Xử lý ngoại lai} (outliers): các phản hồi cực đoan
được xác định bằng phương pháp chuyển đổi logarit và loại bỏ nếu nằm
ngoài 3 độ lệch chuẩn so với trung bình, đồng thời winsorize log năng
suất ở mức 1/99 phân vị trong cụm quốc gia × năm; (c) \textbf{Ngưỡng
quan sát}: tối thiểu 5 quan sát cho trung bình/trung vị và 30 quan sát
cho tốc độ tăng trưởng; nếu không đủ dẫn đến ghi "n.a." (không áp dụng).
Các nguyên tắc này bảo đảm tính so sánh xuyên quốc gia và xuyên thời
gian của 47 nền kinh tế trong pool 101.185 doanh nghiệp (cited World
Bank, 2025, Hệ thống phân loại quốc gia FY26 + n.d., Tài liệu tóm tắt
chỉ số Enterprise Surveys; file 04 v2.5).

\hypertarget{15-ux111uxf3ng-guxf3p-dux1ef1-kiux1ebfn-cux1ee7a-chuyuxean-ux111ux1ec1}{%
\subsubsection{1.5 Đóng góp dự kiến của chuyên
đề}\label{15-ux111uxf3ng-guxf3p-dux1ef1-kiux1ebfn-cux1ee7a-chuyuxean-ux111ux1ec1}}

\textbf{Bốn đóng góp scope chính cho 41 nước châu Á}:

(1) \textbf{Pattern điển hình (stylized facts)} đa chiều cho hiệu quả
doanh nghiệp châu Á 2009--2025 --- phạm vi rộng nhất trong văn liệu IB
hiện hành (Wu et al., 2022 dừng ở 2020; Arte \& Larimo, 2022 dừng ở
2019). Việc bao gồm \textbf{dữ liệu năm 2025} cho 13 nước châu Á (Ấn Độ,
Saudi Arabia, Thái Lan, Sri Lanka, Mongolia, Qatar, Afghanistan,
Maldives, Brunei, Kuwait, Nepal, Solomon Islands*, Fiji*) đặt chuyên đề
ở vị trí cập nhật nhất.

(2) \textbf{Phân nhóm con Advanced} (innovation-driven so với
resource-driven) --- phát hiện lý thuyết MỚI chưa có trong văn liệu hiện
hành (kể cả Đỗ \& Phan, 2026 --- VEFR vốn chỉ tập trung emerging Asia).
Tỷ số phân tán Singapore (1,03) so với Vùng Vịnh (0,49) ≈ \textbf{2,1
lần} (xem Mục 5.2 + Hình 4.1).

(3) \textbf{Nhảy vọt số (digital leapfrog)} ở Frontier và Emerging châu
Á 2018--2025, mở rộng pattern hiệu ứng lá chắn số (digital shield
effect) từ 17 nước (Đỗ \& Phan, 2026, VEFR) lên 47 nước (xem Mục 4.4 và
Mục 4.6). Phân tách hai tầng (Tier) đo lường --- \textbf{Tier-1 Digital
Presence} (website \textbf{nhị phân (binary)} có/không, đo trong Spec 1
2009--2025, đã \textbf{bão hòa (saturation)} ở Advanced từ trước 2018)
vs \textbf{Tier-2 Digital Transformation} (DAI đa thành phần ---
multi-component, 5 thành phần: website + email + thanh toán di động
(smartphone POS) + thương mại điện tử (e-commerce) + điện toán đám mây
(cloud) --- đo trong Spec 2 2018--2025) --- minh bạch hóa rằng hệ số âm
DAI ở Advanced trong Spec 1 (-0,129) là \textbf{biểu hiện giả tạo do bão
hòa Tier-1 (artifact của Tier-1 saturation)}, không phải "số hóa làm
giảm năng suất ở Singapore/Hàn Quốc". Kiểm định loại trừ ICT (ICT
exclusion test) ở Mục 4.4.5 file 15 xác nhận đây là biểu hiện giả tạo:
khi loại bỏ doanh nghiệp ICT khỏi Advanced phân nhóm đổi mới sáng tạo
dẫn dắt, hệ số âm biến mất dẫn đến biến website không còn \textbf{phân
biệt (discriminate)} năng suất giữa các doanh nghiệp ở Advanced vì hầu
hết đã có website từ lâu.

(4) \textbf{Bằng chứng cho khung lý thuyết phi tuyến + điều tiết đa
tầng} --- bốn lần đảo dấu xuyên phân nhóm con thể chế (xem Mục 6 và Hình
6.1) là cơ sở thực tiễn cho \textbf{6 giả thuyết H1-H6} của Chuyên đề 2.
Bốn lần đảo dấu này được lý giải sâu hơn qua hai hướng lý thuyết bổ
sung. \textbf{Thứ nhất}, \textbf{Xu (2024, \emph{Global Strategy
Journal})} phân tách "exposure de jure" (luật trên giấy) vs "enforcement
de facto" (thực thi thực tế), gợi ý rằng các điểm chuyển dấu
cross-regime không thể chỉ giải thích bằng điểm số WGI hay khoảng cách
thể chế tổng hợp; cần phân biệt formal rules và implementation capacity,
đặc biệt cho Việt Nam, Trung Quốc, Indonesia, Mongolia có gap de
jure--de facto lớn (xem Mục 6 và Mục 7.3.3). \textbf{Thứ hai},
\textbf{Kafouros et al. (2023, \emph{Global Strategy Journal})} đề xuất
chất lượng thể chế tương tác hai chiều với động học ngành, khuếch đại
tích cực ở ngành công nghệ-động (technological dynamism) và suy yếu ở
ngành thị trường-động (market dynamism). Hai hướng này mở đường cho
\textbf{giả thuyết H7 mới của CĐ2} --- \emph{industry × institutional
dynamism interaction effect} --- và cho biểu đồ Bảng 4.8.1 dynamism
profile theo Kafouros (xem Mục 4.8 file 15).

\textbf{Một đóng góp boundary case mở rộng (SIDS Pacific)}:

(5) \textbf{Trường hợp biên (boundary case) chi phí buộc phải quốc tế
hóa (forced internationalization penalty)} ở \textbf{7 SIDS Thái Bình
Dương} (Briguglio, 1995; Bertram, 2006) --- bằng chứng firm-level đầu
tiên trong văn liệu IB cho hiện tượng này (xem Mục 5.7). Kiribati 2025
là extreme nhất với FSTS 1,03\%, FDI 0,7\%, website 18,7\%. \textbf{Đóng
góp này vượt scope chính của tiêu đề luận án nhưng làm sáng tỏ điều kiện
biên của lý thuyết Uppsala} (Johanson \& Vahlne, 1977) --- minh họa qua
tương phản. Chi tiết phân tích cho SIDS Pacific tách thành bản thảo P8
riêng (Đỗ \& Phan, đang chuẩn bị, target \emph{World Development}).

\hypertarget{16-kux1ebft-cux1ea5u-chuyuxean-ux111ux1ec1}{%
\subsubsection{1.6 Kết cấu chuyên
đề}\label{16-kux1ebft-cux1ea5u-chuyuxean-ux111ux1ec1}}

Chuyên đề gồm bảy chương: Ch.1 Giới thiệu; Ch.2 Cơ sở lý luận; Ch.3
Khung phân loại 6 phân nhóm con ICRV; Ch.4 Thực trạng từ WBES (10 mục
Mục 4.1-Mục 4.10); Ch.5 Bảy tiểu cảnh điển hình (6 châu Á + 1 SIDS
boundary case); Ch.6 Yếu tố giải thích sơ bộ; Ch.7 Khoảng trống thực
tiễn và kết luận.

\begin{center}\rule{0.5\linewidth}{0.5pt}\end{center}

\hypertarget{chux1b0ux1a1ng-2--cux1a1-sux1edf-luxfd-luux1eadn-vux1ec1-hiux1ec7u-quux1ea3-houx1ea1t-ux111ux1ed9ng-kinh-doanh}{%
\subsection{CHƯƠNG 2 --- CƠ SỞ LÝ LUẬN VỀ HIỆU QUẢ HOẠT ĐỘNG KINH
DOANH}\label{chux1b0ux1a1ng-2--cux1a1-sux1edf-luxfd-luux1eadn-vux1ec1-hiux1ec7u-quux1ea3-houx1ea1t-ux111ux1ed9ng-kinh-doanh}}

\hypertarget{21-khuxe1i-niux1ec7m-hiux1ec7u-quux1ea3-houx1ea1t-ux111ux1ed9ng-kinh-doanh}{%
\subsubsection{2.1 Khái niệm hiệu quả hoạt động kinh
doanh}\label{21-khuxe1i-niux1ec7m-hiux1ec7u-quux1ea3-houx1ea1t-ux111ux1ed9ng-kinh-doanh}}

Hiệu quả hoạt động kinh doanh (firm performance) là khái niệm \textbf{đa
chiều} (multidimensional) phản ánh mức độ đạt được các mục tiêu chiến
lược của doanh nghiệp theo thời gian và trong bối cảnh thể chế cụ thể.
Venkatraman và Ramanujam (1986) phân biệt hai tầng: (a) \textbf{hiệu quả
tài chính} (financial performance) --- doanh thu, lợi nhuận, tỷ suất
sinh lợi; và (b) \textbf{hiệu quả hoạt động rộng hơn} (broader
organizational effectiveness) --- thị phần, đổi mới, năng suất, tăng
trưởng việc làm. Lý thuyết dựa trên nguồn lực (Barney, 1991; Wernerfelt,
1984) xem hiệu quả doanh nghiệp là kết quả của việc tích lũy và khai
thác \textbf{nguồn lực chiến lược} (valuable, rare, inimitable,
non-substitutable --- VRIN), trong khi quan điểm thể chế (North, 1990;
Scott, 1995; Peng, 2001) nhấn mạnh rằng hiệu quả này bị định hình mạnh
bởi \textbf{môi trường thể chế} nơi doanh nghiệp hoạt động.

Trong bối cảnh nghiên cứu xuyên quốc gia, khái niệm hiệu quả cần được
\textbf{tiêu chuẩn hóa} để so sánh được. Theo Cusolito và Maloney
(2018), \textbf{năng suất lao động} (labor productivity, LP), đo bằng
giá trị gia tăng trên lao động hoặc doanh thu trên lao động --- là thước
đo phổ biến nhất trong các nghiên cứu WBES vì: (a) tính khả dụng cao
xuyên 3 thế hệ schema; (b) ít bị biến dạng bởi cơ cấu vốn hay cơ chế
thuế quốc gia; (c) phản ánh trực tiếp năng lực chuyển đổi đầu vào thành
đầu ra. Trong chuyên đề này, LP được tính bằng:

\textbf{LP = ln(Doanh thu bán hàng hàng năm PPP / Số lao động toàn thời
gian)}

Giá trị LP được winsorize ở mức 1/99 phân vị trong cụm quốc gia × năm để
kiểm soát ngoại lai và được điều chỉnh theo sức mua tương đương (PPP)
nhằm đảm bảo tính so sánh xuyên quốc gia (World Bank Enterprise Surveys,
2025).

Bên cạnh LP, chuyên đề tiếp cận hiệu quả theo \textbf{bốn chiều} (Mục
2.2) để cung cấp bức tranh toàn diện hơn. Cách tiếp cận này kế thừa
khung đo lường đa chiều của Murphy et al. (1996) và được điều chỉnh để
phù hợp với độ bao phủ biến của WBES.

\hypertarget{22-bux1ed1n-chiux1ec1u-ux111o-lux1b0ux1eddng-hiux1ec7u-quux1ea3}{%
\subsubsection{2.2 Bốn chiều đo lường hiệu
quả}\label{22-bux1ed1n-chiux1ec1u-ux111o-lux1b0ux1eddng-hiux1ec7u-quux1ea3}}

Hiệu quả doanh nghiệp trong chuyên đề này được đo lường theo \textbf{bốn
chiều} tương ứng với bốn nhóm biến WBES:

\textbf{Chiều 1 --- Năng suất lao động (Labor Productivity)}. Biến đo
lường chính: LP = ln(doanh thu hàng năm PPP / số lao động toàn thời
gian). Đây là \textbf{thước đo chính} (primary metric) xuyên suốt phân
tích, được lựa chọn vì: (a) có thể tính cho hầu hết các nền kinh tế
trong pool; (b) không bị nhiễu bởi cơ cấu thuế và chính sách chuyển giá
quốc gia; (c) có giá trị tham chiếu rõ ràng trong văn liệu quốc tế
(Hsieh \& Klenow, 2009, 2014).

\textbf{Chiều 2 --- Lợi nhuận và biên lợi nhuận (Profitability)}. Biến
đo lường: (a) tỷ suất lợi nhuận hoạt động từ câu hỏi WBES về doanh thu
sau chi phí lao động; (b) chỉ số gián tiếp qua khả năng tái đầu tư
(self-financing capacity). Hạn chế: nhiều doanh nghiệp trong WBES không
tiết lộ lợi nhuận cụ thể --- do đó chiều này mang tính hỗ trợ thay vì
chủ đạo.

\textbf{Chiều 3 --- Tăng trưởng (Growth)}. Biến đo lường: (a) CAGR tăng
trưởng việc làm (employment growth) giữa 3 năm trước và thời điểm khảo
sát; (b) tỷ lệ doanh nghiệp xuất khẩu (export participation) như proxy
cho mở rộng thị trường. Tăng trưởng việc làm được ưu tiên vì: (a) ít
nhạy cảm với thay đổi giá cả so với doanh thu; (b) phản ánh khả năng tạo
việc làm bền vững --- quan trọng cho mục tiêu chính sách.

\textbf{Chiều 4 --- Đổi mới sáng tạo và cấu trúc doanh nghiệp
(Innovation and Structure)}. Biến đo lường: (a) R\&D (tỷ lệ doanh nghiệp
có chi tiêu R\&D dương); (b) đổi mới sản phẩm mới; (c) đổi mới quy trình
mới; (d) chứng nhận ISO; (e) website (Tier-1 Digital Presence); (f) FDI
≥10\% (cơ cấu sở hữu). Chiều thứ tư này phản ánh \textbf{năng lực tương
lai} (forward-looking capacity) --- khả năng duy trì hiệu quả trong dài
hạn qua đổi mới và nâng cấp công nghệ.

\hypertarget{23-nux103m-meta-analysis-lux1edbn-vux1ec1-quan-hux1ec7-quux1ed1c-tux1ebf-huxf3a--hiux1ec7u-quux1ea3-19802024}{%
\subsubsection{2.3 Năm meta-analysis lớn về quan hệ quốc tế hóa -- hiệu
quả
(1980--2024)}\label{23-nux103m-meta-analysis-lux1edbn-vux1ec1-quan-hux1ec7-quux1ed1c-tux1ebf-huxf3a--hiux1ec7u-quux1ea3-19802024}}

Một meta-analysis 2024 bổ sung của Aydemir và cộng sự (2024) trên k=170
mẫu (n=334,855 MNCs từ SSCI và Scopus) củng cố quy mô hiệu ứng dương nhỏ
giữa quốc tế hoá và hiệu quả doanh nghiệp, với biến điều tiết quan trọng
gồm thế hệ phát triển, đòn bẩy tài chính, và R\&D intensity. Wagner
(2025) tổng hợp bằng chứng cấp doanh nghiệp Nordic, nhấn mạnh rằng các
pattern xuyên quốc gia khó tổng quát hoá thành stylised facts cho chính
sách. Hai cập nhật này hỗ trợ luận điểm về tính không đồng nhất xuyên
bối cảnh trong phần §3 ICRV.

Văn liệu về quan hệ quốc tế hóa dẫn đến hiệu quả (I→P) là một trong
những dòng nghiên cứu tích lũy nhiều nhất trong lý thuyết kinh doanh
quốc tế (International Business --- IB). Năm meta-analysis lớn cung cấp
nền tảng thực nghiệm:

\textbf{(1) Bausch \& Krist (2007)}: Tổng hợp \textbf{68 nghiên cứu
(k=68)} giai đoạn 1980--2005. Hệ số tương quan trung bình r=0,045 ---
tích cực nhưng nhỏ. Phát hiện chính: đo lường I (internationalization)
và đo lường P (performance) rất không đồng nhất giữa các nghiên cứu dẫn
đến làm hạn chế khả năng so sánh. Khoảng trống: ít nghiên cứu châu Á mới
nổi.

\textbf{(2) Kirca et al. (2012)}: Tổng hợp \textbf{180 nghiên cứu} với
phân tích điều tiết meta-analytic (MARA). Phát hiện: quan hệ I→P dương
và có ý nghĩa thống kê; nhưng bị điều tiết bởi: loại doanh nghiệp,
phương thức đo lường, bối cảnh quốc gia. Quan trọng: các doanh nghiệp
châu Á và các nền kinh tế mới nổi biểu hiện pattern khác biệt so với
doanh nghiệp OECD.

\textbf{(3) Marano et al. (2016)}: Tập trung vào \textbf{thể chế hóa}
(institutionalization) nghiên cứu I→P. Phát hiện: chất lượng thể chế
nước chủ nhà là biến điều tiết quan trọng nhất --- giải thích phần lớn
phương sai không giải thích được trong các meta-analysis trước. Đặt nền
tảng cho khung ICRV trong chuyên đề này.

\textbf{(4) Schwens et al. (2018)}: Tổng hợp \textbf{các yếu tố điều
tiết} mức độ tác động của internationalization. Phát hiện: năng lực hấp
thụ (absorptive capacity --- Cohen \& Levinthal, 1990), chất lượng CEO,
và phát triển thể chế là các moderators quan trọng nhất. Gợi ý: cần đo
lường TCI và TMT characteristics --- trực tiếp liên quan đến Chuyên đề
2.

\textbf{(5) Wu, Fan \& Chen (2022)}: Cập nhật meta-analysis đến 2020,
tập trung vào \textbf{doanh nghiệp mới nổi châu Á (EMNEs)}. Phát hiện:
quan hệ I→P ở EMNEs có hình dạng U ngược (inverted-U) rõ nét hơn so với
doanh nghiệp OECD; điểm uốn (turning point) thường nằm trong phạm vi
30-50\% FSTS. Phát hiện này là nền tảng cho giả thuyết H1 của Chuyên đề
2.

\textbf{Khoảng trống từ 5 meta-analyses}: Không một meta-analysis nào
trong số trên có dữ liệu WBES xuyên 47 nền kinh tế với phân loại ICRV 6
nhóm. Không có nghiên cứu nào phân biệt tường minh Tier-1 vs Tier-2
digital adoption như biến điều tiết riêng biệt. Khoảng trống này là nền
tảng cho Mục 7.1 (5 khoảng trống nghiên cứu).

\hypertarget{24-huxe0i-huxf2a-huxf3a-ux111o-lux1b0ux1eddng-xuyuxean-quux1ed1c-gia-trong-wbes}{%
\subsubsection{2.4 Hài hòa hóa đo lường xuyên quốc gia trong
WBES}\label{24-huxe0i-huxf2a-huxf3a-ux111o-lux1b0ux1eddng-xuyuxean-quux1ed1c-gia-trong-wbes}}

World Bank Enterprise Surveys trải qua \textbf{ba thế hệ schema} trong
giai đoạn nghiên cứu --- tạo ra thách thức hài hòa hóa (harmonization)
quan trọng:

\textbf{Thế hệ 1 --- PICS3/MENA-WBES (2009--2012, n=14.171)}: Bảng hỏi
PICS3 dùng cho khu vực Đông Nam Á và Trung Á; MENA-WBES cho Trung Đông
và Bắc Phi. Câu hỏi xuất khẩu hỏi tổng quát hơn (d3a = "có xuất khẩu
không?"; d3b+d3c = "\% doanh thu từ xuất khẩu trực tiếp + gián tiếp").
Website được đo qua câu e8.

\textbf{Thế hệ 2 --- Standardized (2013--2017, n=24.564)}: WBES
Standardized Question Core được áp dụng đồng nhất xuyên toàn cầu từ
2013. Cấu trúc câu hỏi xuất khẩu giữ nguyên (d3b+d3c). Website = c22b.
Thanh toán điện tử lần đầu xuất hiện (k33/k38 ở một số quốc gia).

\textbf{Thế hệ 3 --- BREADY/BEE/EAP Core (2018--2025, n=62.450)}: Schema
BREADY (Business Readiness and Enterprise Data) tích hợp e-payment (k33,
k38), cloud computing, mobile banking vào các modules bổ sung. Câu hỏi
FSTS được viết lại rõ ràng hơn dẫn đến gây \textbf{đứt gãy schema}
(schema discontinuity) quan sát được: IND FSTS sụt 7,7\%→2,7\% trong đợt
2025 (Bảng 4.10).

\textbf{Giao thức hài hòa hóa} áp dụng trong chuyên đề này:

\begin{itemize}
\tightlist
\item
  \textbf{Biến FSTS}: \texttt{FSTS\ =\ (d3b\ +\ d3c)\ /\ 100}. Với
  PICS3: d3b = \% xuất khẩu trực tiếp; d3c = \% xuất khẩu gián tiếp.
  Tổng FSTS ≤ 100\%.
\item
  \textbf{Biến LP}: \texttt{ln(annual\_sales\ /\ l1\_employees)}. Doanh
  thu winsorize ở 1/99 phân vị trong cụm quốc gia × năm.
\item
  \textbf{Biến TCI (Technological Capability Index)}: z-standardized
  mean của ≥3 trong 4 thành phần: ISO (b8), R\&D (h8),
  innovation\_product (h1), foreign\_tech (e6). Yêu cầu ≥3 non-missing
  để tính tổng hợp.
\item
  \textbf{Biến DAI (Digital Adoption Index)}: Spec 1 (full coverage
  2009--2025): website binary (c22b/e8) --- \textbf{Tier-1 Digital
  Presence}. Spec 2 (BREADY 2018+ only): z-standardized composite
  website + e-payment intensity (k33/k38) --- \textbf{Tier-2 Digital
  Transformation}.
\item
  \textbf{Pipeline Python}: \texttt{wbes/02\_harmonize.py} tái lập được
  --- 5 bước xử lý từ raw WBES đến clean pool 101.185 doanh nghiệp.
\end{itemize}

\textbf{Xử lý missing values}: Quan sát thiếu FSTS = 0 nếu doanh nghiệp
không xuất khẩu (d3a = "No"); gán missing nếu không trả lời. TCI yêu cầu
≥3/4 components. LP gán missing nếu doanh thu hoặc số lao động = 0 hoặc
không có. Ngưỡng quan sát tối thiểu: 5 cho mean/median; 30 cho CAGR.

\hypertarget{25-ba-khung-luxfd-thuyux1ebft-nux1ec1n}{%
\subsubsection{2.5 Ba khung lý thuyết
nền}\label{25-ba-khung-luxfd-thuyux1ebft-nux1ec1n}}

Ba khung lý thuyết cung cấp nền tảng giải thích cho các pattern thực
trạng được quan sát trong Chương 4--6:

\textbf{Khung 1 --- Lý thuyết Uppsala (Uppsala Model)} (Johanson \&
Vahlne, 1977, 2009): Doanh nghiệp quốc tế hóa theo \textbf{lộ trình tiệm
tiến (incremental path)}, bắt đầu từ các thị trường gần về tâm lý
(psychic distance nhỏ) rồi mở rộng dần. Phiên bản 2009 cập nhật: mạng
lưới quan hệ kinh doanh (business network) thay thế khoảng cách tâm lý
như cơ chế dẫn dắt. Hàm ý cho chuyên đề: (a) pattern FSTS thấp ở SIDS
Pacific (6,3\%) và FDI-linked high FSTS ở Vietnam (23,2\%) phù hợp với
đường cong học hỏi Uppsala; (b) DAI (digital tools) như Banalieva \&
Dhanaraj (2019) lập luận, có thể rút ngắn lộ trình tiệm tiến ở giai đoạn
đầu --- nhưng cần Tier-2 capability để phát huy đầy đủ.

\textbf{Khung 2 --- Quan điểm dựa trên nguồn lực và năng lực động}
(Resource-Based View + Dynamic Capabilities) (Barney, 1991; Teece et
al., 1997; Wernerfelt, 1984): Doanh nghiệp đạt lợi thế cạnh tranh bền
vững thông qua nguồn lực VRIN và khả năng tái cấu hình năng lực theo
biến động môi trường. Hàm ý cho chuyên đề: (a) TCI (năng lực công nghệ)
là proxy cho nguồn lực khó sao chép --- giải thích tại sao Advanced
innovation-driven (ISO 29,9\%, R\&D 16,7\%) có năng suất vượt trội; (b)
khả năng hấp thụ (Cohen \& Levinthal, 1990) giải thích tại sao
FDI-spillover tác động tích cực hơn ở Ấn Độ với R\&D 19,8\% so với Việt
Nam 6,1\%.

\textbf{Khung 3, Lý thuyết thể chế} (Institutional Theory) (North, 1990;
Scott, 1995; Peng, 2001, 2003): Thể chế, luật pháp, chuẩn mực, quy tắc
phi chính thức --- định hình chi phí giao dịch và cơ hội chiến lược cho
doanh nghiệp. Peng (2001, 2003) xây dựng \textbf{khung thể chế dựa trên
giao dịch} cho các nền kinh tế mới nổi --- lý giải tại sao quan hệ I→P
không đồng nhất xuyên chế độ thể chế. Hàm ý: (a) bốn lần đảo dấu của hệ
số FDI/FSTS/TCI/DAI xuyên 5 phân nhóm con ICRV (Bảng 6.1) là bằng chứng
trực tiếp cho institutional moderation; (b) Xu (2024) làm sâu thêm bằng
phân biệt \textbf{de jure exposure} (luật trên giấy) vs \textbf{de facto
enforcement} (thực thi thực tế) --- giải thích tại sao Vietnam và China
có gap lớn giữa hai tầng này; (c) khung ICRV 6 nhóm trong Chương 3 của
chuyên đề là operationalization của Institutional Theory theo hướng này.

\hypertarget{mux1ee5c-26-ux111iux1ec1u-kiux1ec7n-biuxean-cho-u-curve-lu--beamish-2004}{%
\subsubsection{Mục 2.6 ĐIỀU KIỆN BIÊN CHO U-CURVE LU \& BEAMISH
(2004)}\label{mux1ee5c-26-ux111iux1ec1u-kiux1ec7n-biuxean-cho-u-curve-lu--beamish-2004}}

Phân tích nội bộ khuyến nghị: việc nhét khái niệm "chi phí buộc phải
quốc tế hóa SIDS Pacific" vào Mục 4.4 file 15 làm mờ ranh giới giữa hai
loại điều tiết thể chế --- (i) thể chế \textbf{điều tiết cường độ lợi
nhuận} từ quốc tế hóa (U-curve operating); (ii) thể chế \textbf{quyết
định sự tồn tại của quy luật U-curve} (boundary case). Cần thiết lập 3
điều kiện cấu trúc cho U-curve từ sớm trong Chương 2 thì khi Chương 4--5
phân tích sự thất bại, mảnh ghép sẽ hoàn hảo. SIDS Pacific là trường hợp
\textbf{đường cong U ngược (Lu \& Beamish, 2004) bị biến dạng thành
đường cong đảo ngược} --- không phải dị biệt ngẫu nhiên mà là điều kiện
biên (boundary condition --- Bello \& Kostova, 2012) của lý thuyết.

\textbf{Ba điều kiện cấu trúc bắt buộc cho U-curve vận hành}
(theoretical foundation cho Chương 4--5 file 15/16 + CĐ2 H6):

(1) \textbf{Điều kiện thị trường nội địa khả thi} (viable domestic
market): Doanh nghiệp có thể bán hàng nội địa ở quy mô đủ để tích lũy
năng lực kinh doanh trước khi quốc tế hóa. Điều kiện này thỏa ở
Singapore (5,9 triệu dân, GDP per capita \textasciitilde\$80.000), Trung
Quốc (1,4 tỷ dân), Việt Nam (98 triệu dân), Indonesia (273 triệu). KHÔNG
thỏa ở Kiribati (130.000 dân, GDP per capita \textasciitilde\$1.700),
Tuvalu (\textasciitilde12.000 dân) hoặc Tonga (\textasciitilde106.000);

(2) \textbf{Điều kiện chi phí thương mại kiểm soát được} (manageable
trade costs): Logistics, vận tải biển/hàng không, regulations cho xuất
nhập khẩu nằm trong khoảng kinh tế. Điều kiện này thỏa ở Châu Á đại lục
--- KHÔNG thỏa đầy đủ ở SIDS Pacific xa hub thương mại (Fiji-Auckland
2.100 km, Kiribati-Suva 2.700 km, không có direct flight thường xuyên;
Vanuatu-Sydney 2.500 km);

(3) \textbf{Điều kiện thể chế vận hành} (functional institutions): Khung
pháp lý, IPR, contract enforcement, banking infrastructure ở mức tối
thiểu cho hoạt động kinh doanh. Điều kiện này thỏa ở Advanced
(Singapore, Korea, Israel) và Upper-middle (Trung Quốc, Malaysia) ---
KHÔNG thỏa đầy đủ ở SIDS isolated (Kiribati ISO 1,3\%, FDI 0,7\%) hay
frontier conflict zones (Afghanistan, Yemen, Iraq); riêng nhóm Emerging
có gap de jure--de facto lớn theo Xu (2024) gây bẫy chuyển hướng nguồn
lực.

\textbf{Bảng 2.6.1}. \emph{Ma trận đối chiếu lý thuyết 3 trạng thái
U-curve theo 3 điều kiện cấu trúc.}

\begin{longtable}[]{@{}llcccl@{}}
\toprule\noalign{}
Trạng thái & Đại diện & ĐK 1 (thị trường) & ĐK 2 (chi phí TM) & ĐK 3
(thể chế) & Hành vi U-curve \\
\midrule\noalign{}
\endhead
\bottomrule\noalign{}
\endlastfoot
\textbf{A. Điểm tối ưu của U-curve} & Singapore, Hong Kong, Hàn Quốc,
Đài Loan & ✓ & ✓ & ✓ & U-curve vận hành chuẩn --- FSTS có optimal level
cho hiệu quả tối đa (Lu \& Beamish, 2004) \\
\textbf{B. Bẫy chuyển hướng nguồn lực (RDT)} & Việt Nam, Indonesia,
Trung Quốc, Mongolia & ✓ & ✓ & ⚠ \emph{(gap de jure--de facto Xu 2024)}
& U-curve vận hành nhưng có \textbf{bẫy chuyển hướng nguồn lực}
(resource diversion trap) --- chi phí thể chế thực thi \textgreater{}
lợi ích quốc tế hóa ở mức trung; cần ICRV phân biệt formal rules vs
implementation capacity \\
\textbf{C. Sự sụp đổ của đường cong} & SIDS Pacific (Fiji, Kiribati,
Tuvalu, Tonga) & ✗ \emph{(thị trường \textless1 triệu dân)} & ✗
\emph{(cô lập địa lý 2.000+ km tới hub)} & ⚠--✗ \emph{(ISO \textless2\%,
banking limited)} & \textbf{U-curve không vận hành} --- quốc tế hóa là
\textbf{buộc phải} (forced internationalization penalty Briguglio 1995;
Bertram 2006); Fiji 74,8\% website \textgreater{} Singapore 66\% là minh
chứng \\
\end{longtable}

\textbf{Tham chiếu}: Lu \& Beamish (2004); Briguglio (1995); Bertram
(2006); Bello \& Kostova (2012); Cuervo-Cazurra et al. (2018); Banalieva
\& Dhanaraj (2019); Xu (2024) cho de jure--de facto gap ở trạng thái B;

Ma trận 3 trạng thái này thiết lập \textbf{tiền đề lý thuyết bắt buộc}
cho:

\begin{itemize}
\tightlist
\item
  \textbf{Mục 3 ICRV regime classification}: Singapore = Advanced
  innovation (trạng thái A); Việt Nam = Emerging (trạng thái B); SIDS
  Pacific = boundary case extension (trạng thái C) --- tách 3 trạng thái
  thay vì gộp;
\item
  \textbf{Mục 5.2 (Singapore archetype)} + \textbf{Mục 5.7 (Pacific SIDS
  case)} trong file 16;
\item
  \textbf{CĐ2 H6 (forced internationalization penalty)}: dự đoán ở SIDS
  Pacific (trạng thái C), FSTS KHÔNG correlate với năng suất vì U-curve
  không vận hành --- hệ số FSTS ở SIDS sẽ không-significant hoặc đảo dấu
  so với Advanced/Emerging.
\end{itemize}

\hypertarget{mux1ee5c-27-khung-4-tier-chuyux1ec3n-ux111ux1ed5i-sux1ed1--muxf4-huxecnh-nux103ng-lux1ef1c-sux1ed1-phux1ee5-thuux1ed9c-bux1ed1i-cux1ea3nh}{%
\subsubsection{Mục 2.7 KHUNG 4-TIER CHUYỂN ĐỔI SỐ + MÔ HÌNH NĂNG LỰC SỐ
PHỤ THUỘC BỐI
CẢNH}\label{mux1ee5c-27-khung-4-tier-chuyux1ec3n-ux111ux1ed5i-sux1ed1--muxf4-huxecnh-nux103ng-lux1ef1c-sux1ed1-phux1ee5-thuux1ed9c-bux1ed1i-cux1ea3nh}}

Phân tích đối chiếu nội bộ (đối chiếu hai bản thảo của tác giả: P3
Singapore, Mar et al., 2026, \emph{MIR}; P4 Việt Nam, Đỗ \& Phan, 2026
--- \emph{APJM}) chỉ ra một construct validity gap quan trọng: DAI đo
lường trong cả 2 papers chỉ chạm tầng cơ bản (Tier 1 = website cho Việt
Nam; Tier 1+2 = website + e-payment cho Singapore), trong khi lý thuyết
về \emph{super-linear coordination cost} (Brynjolfsson \& McAfee, 2014)
đòi hỏi tầng tích hợp sâu (Tier 3) hoặc năng lực động (Tier 4). Cần
thiết lập explicit \textbf{construct boundary} từ Chương 2, trước khi
báo cáo descriptive results, để phòng tránh "overclaiming" trong khi nộp
Hội đồng CTU + tránh "desk-reject" khi submit Q1/Q2 IB journals.

\textbf{Bốn tầng bậc số hóa theo Verhoef et al. (2021 --- \emph{Journal
of Business Research})}:

\begin{longtable}[]{@{}lllll@{}}
\toprule\noalign{}
Tầng & Tên (Việt) & Tên (Anh) & Đại diện đo lường & Quan sát được trong
WBES? \\
\midrule\noalign{}
\endhead
\bottomrule\noalign{}
\endlastfoot
\textbf{Tier 1} & Sự hiện diện số cơ bản & \emph{Digital Presence} &
website binary, email & ✓ Có (toàn pool 2009-2025) \\
\textbf{Tier 2} & Giao tiếp số + thương mại điện tử cơ bản &
\emph{Digital communication + basic e-commerce} & e-payment intensity
(k33, k38), online ordering & ✓ Có một phần (BREADY 2018+) \\
\textbf{Tier 3} & Tích hợp quy trình số & \emph{Digital process
integration} & ERP, CRM, supply-chain digitization & ✗ KHÔNG quan sát \\
\textbf{Tier 4} & Năng lực số động & \emph{Dynamic digital capability} &
AI integration, platform orchestration, data-driven decision & ✗ KHÔNG
quan sát \\
\end{longtable}

\textbf{Tuyên bố ranh giới đo lường (Construct Boundary Statement)}:
Pool 101.185 doanh nghiệp WBES quan sát Tier 1-2 (website cho toàn giai
đoạn; e-payment cho wave BREADY 2018+). KHÔNG quan sát Tier 3-4. Do đó,
mọi lập luận về "năng lực số" trong CĐ1 + papers IB thuộc luận án phải
hạn chế ở phạm vi \textbf{Foundational Digital Adoption (FDA)}, KHÔNG mở
rộng thành "Digital Capability" hay "Digital Transformation". Đây là
tiêu chuẩn nhất quán xuyên CĐ1 / CĐ2 / luận án 5 chương.

\textbf{Mô hình Năng lực số phụ thuộc bối cảnh (Context-Contingent
Digital Capability Model, CDCM)}, đóng góp lý thuyết MỚI hợp nhất hai
papers:

CDCM hợp nhất phát hiện Singapore P3 (Mar et al., 2026) + Việt Nam P4
(Đỗ \& Phan, 2026 --- \emph{APJM}) thành một khung xuyên quốc gia 3
chiều:

\textbf{(a) Hình dáng đường cong I-P
(Internationalization-Performance)}: Việt Nam dẫn đến Inverted-U rõ nét,
điểm uốn (TP) ≈ 39-46\% FSTS xuyên 3 wave 2009-2023; Lind-Mehlum
p\textless0,013. Singapore dẫn đến tuyến tính dương + bậc hai nhẹ, TP ≈
82\% FSTS (vùng dữ liệu thưa, 3,2\% firms vượt 70\% FSTS); Lind-Mehlum
p=0,303 KHÔNG xác nhận.

\textbf{(b) Cơ chế DAI}: Việt Nam (Tier 1 = website only) dẫn đến
\textbf{stage-contingent}: mạnh 2009 dẫn đến vô tác dụng 2015 dẫn đến
phục hồi 2023 nhưng tương tác DAI×FSTS mang dấu âm (β=-0,912, p=,043) ở
mức xuất khẩu cao dẫn đến Tier 1 trở thành \textbf{điểm nghẽn}.
Singapore (Tier 1+2 = website + e-payment) dẫn đến \textbf{conditional
scaling}: tương tác DAI×FSTS² dương mạnh (β=3,119, p=,005) dẫn đến Tier
1+2 trở thành \textbf{đòn bẩy} ở xuất khẩu cao.

\textbf{(c) Vai trò TCI (Năng lực công nghệ)}: Việt Nam, \emph{scarce
advantage} giúp doanh nghiệp vượt rào cản quốc tế hóa (TCIz IV β=1,639,
p\textless,001 robust 2SLS, Đỗ \& Phan, 2026 --- \emph{APJM}). Singapore
--- \emph{hygiene factor} nâng cao productivity floor nhưng KHÔNG điều
tiết I-P curve.

\textbf{Lý giải xuyên quốc gia --- Institutional Transaction Costs (TCE
Coase-Williamson, mở rộng Banalieva \& Dhanaraj 2019)}: Ở Việt Nam, hạ
tầng thương mại số yếu khiến chi phí điều phối xuyên biên giới tăng vọt
ở mức xuất khẩu trung bình (TP 39-46\%). Ở Singapore, hạ tầng số +
logistics đã hấp thụ phần lớn chi phí này, đẩy ranh giới quá tải
(overload) về vùng FSTS cực cao (\textasciitilde82\%).

\textbf{Hàm ý phương pháp luận cho CĐ2 + luận án}:

\begin{itemize}
\tightlist
\item
  \textbf{Specification 1} (Tier 1-only website, full pool 2009-2025):
  expected hệ số DAI âm hoặc null ở Advanced phân nhóm innovation-driven
  do bão hòa Tier 1 (xem Mục 1.5 contribution \#3 + Mục 4.4.5 ICT
  exclusion test file 15);
\item
  \textbf{Specification 2} (Tier 1+2 multi-component, BREADY 2018+
  wave): expected positive coefficient + interaction với FSTS² dương ở
  Advanced --- replication pattern Singapore P3;
\item
  \textbf{Robustness check \#10} (mới): cross-validation construct
  boundary --- chạy Spec 2 chỉ ở 41 nước có đủ Tier 1+2 data (BREADY+
  wave); confirm rằng pattern Singapore-style emerges khi có composite
  digital adoption metric;
\item
  \textbf{Hạn chế minh bạch (Limitations)}: WBES không quan sát Tier 3
  (ERP/CRM) và Tier 4 (AI). Đề xuất nghiên cứu tương lai sử dụng panel
  data đo trực tiếp các hệ thống ERP/AI để xác định ranh giới giữa "đòn
  bẩy" (lever) và "điểm nghẽn" (bottleneck) của số hóa theo cấp tầng và
  cường độ xuất khẩu.
\end{itemize}

\textbf{Tham chiếu}: Verhoef et al. (2021); Bharadwaj et al. (2013);
Brynjolfsson \& McAfee (2014); Mar et al. (2026, P3 Singapore
manuscript, \emph{MIR} under review); Đỗ \& Phan (2026, P4 Việt Nam ---
\emph{APJM}); Coase (1937); Williamson (1985); Banalieva \& Dhanaraj
(2019).

\begin{center}\rule{0.5\linewidth}{0.5pt}\end{center}

\hypertarget{chux1b0ux1a1ng-3--khung-phuxe2n-loux1ea1i-6-phuxe2n-nhuxf3m-con-icrv}{%
\subsection{CHƯƠNG 3 --- KHUNG PHÂN LOẠI 6 PHÂN NHÓM CON
ICRV}\label{chux1b0ux1a1ng-3--khung-phuxe2n-loux1ea1i-6-phuxe2n-nhuxf3m-con-icrv}}

\hypertarget{31-luxfd-do-phuxe2n-loux1ea1i}{%
\subsubsection{3.1 Lý do phân
loại}\label{31-luxfd-do-phuxe2n-loux1ea1i}}

Châu Á không phải là một khối đồng nhất về thể chế. Các nghiên cứu
meta-analysis (Marano et al., 2016; Wu et al., 2022) đã ghi nhận rằng
bối cảnh thể chế là biến điều tiết quan trọng nhất trong quan hệ I→P ---
nhưng không có nghiên cứu nào thiết lập hệ thống phân loại thể chế toàn
diện cho 47 nền kinh tế châu Á. Khung phân loại ICRV (Institutional
Context Regime Variation) được đề xuất trong chuyên đề này nhằm lấp
khoảng trống đó, dựa trên bốn nguyên tắc:

(a) \textbf{Tính phân biệt} (discriminating power): mỗi nhóm phải có mô
hình hiệu quả doanh nghiệp khác biệt thực sự --- không chỉ khác nhau về
mức GDP per capita;

(b) \textbf{Tính đại diện lý thuyết}: phân loại phải phản ánh được các
cơ chế thể chế khác nhau (thể chế đổi mới sáng tạo vs thể chế tài nguyên
vs thể chế chuyển đổi, v.v.) theo Varieties of Capitalism (Hall \&
Soskice, 2001) và Institutional Theory (North, 1990);

(c) \textbf{Tính khả dụng dữ liệu}: mỗi nhóm phải có đủ quan sát WBES để
phân tích có ý nghĩa thống kê (tối thiểu 500 doanh nghiệp);

(d) \textbf{Tính minh bạch và tái tạo}: tiêu chí phân loại phải rõ ràng
và có thể áp dụng nhất quán.

\hypertarget{32-tiuxeau-chuxed-phuxe2n-loux1ea1i-suxe1u-nhuxf3m}{%
\subsubsection{3.2 Tiêu chí phân loại sáu
nhóm}\label{32-tiuxeau-chuxed-phuxe2n-loux1ea1i-suxe1u-nhuxf3m}}

\textbf{Nhóm I --- Advanced đổi mới sáng tạo dẫn dắt
(Innovation-driven)}: WGI Rule of Law \textgreater+0,80; GNI per capita
\textgreater\$30.000 (WB Atlas method); GDP tăng trưởng dựa trên đổi mới
công nghệ, dịch vụ tài chính và hàm lượng tri thức cao. Bằng chứng WBES:
R\&D \textgreater10\%, ISO \textgreater20\%, sd log năng suất
\textasciitilde1,03 (diversified).

\textbf{Nhóm II --- Advanced tài nguyên dẫn dắt (Resource-driven)}: WGI
Rule of Law trung bình (+0,00 dẫn đến +0,50); GNI per capita
\textgreater\$20.000; GDP tăng trưởng dựa trên xuất khẩu tài nguyên
thiên nhiên (dầu mỏ, khí đốt, khoáng sản). Bằng chứng WBES: FDI cao
hướng khai thác tài nguyên, sd log năng suất thấp (\textasciitilde0,49
--- đặc trưng nhà nước tô rentier state).

\textbf{Nhóm III --- Trung bình cao (Upper-middle)}: WGI Rule of Law
0,00 dẫn đến +0,80; GNI per capita \$5.000--\$20.000; nền kinh tế chuyển
đổi từ manufacturing sang dịch vụ, với R\&D đang tăng (đặc biệt Trung
Quốc). WBES pattern: manufacturing FSTS cao, R\&D tăng mạnh giai đoạn
2018--2025 (+21,5 đpt Upper-middle).

\textbf{Nhóm IV --- Đang nổi (Emerging)}: WGI Rule of Law -0,50 dẫn đến
0,00; GNI per capita \$1.000--\$5.000; tốc độ tăng trưởng cao nhưng gap
de jure--de facto lớn (Xu, 2024); manufacturing FDI-driven dẫn dắt xuất
khẩu. WBES pattern: FSTS phân cực cao/thấp trong nội bộ nhóm; FDI thấp
nhất (4,7\%) --- do doanh nghiệp nội địa chi phối sample.

\textbf{Nhóm V --- Cận biên (Frontier)}: WGI Rule of Law \textless{}
-0,50; GNI per capita \textless\$1.000 hoặc \textless\$3.000 với voids
thể chế lớn; nhiều nước đang qua giai đoạn xung đột hoặc chuyển đổi sau
xung đột. WBES pattern: phân tán năng suất cao nhất (sd log 1,36), R\&D
thấp nhất, nhưng đổi mới sản phẩm cao (bricolage innovation).

\textbf{Nhóm VI --- SIDS Thái Bình Dương (Boundary case extension)}: Dân
số \textless1 triệu; cô lập địa lý \textgreater2.000 km từ hub thương
mại; phụ thuộc ODA + remittances \textgreater{} 20\% GDP; MIRAB model
(Bertram, 2006). WBES pattern: n=1.371 (1,4\% pool); FDI cao nhất
(23,5\% --- do MNE tourism); innovation product cao nhất (41,5\% ---
bricolage); website cao trừ Kiribati (18,7\%).

\hypertarget{33-danh-suxe1ch-47-nux1ec1n-kinh-tux1ebf-theo-nhuxf3m-icrv-bux1ea3ng-31}{%
\subsubsection{3.3 Danh sách 47 nền kinh tế theo nhóm ICRV (Bảng
3.1)}\label{33-danh-suxe1ch-47-nux1ec1n-kinh-tux1ebf-theo-nhuxf3m-icrv-bux1ea3ng-31}}

\textbf{Bảng 3.1}. \emph{Phân loại 47 nền kinh tế vào 6 nhóm ICRV ---
tiêu chí và n\_firms.}

\begin{longtable}[]{@{}llll@{}}
\toprule\noalign{}
Nhóm ICRV & Tên nhóm & Quốc gia (ISO3) & n\_firms (ước tính) \\
\midrule\noalign{}
\endhead
\bottomrule\noalign{}
\endlastfoot
\textbf{I} & Advanced --- innovation-driven & SGP, HKG, KOR, TWN, ISR &
\textasciitilde4.220 \\
\textbf{II} & Advanced --- resource-driven & SAU, QAT, KWT, BHR, BRN &
\textasciitilde1.932 \\
\textbf{III} & Upper-middle & CHN, MYS, THA, KAZ, ARM, GEO &
\textasciitilde16.693 \\
\textbf{IV} & Emerging & VNM, IDN, PHL, IND, LKA, JOR, MNG &
\textasciitilde47.803 \\
\textbf{V} & Frontier & BGD, PAK, LAO, KHM, MMR, NPL, BTN, MDV, UZB,
TJK, KGZ, TKM, AFG, TLS, IRQ, LBN, YEM & \textasciitilde28.678 \\
\textbf{VI} & SIDS Pacific (boundary) & FJI, PNG, SLB, TON, VUT, WSM,
KIR & \textasciitilde1.371 \\
\textbf{Tổng} & & 47 nền kinh tế & \textbf{101.185} \\
\end{longtable}

\emph{Ghi chú: n\_firms là phân bố xấp xỉ theo pool 101.185; tổng 5 nhóm
chính (I--V) = 99.814 (scope chính 41 nước châu Á); SIDS Nhóm VI = 1.371
(1,4\% pool).}

\hypertarget{34-ux111ux1eb7c-ux111iux1ec3m-cuxe1c-nhuxf3m-icrv-bux1ea3ng-32}{%
\subsubsection{3.4 Đặc điểm các nhóm ICRV (Bảng
3.2)}\label{34-ux111ux1eb7c-ux111iux1ec3m-cuxe1c-nhuxf3m-icrv-bux1ea3ng-32}}

\textbf{Bảng 3.2}. \emph{Đặc điểm điển hình của 6 nhóm ICRV --- WGI,
GNI, cơ chế tăng trưởng, pattern WBES dự kiến.}

\begin{longtable}[]{@{}lllllll@{}}
\toprule\noalign{}
Nhóm & WGI Rule of Law & GNI/capita & Cơ chế tăng trưởng & FSTS dự kiến
& sd log LP & Cơ chế I→P \\
\midrule\noalign{}
\endhead
\bottomrule\noalign{}
\endlastfoot
I --- Innovation & \textgreater+0,80 & \textgreater\$30k & R\&D + dịch
vụ tri thức & Trung bình (6--15\%) & \textasciitilde1,03 & Inverted-U
nhẹ / tuyến tính dương; TP \textasciitilde82\% (SG) \\
II --- Resource & +0,00→+0,50 & \textgreater\$20k & Xuất khẩu dầu/khí &
Thấp (1--3\%) & \textasciitilde0,49 & Phi tuyến suy yếu; rentier state
redistribution \\
III --- Upper-mid & 0,00→+0,80 & \$5k--\$20k & Manufacturing + chuyển
đổi & Trung bình (8--12\%) & \textasciitilde1,29 & Inverted-U; TP
\textasciitilde47\% (CHN Đỗ \& Phan, 2026 JFAR) \\
IV --- Emerging & -0,50→0,00 & \$1k--\$5k & FDI-driven manufacturing &
Biến động (7--23\%) & \textasciitilde1,24 & Inverted-U; TP 39--46\%
(VNM) --- bẫy RDT \\
V --- Frontier & \textless-0,50 & \textless\$1k--\$3k & Subsistence +
informal & Thấp--trung bình & \textasciitilde1,36 & Yếu; institutional
voids → bricolage innovation \\
VI --- SIDS & Đặc biệt & \$1k--\$8k & MIRAB + tourism & Rất thấp
(1--12\%) & \textasciitilde1,32 & Forced internationalization ---
U-curve không vận hành \\
\end{longtable}

\emph{Nguồn: Tổng hợp từ World Bank Enterprise Surveys (2025); WGI từ
Kaufmann et al. (2011); GNI từ World Bank (2025, July). Pattern WBES từ
phân tích Chương 4 và Đỗ \& Phan (2026 --- VEFR, APJM, JFAR, MIR under
review).}

\begin{center}\rule{0.5\linewidth}{0.5pt}\end{center}

\begin{center}\rule{0.5\linewidth}{0.5pt}\end{center}

\hypertarget{phux1ee5-lux1ee5c-h--bux1ea3ng-thuux1eadt-ngux1eef-anh-viux1ec7t-ruxfat-gux1ecdn}{%
\subsection{PHỤ LỤC H --- BẢNG THUẬT NGỮ ANH-VIỆT (RÚT
GỌN)}\label{phux1ee5-lux1ee5c-h--bux1ea3ng-thuux1eadt-ngux1eef-anh-viux1ec7t-ruxfat-gux1ecdn}}

Phụ lục H trình bày bản đối chiếu Việt -- Anh -- Hán-Việt cho các thuật
ngữ học thuật chuyên ngành sử dụng trong chuyên đề. Sáu nhóm thuật ngữ:
(i) khái niệm hiệu quả; (ii) thước đo năng suất; (iii) khung lý thuyết;
(iv) chỉ số WBES; (v) phương pháp thống kê; (vi) thể chế. Chi tiết bản
đầy đủ xem \texttt{thesis/09b\_vn\_term\_glossary.md}.

\begin{center}\rule{0.5\linewidth}{0.5pt}\end{center}

\emph{Tiếp tục ở Phần 2 (Chương 4 --- thực trạng từ dữ liệu WBES với
pool 101.185 doanh nghiệp) trong file
\texttt{thesis/15\_cd1\_part2\_findings\_vi.md}.}

\hypertarget{chuyuxean-ux111ux1ec1-tiux1ebfn-sux129-sux1ed1-1-bux1ea3n-nhuxe1p-ux111ux1ea7y-ux111ux1ee7-phux1ea7n-2-chux1b0ux1a1ng-4-thux1ef1c-trux1ea1ng-tux1eeb-wbes}{%
\section{CHUYÊN ĐỀ TIẾN SĨ SỐ 1, BẢN NHÁP ĐẦY ĐỦ (PHẦN 2: CHƯƠNG 4, THỰC
TRẠNG TỪ
WBES)}\label{chuyuxean-ux111ux1ec1-tiux1ebfn-sux129-sux1ed1-1-bux1ea3n-nhuxe1p-ux111ux1ea7y-ux111ux1ee7-phux1ea7n-2-chux1b0ux1a1ng-4-thux1ef1c-trux1ea1ng-tux1eeb-wbes}}

\begin{quote}
Tiếp nối \texttt{thesis/14\_cd1\_part1\_intro\_theory\_vi.md}. Phần 3
(Chương 5--7 + TLTK):
\texttt{thesis/16\_cd1\_part3\_cases\_conclusion\_vi.md}. Bảng thuật ngữ
Anh-Việt: \texttt{thesis/09b\_vn\_term\_glossary.md}. Hình minh họa:
\texttt{thesis/figures/} (\textbf{11 hình}; chạy
\texttt{python3\ generate\_figures.py} để regen).
\end{quote}

\begin{center}\rule{0.5\linewidth}{0.5pt}\end{center}

\hypertarget{chux1b0ux1a1ng-4--thux1ef1c-trux1ea1ng-hiux1ec7u-quux1ea3-doanh-nghiux1ec7p-chuxe2u-uxe1-20092025}{%
\subsection{CHƯƠNG 4 --- THỰC TRẠNG HIỆU QUẢ DOANH NGHIỆP CHÂU Á
2009--2025}\label{chux1b0ux1a1ng-4--thux1ef1c-trux1ea1ng-hiux1ec7u-quux1ea3-doanh-nghiux1ec7p-chuxe2u-uxe1-20092025}}

\hypertarget{41-nguux1ed3n-dux1eef-liux1ec7u-world-bank-enterprise-surveys}{%
\subsubsection{4.1 Nguồn dữ liệu World Bank Enterprise
Surveys}\label{41-nguux1ed3n-dux1eef-liux1ec7u-world-bank-enterprise-surveys}}

\textbf{Phạm vi tổng hợp dữ liệu}. Nhóm dữ liệu (pool) gồm
\textbf{101.185 doanh nghiệp}, 47 nền kinh tế, \textbf{108 cặp quốc gia
× năm}, giai đoạn 2009--2025. Tổng hợp này kế thừa và mở rộng từ 17 nước
châu Á mới nổi (\textasciitilde40.633 doanh nghiệp) của Đỗ \& Phan (2026
--- VEFR), gấp \textasciitilde2,5 lần. Phân bố theo phân nhóm con
(sub-regime) ICRV: Emerging 47.803 (47\%), Frontier 28.678 (28\%),
Upper-middle 16.693 (17\%), Advanced 6.640 (7\%), \textbf{SIDS 1.371
(1,4\%, gồm Kiribati 2025)}. Có \textbf{14 đợt khảo sát năm 2025} với
16.979 doanh nghiệp.

\includegraphics{figures/fig_4_1_pool_composition.png}

\emph{Hình 4.1.1. Bar chart phân bố theo 5 phân nhóm con: Emerging
(47.803, 47,2\%), Frontier (28.678, 28,3\%), Upper-middle (16.693,
16,5\%), Advanced (6.640, 6,6\%), SIDS Thái Bình Dương (1.371, 1,4\%,
gồm Kiribati 2025). Tái lập:
\texttt{thesis/figures/generate\_figures.py} function
\texttt{fig\_4\_1\_pool\_composition()}.}

\includegraphics{figures/fig_4_0_pool_by_year.png}

\emph{Hình 4.0. Phân bố mẫu doanh nghiệp xuyên 17 năm theo 3 thế hệ
schema WBES: PICS3 (2009-2012, n=14.171), Standardized (2013-2017,
n=24.564), BREADY (2018-2025, n=62.450 gồm Kiribati 2025). Đợt 2025 đột
biến với 16.979 doanh nghiệp = 16,8\% pool --- đợt khảo sát đơn năm lớn
nhất. Tái lập: \texttt{generate\_figures.py} function
\texttt{fig\_4\_0\_pool\_by\_year()}.}

\textbf{Trường hợp biên (boundary cases)}: \textbf{7 SIDS Thái Bình
Dương đầy đủ (FJI, PNG, SLB, TON, VUT, WSM, KIR)} + Tây Á 9 nước. Phân
bố thời gian: 2009--2012 (n=14.171), 2013--2017 (n=24.564), 2018--2025
(n=62.450, chiếm 62\%). Ba thế hệ khung dữ liệu (schema):
PICS3/MENA-WBES, Standardized, Standardized2018+/BREADY.

\textbf{Hài hòa}. Pipeline Python (\texttt{wbes/02\_harmonize.py}); FSTS
= \texttt{d3b\ +\ d3c}; giới hạn cực trị (winsorize) log năng suất ở mức
1/99 phân vị trong cụm quốc gia × năm. Doanh thu chưa quy đổi sang USD
PPP.

\hypertarget{42-thux1ef1c-trux1ea1ng-nux103ng-suux1ea5t-lao-ux111ux1ed9ng-phuxe2n-tuxe1n-trong-tux1eebng-quux1ed1c-gia}{%
\subsubsection{4.2 Thực trạng năng suất lao động, phân tán trong từng
quốc
gia}\label{42-thux1ef1c-trux1ea1ng-nux103ng-suux1ea5t-lao-ux111ux1ed9ng-phuxe2n-tuxe1n-trong-tux1eebng-quux1ed1c-gia}}

\textbf{Bảng 4.1}. \emph{Phân tán (dispersion) năng suất lao động theo
phân nhóm con thể chế (n=108 cặp QG×năm, v3.1; nguồn dữ liệu: P7 pooled,
dispersion tính trên ln(LP) đã loại trừ trung bình quốc gia × năm).}

\begin{longtable}[]{@{}llllll@{}}
\toprule\noalign{}
Phân nhóm con & Cặp QG×năm & n\_firms & sd log & P90/P10 & P75/P25 \\
\midrule\noalign{}
\endhead
\bottomrule\noalign{}
\endlastfoot
\textbf{Advanced, innovation-driven} \emph{(SG, HK, KOR, TWN, ISR, CYP)}
& 8 & 3.881 & \textbf{1,13} & \textbf{14,05} & \textbf{3,55} \\
\textbf{Advanced, resource-driven} \emph{(SAU, QAT, KWT, BHR, BRN)} & 5
& 1.810 & \textbf{0,72} & \textbf{3,99} & \textbf{1,76} \\
\emph{Advanced (gộp, tham chiếu)} & 13 & 5.691 & 1,01 & 10,8 & 3,1 \\
Upper-middle & 18 & 15.174 & 1,29 & 27,7 & 5,4 \\
Emerging & 20 & 45.388 & 1,24 & 30,6 & 5,1 \\
Frontier & 42 & 18.877 & 1,36 & 39,6 & 6,1 \\
\textbf{SIDS Thái Bình Dương (v3.1: gồm Kiribati 2025)} & \textbf{13} &
\textbf{1.352} & \textbf{1,50} & \textbf{32,62} & \textbf{6,15} \\
\end{longtable}

\emph{Ghi chú: Advanced tách 2 phân nhóm con ở v3.1; tỷ số phân tán tiên
tiến đổi mới sáng tạo (sd log 1,13) so với tiên tiến tài nguyên dẫn dắt
(sd log 0,72) ≈ \textbf{1,6 lần}. SIDS Thái Bình Dương có phân tán cao
nhất pool (P90/P10 = 32,62), phù hợp với pattern survival-based
internationalization (Bertram, 2006). Dispersion tính trên ln(LP)
demeaned theo cặp quốc gia × năm để loại trừ chênh lệch tiền tệ danh
nghĩa xuyên nước. Nguồn: P7 pooled clean (49 nền kinh tế; Đỗ \& Phan,
2026, đang revision JIBS).}

\includegraphics{figures/fig_4_2_productivity_density.png}

\emph{Hình 4.2. Đồ thị mật độ kernel cho 6 phân nhóm con: Advanced
innovation (sd 1,03), Advanced resource (sd 0,49), Upper-middle (1,29),
Emerging (1,24), Frontier (1,36), SIDS (1,32). Phân nhóm tài nguyên dẫn
dắt (Vùng Vịnh) phân tán hẹp nhất; Frontier rộng nhất. Tái lập:
\texttt{generate\_figures.py} function
\texttt{fig\_4\_2\_productivity\_density()}.}

5 phát hiện: (1) Phân tán Advanced gộp giảm từ 1,00 xuống 0,86 sau khi
bổ sung Vùng Vịnh, bằng chứng dị biệt nội bộ; (2) Frontier cao nhất, phù
hợp \textbf{giả thuyết phân bổ sai nguồn lực (misallocation hypothesis)}
Hsieh \& Klenow (2009, 2014); (3) SIDS ở mức trung bình
\textasciitilde1,32 (v3.1); (4) tỷ số P90/P10 tăng đơn điệu theo phân
nhóm con; (5) bằng chứng cho H5 --- điều tiết thể chế (institutional
moderation).

\hypertarget{43-thux1ef1c-trux1ea1ng-quux1ed1c-tux1ebf-huxf3a-vuxe0-tux103ng-trux1b0ux1edfng-viux1ec7c-luxe0m}{%
\subsubsection{4.3 Thực trạng quốc tế hóa và tăng trưởng việc
làm}\label{43-thux1ef1c-trux1ea1ng-quux1ed1c-tux1ebf-huxf3a-vuxe0-tux103ng-trux1b0ux1edfng-viux1ec7c-luxe0m}}

\textbf{Bảng 4.3}. \emph{FSTS, doanh nghiệp xuất khẩu và CAGR việc làm
theo phân nhóm con.}

\begin{longtable}[]{@{}llll@{}}
\toprule\noalign{}
Phân nhóm con & FSTS (\%) & Doanh nghiệp xuất khẩu (\%) & CAGR việc làm
(\%) \\
\midrule\noalign{}
\endhead
\bottomrule\noalign{}
\endlastfoot
Advanced & 10,2 & 23,0 & 3,15 \\
Upper-middle & 10,3 & 21,7 & 4,25 \\
Emerging & 8,6 & 15,5 & 2,81 \\
Frontier & 10,1 & 16,6 & 3,65 \\
SIDS & 6,3 & 16,3 & 5,77 \\
\end{longtable}

Trung vị FSTS bằng 0\%, phân phối phân cực mạnh; SIDS có CAGR việc làm
cao nhất.

\hypertarget{44-thux1ef1c-trux1ea1ng-ux111ux1ed5i-mux1edbi-suxe1ng-tux1ea1o-vuxe0-nux103ng-lux1ef1c-sux1ed1}{%
\subsubsection{4.4 Thực trạng đổi mới sáng tạo và năng lực
số}\label{44-thux1ef1c-trux1ea1ng-ux111ux1ed5i-mux1edbi-suxe1ng-tux1ea1o-vuxe0-nux103ng-lux1ef1c-sux1ed1}}

\hypertarget{440-khung-5-lux129nh-vux1ef1c-chux1ec9-sux1ed1-wbes--ux111ux1ecbnh-vux1ecb-mux1ee5c-44-trong-tux1ed5ng-thux1ec3}{%
\paragraph{4.4.0 Khung 5 lĩnh vực chỉ số WBES --- định vị Mục 4.4 trong
tổng
thể}\label{440-khung-5-lux129nh-vux1ef1c-chux1ec9-sux1ed1-wbes--ux111ux1ecbnh-vux1ecb-mux1ee5c-44-trong-tux1ed5ng-thux1ec3}}

\begin{quote}
\textbf{Nguồn chuẩn}: World Bank (n.d.). \emph{Tài liệu tóm tắt: Các chỉ
số khảo sát doanh nghiệp Enterprise Surveys} (file 04 v2.5 Section L).
Cross-cite glossary v1.2 Nhóm 6 (commit e9a73ae) + Nhóm 7 (commit
beb47d3).
\end{quote}

WBES tổ chức các chỉ số doanh nghiệp thành \textbf{5 lĩnh vực chính thức
(5 indicator domains)} --- đây là khung phân tích chuẩn của World Bank
được dùng nhất quán xuyên các báo cáo WBES châu Á và châu lục khác. Bảng
dưới định vị Mục 4.4 hiện tại trong tổng thể 5 lĩnh vực:

\textbf{Bảng 4.4.0}. \emph{Khung 5 lĩnh vực chỉ số WBES --- định vị phạm
vi đo lường của Mục 4.4.}

\begin{longtable}[]{@{}lllll@{}}
\toprule\noalign{}
\# & Lĩnh vực & Ví dụ chỉ số WBES & Mục CĐ1 đo trực tiếp & Hàm ý CĐ2 \\
\midrule\noalign{}
\endhead
\bottomrule\noalign{}
\endlastfoot
(1) & \textbf{Quy định pháp lý + Thuế} & Thuế thời gian (Time Tax),
h/năm tuân thủ thuế; Tuân thủ thuế (Tax compliance); Cấp phép kinh
doanh; Thanh tra & (gợi ý CĐ2) & CĐ2 sẽ deep-dive, biến
\texttt{time\_tax\_hours} và \texttt{tax\_inspection\_freq} \\
(2) & \textbf{Tài chính + Tín dụng} & Hạn chế tín dụng (Credit
constraint); Số hóa tài chính (Financial digitalization); Tỷ lệ DN có
tài khoản ngân hàng; Tỷ lệ DN có khoản vay & (gợi ý CĐ2, link Mục 4.5.5
rào cản \#1) & CĐ2 deep-dive, biến \texttt{credit\_constraint\_dummy} \\
(3) & \textbf{Cơ sở hạ tầng + Khí hậu} & Quản lý năng lượng (Energy
management); Mất điện (Power outages); Theo dõi CO2 (CO2 monitoring);
Chất lượng đường giao thông & \textbf{Mục 4.4 đo gián tiếp qua quy trình
mới}; \textbf{Mục 4.5.5 rào cản \#4 đo trực tiếp} & CĐ2, biến
\texttt{power\_outage\_intensity} từ question c30 \\
(4) & \textbf{Thương mại + Cạnh tranh} & Thời gian thông quan (Customs
clearance time); Đối thủ phi chính thức; Xuất khẩu trực tiếp/gián tiếp;
FSTS & (FSTS đã đo Mục 4.3); cạnh tranh phi chính thức ở Mục 4.5.5 \#3 &
CĐ2, biến \texttt{customs\_days} + \texttt{informal\_competition} \\
(5) & \textbf{Tham nhũng + Phi chính thức} & Hối lộ (Bribery): Bribery
Incidence + Bribery Depth; Tham ô (Graft Index); Đối thủ không chính
thức & (sẽ phân biệt rõ trong Mục 4.7.5 v3.8c, NEW) & CĐ2 deep-dive,
biến \texttt{bribery\_incidence\_pct} \\
\end{longtable}

\textbf{Phạm vi đo lường của Mục 4.4} \emph{(đổi mới sáng tạo + năng lực
số)}: Chỉ tập trung vào \textbf{lĩnh vực (5) phần đổi mới + năng lực số}
(gồm Sản phẩm mới, Quy trình mới, R\&D, ISO, Website) và một phần
\textbf{lĩnh vực (3) hạ tầng} (Website đại diện cho ICT infrastructure
access). Các lĩnh vực (1) (2) (4) và phần lớn (5) tham nhũng được
\textbf{đề cập gián tiếp} qua bảng kết quả mức độ chi tiết quốc gia (Mục
4.10) hoặc rào cản hàng đầu (Mục 4.5.5), nhưng chưa được phân tích sâu
--- đây là \textbf{scope mở rộng cho CĐ2} với 5 control variables mới
(xem Mục 4.5.5 hàm ý CĐ2).

\textbf{Lý do giới hạn phạm vi CĐ1}: (a) tính khả dụng của biến, biến
trong (1) (2) (4) có nhiều missing values trong PICS3 schema cũ
2009-2012, gây mất cân đối panel; (b) tập trung vào hai construct cốt
lõi của luận án (TCI + DAI), Đỗ \& Phan (2026 --- VEFR) đã thiết kế hai
biến này từ lĩnh vực (5) đổi mới + năng lực số; (c) CĐ1 mang tính
descriptive và gợi mở, để dành CĐ2 cho phân tích đa biến với
specification đầy đủ 5 lĩnh vực.

\textbf{Bảng 4.4}. \emph{Đổi mới sáng tạo và áp dụng số (\%).}

\begin{longtable}[]{@{}llllll@{}}
\toprule\noalign{}
Phân nhóm con & Sản phẩm mới & Quy trình mới & R\&D & ISO & Website \\
\midrule\noalign{}
\endhead
\bottomrule\noalign{}
\endlastfoot
Advanced & 22,3 & 52,3 & 16,7 & 29,9 & 59,3 \\
Upper-middle & 26,7 & 71,7 & 21,0 & 31,4 & 56,9 \\
Emerging & 17,5 & 65,2 & 16,4 & 24,9 & 49,2 \\
Frontier & 23,1 & 68,9 & 14,2 & 20,7 & 38,0 \\
SIDS Thái Bình Dương & \textbf{41,5} & 65,1 & 11,8 & 16,5 &
\textbf{58,9} \\
\end{longtable}

SIDS Thái Bình Dương thể hiện pattern \textbf{thích nghi và nhảy vọt số
(adaptation + digital leapfrog)}; phân tách rõ Năng lực công nghệ (TCI)
so với Năng lực số (DAI), kế thừa Đỗ \& Phan (2026, VEFR).

\hypertarget{445-tuxe1i-ux111ux1ecbnh-huxecnh-nux103ng-lux1ef1c-sux1ed1--tier-1-digital-presence-rebrand}{%
\paragraph{4.4.5 Tái định hình "Năng lực số" --- Tier-1 Digital Presence
rebrand}\label{445-tuxe1i-ux111ux1ecbnh-huxecnh-nux103ng-lux1ef1c-sux1ed1--tier-1-digital-presence-rebrand}}

Spec 1 (full coverage 2009--2025) dùng biến nhị phân duy nhất
(\texttt{website} Y/N) đại diện cho "Năng lực số (DAI)". Trong khi đó
TCI được đo bằng nhiều thành phần vững chắc (R\&D + ISO multi-component
aspirational). \textbf{Hậu quả phân tích}: hệ số -0,129 cho DAI ở
Advanced trong Spec 1 gợi ý "số hóa làm giảm năng suất ở Singapore/Hàn
Quốc/Hong Kong", nhưng đây là \textbf{ảo ảnh thống kê} do đo lường sơ
sài, không phải thực tế kinh tế. Như so sánh: \emph{"Giống như cố đánh
giá hiệu suất của mạng lưới logistic AI hiện đại nhưng dùng bản đồ trạm
điện tín thế kỷ 19."}

\textbf{Đề xuất rebrand v3.5}: Đổi tên biến trong Spec 1 từ "DAI / Năng
lực số / Chuyển đổi số" sang \textbf{"Sự hiện diện số cơ bản (Tier-1
Digital Presence)"}, phạm vi hẹp hơn, khiêm tốn hơn, thừa nhận rằng Spec
1 chỉ đo \emph{"bức số hóa tài liệu"} (website binary), KHÔNG đo thay
đổi mô hình kinh doanh hay digital transformation.

\textbf{ICT exclusion test (promoted từ Mục 4.8 I2 lên thân Chương 4)}:
Khi loại bỏ ICT firms (mã ISIC J 58--63) khỏi Advanced phân nhóm
innovation-driven, hệ số âm DAI (-0,129) \textbf{biến mất} dẫn đến bằng
chứng cho \textbf{bão hòa Tier-1 ở các nước tiên tiến} (Sự hiện diện số
cơ bản đã saturated từ trước 2018 ở Singapore/Hong Kong/Hàn Quốc/Đài
Loan/Israel, Đỗ \& Phan, 2026, VEFR p. 24). Việc đưa ICT exclusion test
lên mạch chính (thay vì đẩy xuống phụ lục/CĐ2) \textbf{không chỉ sửa lỗi
thống kê mà còn tạo ra lập luận sắc bén}: \emph{"Sự hiện diện số cơ bản
đã bão hòa ở các nước tiên tiến từ lâu, vì vậy biến này không còn
discriminate năng suất giữa các firms ở Advanced."}

\textbf{Spec 1 vs Spec 2, biểu đồ đối chiếu hai trục thời gian}
\emph{(đề xuất Hình 4.X mới, đang chuẩn bị)}: Spec 1 (single var website
binary, 2009--2025) vs Spec 2 (composite DAI 5 thành phần: website +
email + smartphone POS + e-commerce + cloud, 2018--2025) trên cùng trục
thời gian, chứng minh hệ sinh thái số chuyển từ \textbf{trạng thái tĩnh
"Tier-1 Digital Presence"} (chỉ documents online) sang \textbf{trạng
thái động "Tier-2 Digital Transformation"} (business models reshaped, kế
thừa Banalieva \& Dhanaraj, 2019 layered framework).

\hypertarget{4451-ux111ux1ed1i-chiux1ebfu-dai-xuyuxean-quux1ed1c-gia-singapore-tier-12-vs-viux1ec7t-nam-tier-1-only}{%
\paragraph{4.4.5.1 Đối chiếu DAI xuyên quốc gia. Singapore (Tier 1+2) vs
Việt Nam (Tier 1
only)}\label{4451-ux111ux1ed1i-chiux1ebfu-dai-xuyuxean-quux1ed1c-gia-singapore-tier-12-vs-viux1ec7t-nam-tier-1-only}}

\begin{quote}
\textbf{Bằng chứng nền tảng}: Mục 2.7 file 14 v3.13 (Khung 4-Tier
Verhoef + CDCM); P3 Singapore manuscript (Mar et al., 2026, \emph{MIR});
P4 Việt Nam (Đỗ \& Phan, 2026, \emph{APJM}).
\end{quote}

Phân tích đối chiếu nội bộ chỉ ra \textbf{mâu thuẫn tử huyệt} giữa hai
papers: Singapore (Tier 1+2 = website + e-payment) cho hệ số tương tác
DAI×FSTS² \textbf{dương mạnh} (β=3,119, p=,005), trong khi Việt Nam
(Tier 1 = website only) cho tương tác DAI×FSTS \textbf{âm} ở wave 2023
(β=-0,912, p=,043). Đây không phải mâu thuẫn ngẫu nhiên mà là
\textbf{bằng chứng cho khung CDCM}, giá trị của công cụ số phụ thuộc độ
tương thích giữa cấp độ số hóa và mật độ giao dịch quốc tế:

\textbf{Bảng 4.4.5.1}. \emph{Đối chiếu cơ chế DAI giữa hai bối cảnh thể
chế (CDCM application).}

\begin{longtable}[]{@{}lll@{}}
\toprule\noalign{}
Chiều & Singapore (Advanced innovation, 2023) & Việt Nam (Emerging,
2009-2023) \\
\midrule\noalign{}
\endhead
\bottomrule\noalign{}
\endlastfoot
\textbf{Cấp DAI} & Tier 1+2 (website + cường độ thanh toán điện tử) &
Tier 1 (website only) \\
\textbf{Mẫu} & n=623, cắt ngang 1 năm & n=2.958 firms × 3 wave \\
\textbf{DAI direct effect} & β=0,168, p\textless,001 (dương ổn định) &
Stage-contingent: 2009 mạnh → 2015 null (β=-0,044, p=,377) → 2023 phục
hồi \\
\textbf{DAI × FSTS² interaction} & \textbf{β=3,119, p=,005 (dương
mạnh)}, conditional scaling resource ở xuất khẩu cao & \textbf{β=-0,912,
p=,043 (âm) ở 2023}, Tier 1 trở thành điểm nghẽn ở xuất khẩu cao \\
\textbf{2SLS robustness (DAI)} & (chưa kiểm định IV) & β co về 0,02,
p=,94 --- \textbf{không xác lập nhân quả} \\
\textbf{Cơ chế giải thích} & Tier 1+2 hấp thụ super-linear coordination
cost xuất khẩu cao (Brynjolfsson \& McAfee, 2014) & Tier 1 không quản lý
nổi giao dịch lớn → khuếch đại quá tải thông tin \\
\end{longtable}

\textbf{Hàm ý cho lập luận Mục 4.4.5 hiện hành}: Hệ số âm DAI ở Advanced
trong Spec 1 (-0,129) không chỉ là \emph{artifact của Tier-1 saturation}
(đã lập luận v3.5) mà còn phản ánh \textbf{gãy đổ đo lường} (measurement
break) khi Tier 1 không đủ sức quản lý giao dịch ở xuất khẩu cao. Khi mở
rộng sang Tier 1+2 (Spec 2 với e-payment), hệ số dương mạnh emerges,
replication pattern Singapore P3.

\textbf{Đề xuất phương pháp luận cho CĐ2} \emph{(robustness check
\#10)}: Cross-validation construct boundary, chạy Spec 2 chỉ ở 41 nước
có đủ Tier 1+2 data (BREADY+ wave); confirm rằng pattern Singapore-style
(DAI × FSTS² dương) emerges khi có composite digital adoption metric.
Nếu confirm dẫn đến đóng góp lý thuyết mạnh cho CDCM (Đỗ \& Phan, 2026
mở rộng từ APJM sang luận án). Nếu không confirm dẫn đến cần Tier 3
(ERP/CRM) hoặc Tier 4 (AI) để nắm cơ chế thực sự, hàm ý cho data
collection tương lai (limitations Mục 7.3.4 file 16).

\hypertarget{45-thux1ef1c-trux1ea1ng-cux1ea5u-truxfac-doanh-nghiux1ec7p}{%
\subsubsection{4.5 Thực trạng cấu trúc doanh
nghiệp}\label{45-thux1ef1c-trux1ea1ng-cux1ea5u-truxfac-doanh-nghiux1ec7p}}

\textbf{Bảng 4.5}. \emph{Cấu trúc doanh nghiệp theo phân nhóm con (\%).}

\begin{longtable}[]{@{}llll@{}}
\toprule\noalign{}
Phân nhóm con & SME & Doanh nghiệp xuất khẩu & FDI ≥10\% \\
\midrule\noalign{}
\endhead
\bottomrule\noalign{}
\endlastfoot
Advanced & 79,1 & 23,0 & 11,1 \\
Upper-middle & 76,2 & 21,7 & 8,4 \\
Emerging & 74,4 & 15,5 & 4,7 \\
Frontier & 85,0 & 16,5 & 5,9 \\
SIDS Thái Bình Dương & \textbf{88,5} & 16,3 & \textbf{23,5} \\
\end{longtable}

Tỷ lệ FDI có dạng chữ U với cực tiểu ở Emerging. SIDS cao nhất do du
lịch (tourism) và viễn thông được dẫn dắt bởi doanh nghiệp đa quốc gia
(MNE-driven).

\includegraphics{figures/fig_4_6_firm_size.png}

\emph{Hình 4.6. Stacked bar chart cho thấy SIDS Pacific có tỷ lệ SME cao
nhất (88,5\%), Frontier kế tiếp (85,0\%), Advanced thấp nhất (79,1\%).
Pattern phù hợp với cấu trúc kinh tế khu vực. SIDS thị trường nhỏ phụ
thuộc SME; Advanced có nhiều doanh nghiệp lớn. Tái lập:
\texttt{generate\_figures.py} function
\texttt{fig\_4\_6\_firm\_size()}.}

\hypertarget{455-bux1ed1n-ruxe0o-cux1ea3n-huxe0ng-ux111ux1ea7u-cux1ee7a-khu-vux1ef1c-tux1b0-nhuxe2n-uxe1-thuxe1i-bux1eb1ng-chux1ee9ng-tux1eeb-wbes}{%
\paragraph{4.5.5 Bốn rào cản hàng đầu của khu vực tư nhân Á-Thái, bằng
chứng từ
WBES}\label{455-bux1ed1n-ruxe0o-cux1ea3n-huxe0ng-ux111ux1ea7u-cux1ee7a-khu-vux1ef1c-tux1b0-nhuxe2n-uxe1-thuxe1i-bux1eb1ng-chux1ee9ng-tux1eeb-wbes}}

\begin{quote}
\textbf{Nguồn chuẩn}: World Bank (n.d.). \emph{Tài liệu tóm tắt: Các chỉ
số khảo sát doanh nghiệp Enterprise Surveys} (file 04 v2.5 Section L).
Cross-cite glossary v1.2 Nhóm 7 \#61 (commit beb47d3).
\end{quote}

WBES chính thức xác định \textbf{4 rào cản hàng đầu (top 4 barriers)}
đối với khu vực tư nhân ở các nền kinh tế đang phát triển và mới nổi
Á-Thái Bình Dương, phù hợp với pool 101.185 doanh nghiệp của chuyên đề
này. Bốn rào cản này tạo thành \textbf{khung phân tích cấu trúc}
(structural framework) bổ sung cho 5 lĩnh vực chỉ số WBES được phân tích
trong Mục 4.4 (đổi mới + năng lực số) và Mục 4.5 (cấu trúc doanh
nghiệp). Bảng phân tầng theo 4 phân nhóm con thể chế:

\textbf{Bảng 4.5.5}. \emph{Bốn rào cản hàng đầu Asia-Pacific WBES, phân
tầng theo phân nhóm con thể chế (cảm quan định tính + vào sâu chương
trình CĐ2).}

\begin{longtable}[]{@{}llllllll@{}}
\toprule\noalign{}
\# & Rào cản & Tiếng Anh chuẩn WB & Cường độ Frontier & Cường độ
Emerging & Cường độ Upper-middle & Cường độ Advanced & Lý thuyết liên
kết \\
\midrule\noalign{}
\endhead
\bottomrule\noalign{}
\endlastfoot
(1) & \textbf{Tiếp cận tài chính} \emph{(Hạn chế tín dụng)} & Credit
constraint / Access to finance & Cao & Cao & Trung bình & Thấp & Aguinis
et al. (2011), formative composite measurement; Coltman et al. (2008)
construct validation cho credit access \\
(2) & \textbf{Lực lượng lao động thiếu kỹ năng} & Skilled labor shortage
/ Workforce skill gap & Cao & Trung bình & Trung bình & Thấp & Cohen \&
Levinthal (1990), absorptive capacity; Kafouros et al. (2023),
institutional quality moderation tech-dynamism \\
(3) & \textbf{Cạnh tranh từ doanh nghiệp phi chính thức} & Informal
sector competition & Cao & Cao & Trung bình & Thấp & Khanna \& Palepu
(2010), institutional voids; La Porta \& Shleifer (2008, 2014) shadow
economy share \\
(4) & \textbf{Nguồn cung điện không đáng tin cậy} & Unreliable
electricity supply / Power outages & Cao (đặc biệt SIDS, Frontier) &
Trung bình & Thấp & Thấp & IFC PSD Blueprint operational resources
pillar; Banerjee \& Duflo (2014) infrastructure productivity links \\
\end{longtable}

\textbf{Bốn nhận xét} (cập nhật theo WBES Asia-Pacific 2018--2025):

(a) \textbf{Rào cản 1 --- Hạn chế tín dụng}: Hơn 30\% doanh nghiệp ở
Frontier và Emerging xem \textbf{hạn chế tiếp cận tín dụng (credit
constraint)} là rào cản trên trung bình. Pattern phù hợp giả thuyết
\textbf{khoảng trống thể chế (institutional voids)} của Khanna \& Palepu
(2010): các thị trường đang phát triển có hệ thống tài chính kém phát
triển dẫn đến doanh nghiệp khó tiếp cận tín dụng chính thức. Liên kết
với Mục 4.4 \emph{Số hóa tài chính (Financial digitalization)}, công
nghệ tài chính (fintech) có thể giảm chi phí giao dịch nhưng chưa thay
thế được kênh ngân hàng truyền thống ở SME.

(b) \textbf{Rào cản 2, Lực lượng lao động thiếu kỹ năng}: Đặc biệt
nghiêm trọng cho doanh nghiệp sản xuất công nghệ cao (manufacturing
high-tech) và ICT (Bảng 4.8.1), đúng với khung \textbf{Schumpeter II về
động học công nghệ (Schumpeter Mark II tech dynamism)} (Kafouros et al.,
2023). Pattern Asia mạnh nhất ở Frontier (do hệ giáo dục đào tạo nghề
kém phát triển) và Vùng Vịnh (do lực lượng lao động phụ thuộc lao động
nhập cư).

(c) \textbf{Rào cản 3 --- Cạnh tranh từ doanh nghiệp phi chính thức}:
Tạo \textbf{cạnh tranh không công bằng (unfair competition)} cho các
doanh nghiệp chính thức vì doanh nghiệp phi chính thức không phải tuân
thủ thuế, lao động, chuẩn an toàn. Mức độ \textbf{kinh tế ngầm (shadow
economy)} ở Frontier và Emerging Á-Thái có thể lên 30--40\% GDP (La
Porta \& Shleifer, 2014). Giả thuyết: cạnh tranh phi chính thức làm giảm
động lực đầu tư R\&D và ISO chứng chỉ, liên kết Mục 4.4 \emph{Δ R\&D
Emerging giảm -42,1 đpt}.

(d) \textbf{Rào cản 4. Nguồn cung điện không đáng tin cậy}: Đặc biệt
nghiêm trọng cho SIDS Pacific và Frontier (e.g., Bangladesh, Pakistan
power outages \textgreater{} 10 hours/tháng theo WBES). Trực tiếp giảm
năng suất lao động (Banerjee \& Duflo, 2014), liên kết Mục 4.2
\emph{phân tán năng suất Frontier 1,36 cao nhất}. \textbf{IFC PSD
Blueprint} classify đây là \emph{operational resources pillar}, yếu tố
hạ tầng vật chất khác biệt với \emph{human resources pillar} (kỹ năng)
và \emph{financial resources pillar} (tín dụng).

\textbf{Hàm ý cho CĐ2}: 4 rào cản này dẫn đến 4 control variables mới
trong Spec 1 (full pool) và Spec 2 (2018-2025): (i)
\texttt{credit\_constraint\_dummy} từ WBES question k30; (ii)
\texttt{labor\_skill\_gap} từ question b8; (iii)
\texttt{informal\_competition} từ question e30; (iv)
\texttt{power\_outage\_intensity} từ question c30 (ngày/tháng mất điện).
Cộng với khung industry FE (Mục 4.8) và phân nhóm con thể chế (Chương 3
file 14), CĐ2 sẽ có \textbf{specification đầy đủ 5 lĩnh vực WBES}, không
chỉ 2 lĩnh vực (đổi mới + năng lực số) như Mục 4.4 hiện đang phân tích.

\hypertarget{456-giux1edbi-tuxednh-quux1ea3n-luxfd-cux1ea5p-cao-vuxe0-sux1edf-hux1eefu-bux1eb1ng-chux1ee9ng-tux1eeb-wbes-pool}{%
\paragraph{4.5.6 Giới tính quản lý cấp cao và sở hữu, bằng chứng từ WBES
pool}\label{456-giux1edbi-tuxednh-quux1ea3n-luxfd-cux1ea5p-cao-vuxe0-sux1edf-hux1eefu-bux1eb1ng-chux1ee9ng-tux1eeb-wbes-pool}}

\begin{quote}
\textbf{Nguồn chuẩn}: Carboni et al. (2024/2025 --- WBES 192.000 doanh
nghiệp, 158 quốc gia, 2006--2023); World Bank (2024) \emph{Unlocking
Global Growth: Closing the Gender Gap in Business} (WBES 167 nền kinh
tế). Biến WBES: \texttt{b7a} = female top manager (nhị phân Y/N);
\texttt{b4} = female majority ownership.
\end{quote}

\textbf{Bảng 4.5.6}. \emph{Tỷ lệ nữ quản lý cấp cao (female top manager)
và nữ sở hữu, theo phân nhóm con ICRV (nguồn: P7 pooled clean, n=87.872
doanh nghiệp; trường WBES b7a + b4).}

\begin{longtable}[]{@{}lllll@{}}
\toprule\noalign{}
Phân nhóm con & n & \% Nữ Top Manager (b7a) & \% Nữ Owner (b4) & Ghi
chú \\
\midrule\noalign{}
\endhead
\bottomrule\noalign{}
\endlastfoot
\textbf{Advanced, innovation-driven} & 4.695 & \textbf{25,6\%} &
\textbf{33,3\%} & Singapore/HK/Korea/Taiwan/Israel/Cyprus \\
\textbf{Advanced, resource-driven} & 1.932 & \textbf{4,0\%} &
\textbf{8,9\%} & Vùng Vịnh, thấp do cấu trúc xã hội-thể chế \\
\textbf{Upper-middle} & 10.349 & \textbf{31,2\%} & \textbf{33,9\%} &
Trung Quốc, Malaysia, Thái Lan dẫn dắt \\
\textbf{Emerging} & 48.297 & \textbf{13,2\%} & \textbf{24,7\%} & Việt
Nam, Indonesia, Philippines, Ấn Độ \\
\textbf{Frontier} & 21.157 & \textbf{13,6\%} & \textbf{21,5\%} & Phân
tán lớn xuyên Bangladesh, Pakistan, Cambodia, Trung Á \\
\textbf{SIDS Thái Bình Dương} & 1.442 & \textbf{31,8\%} &
\textbf{40,8\%} & Cao nhất pool, Fiji/Samoa/Solomon Islands có
matrilineal land tenure culture \\
\end{longtable}

\emph{Lưu ý: Số liệu thực tế đã được tính trực tiếp từ P7 pooled clean
(Đỗ \& Phan, 2026, JIBS revision); cao hơn ước tính ban đầu cho SIDS
Thái Bình Dương (matrilineal culture pattern xác nhận) và thấp hơn ước
tính cho Emerging Asia (chỉ Đông Nam Á cụ thể có pattern cao). Bảng
6.1.1 (file 16 Mục 6.1) cung cấp số liệu theo vùng WB rộng hơn.}

\textbf{Ba phát hiện nổi bật}:

(a) \textbf{EAP paradox, dẫn đầu về representation nhưng không về
performance gap}: Dữ liệu WBES EAP 33,4\% nữ TMT cao nhất toàn cầu,
nhưng Carboni et al. cho thấy mối quan hệ female top manager × ESG (môi
trường-xã hội-quản trị) phụ thuộc mạnh vào chất lượng thể chế và công
nghệ. Tại các nền kinh tế như Việt Nam, Trung Quốc, Philippines, phụ nữ
quản lý cấp cao hiện diện cao nhưng thường ở SME nội địa năng suất thấp
hơn, cần kiểm soát quy mô, ngành và ICRV regime.

(b) \textbf{Vùng Vịnh, gap lớn nhất}: Advanced resource-driven (Vùng
Vịnh Saudi/Qatar/Kuwait) có tỷ lệ nữ TMT thấp nhất trong toàn pool
(3--8\%), phản ánh cấu trúc thể chế-xã hội đặc thù. Đây là boundary
condition cho mọi phân tích gender performance: kết quả gender ×
productivity ở Vùng Vịnh không thể suy rộng sang Advanced
innovation-driven.

(c) \textbf{Frontier, phân tán nội bộ lớn}: Kyrgyz Republic + Tajikistan
(Frontier/Trung Á, nữ TMT \textasciitilde35\%) so với Yemen +
Afghanistan (Frontier/MENA, nữ TMT \textless5\%), minh họa cho lý do
ICRV phân nhóm Frontier 17 nước thay vì gộp chung.

\textbf{Biến WBES và hàm ý CĐ2}: Biến \texttt{b7a} (female top manager)
trong CĐ2 có thể đóng vai trò (i) \textbf{biến kiểm soát} trong
specification cơ bản để loại tác động cơ cấu giới lên năng suất; (ii)
\textbf{biến điều tiết (moderating variable)} cho H4 TMT moderation theo
Mardones-Ibáñez (2025, \emph{SAGE Open}) + Al-Najjar et al. (2025,
\emph{IJHRM}). Carboni et al. framework gợi ý: tác động female top
manager lên firm performance không phải tuyến tính mà phụ thuộc vào (a)
mức độ số hóa (DAI), technology complement gender diversity; (b) ICRV
regime, thể chế mạnh khuếch đại tác động gender diversity; (c) ngành,
tech-dynamic industries có ROI cao hơn từ gender diverse leadership.

\hypertarget{46-bux1ee9c-tranh-thay-ux111ux1ed5i-theo-thux1eddi-gian}{%
\subsubsection{4.6 Bức tranh thay đổi theo thời
gian}\label{46-bux1ee9c-tranh-thay-ux111ux1ed5i-theo-thux1eddi-gian}}

\textbf{Bảng 4.6}. \emph{Δ điểm phần trăm (đpt) khi so sánh 2018--2025
với 2009--2012.}

\begin{longtable}[]{@{}llllll@{}}
\toprule\noalign{}
Phân nhóm con & Δ Website & Δ Doanh nghiệp xuất khẩu & Δ FDI & Δ R\&D &
Δ ISO \\
\midrule\noalign{}
\endhead
\bottomrule\noalign{}
\endlastfoot
Upper-middle & -9,9 & +1,4 & +2,3 & +21,5 & -25,4 \\
Emerging & +20,3 & -7,5 & -10,9 & -42,1 & +1,9 \\
Frontier & +22,1 & +1,5 & -6,5 & -18,9 & +18,5 \\
SIDS Thái Bình Dương & +35--43 & +6--11 & -5 đến -10 & (cập nhật) & -20
đến -25 \\
\end{longtable}

Nhảy vọt số (digital leapfrog): website tăng +20--43 đpt ở Frontier,
Emerging và SIDS Thái Bình Dương, phù hợp luận điểm digital leapfrog của
Banalieva \& Dhanaraj (2019).

\includegraphics{figures/fig_4_4_growth_pathway.png}

\emph{Hình 4.4. Slope chart minh họa thay đổi 5 chỉ số (Website, DN xuất
khẩu, FDI, R\&D, ISO) giữa 2 mốc thời gian cho 4 phân nhóm con. Pattern
nổi bật: Website tăng mạnh +20-40 đpt ở Frontier/SIDS (digital
leapfrog); FDI giảm phổ biến (-5 đến -11 đpt); R\&D giảm mạnh ở Emerging
(-42,1 đpt) cảnh báo schema effect. Tái lập:
\texttt{generate\_figures.py} function
\texttt{fig\_4\_4\_growth\_pathway()}.}

\includegraphics{figures/fig_4_7_spider_chart.png}

\emph{Hình 4.7. Radar/spider chart cho 3 phân nhóm con (Advanced,
Emerging, SIDS) trên 5 chiều (FSTS, Innovation product, R\&D, Website,
ISO) với 2 mốc thời gian (2009--2012 vs 2018--2025). Quan sát: SIDS có
Innovation product cao đột biến + Website tăng mạnh giai đoạn 2018-2025
(digital leapfrog confirmed). Advanced ổn định high level. Emerging: ISO
+ R\&D thay đổi mạnh, FSTS giảm. Tái lập: \texttt{generate\_figures.py}
function \texttt{fig\_4\_7\_spider\_chart()}.}

\hypertarget{47-tux1ed5ng-hux1ee3p-chux1b0ux1a1ng-4-mux1edf-rux1ed9ng-d4-10-kux1ebft-luux1eadn-chuxednh}{%
\subsubsection{4.7 Tổng hợp Chương 4 (mở rộng D4, 10 kết luận
chính)}\label{47-tux1ed5ng-hux1ee3p-chux1b0ux1a1ng-4-mux1edf-rux1ed9ng-d4-10-kux1ebft-luux1eadn-chuxednh}}

Chương 4 cung cấp bức tranh thực trạng đa chiều dựa trên \textbf{101.185
doanh nghiệp ở 47 nền kinh tế châu Á và Thái Bình Dương (108 cặp QG×năm)
trong giai đoạn 2009--2025} từ nhóm dữ liệu WBES sau hài hòa. Đây là
phạm vi rộng nhất từng có cho nghiên cứu I→P trong văn liệu IB, mở rộng
từ 17 nước châu Á mới nổi (Đỗ \& Phan, 2026 --- VEFR) khoảng 2,5 lần.

\textbf{Mười kết luận chính}:

\textbf{(i) Phân tán năng suất nội bộ tăng đơn điệu khi phân nhóm con
suy giảm}: Advanced 0,86 dẫn đến Upper-middle 1,29 ≈ Emerging 1,24 ≈
SIDS 1,32 dẫn đến Frontier 1,36. Tỷ số P90/P10 leo từ 10,8 lần lên 39,6
lần. Pattern khẳng định \textbf{giả thuyết phân bổ sai nguồn lực} (Hsieh
\& Klenow, 2009, 2014).

\textbf{(ii) Dị biệt nội bộ phân nhóm Advanced}. Sự sụt giảm phân tán từ
1,00 (chỉ innovation-driven) xuống 0,86 (sau khi thêm Vùng Vịnh
resource-driven) gợi ý cần \textbf{phân nhóm con (sub-grouping)
Advanced} trong CĐ2. Tỷ số phân tán Singapore (1,03) so với Vùng Vịnh
(0,49) ≈ \textbf{2,1 lần} --- đã được tách trực quan trong \textbf{Bảng
4.1 (v3.1)}.

\textbf{(iii) Dị biệt nội bộ phân nhóm Emerging, 3 phân nhóm con}
\emph{(NEW D1, Mục 4.9)}. FDI dẫn dắt SEA (VNM+IDN+PHL, FSTS 13,2\%) ≠
dân số lớn (IND+LKA+JOR, FSTS 7,2\%) ≠ tài nguyên (MNG, FSTS 5,0\%). TCI
ngược dấu --- \textbf{học hỏi do FDI dẫn dắt vs học hỏi tự thân} (Cohen
\& Levinthal, 1990).

\textbf{(iv) SIDS Thái Bình Dương (7 nước, n=1.371)}: phân tán 1,32, đổi
mới sản phẩm cao nhất (41,5\%), website 58,9\% gần Advanced.
\textbf{Kiribati 2025 là trường hợp biên CỰC ĐOAN nhất với FSTS 1,03\%,
FDI 0,7\%, website 18,7\%, ISO 1,3\%, gợi ý dị biệt nội bộ SIDS giữa
"high digital leapfrog" (Fiji, Maldives) và "isolated rural"
(Kiribati).} Cơ sở cho \textbf{H6 (forced internationalization penalty)}
(Briguglio, 1995; Bertram, 2006).

\textbf{(v) Quốc tế hóa là hiện tượng phân cực}: trung vị FSTS = 0\%
xuyên năm phân nhóm con; chỉ 15--23\% doanh nghiệp xuất khẩu. Cần
\textbf{lựa chọn hai giai đoạn (2-stage selection)} trong CĐ2.

\textbf{(vi) Nhảy vọt số (digital leapfrog) 2018--2025} ở
Frontier/Emerging/SIDS (+20--43 đpt website), bằng chứng tái định vị
Uppsala (Banalieva \& Dhanaraj, 2019). \emph{Kiribati 2025 không thể
hiện leapfrog (website 18,7\%), cần điều kiện hạ tầng tối thiểu mà
Kiribati chưa đạt.} **: Cần phân biệt \textbf{Tier-1 Digital Presence}
(website binary, đo trong Spec 1, đã bão hòa ở Advanced từ trước 2018)
vs \textbf{Tier-2 Digital Transformation} (DAI multi-component,
e-commerce, cloud, AI adoption, đo trong Spec 2 với 5 thành phần). Hệ số
âm DAI ở Advanced trong Spec 1 (-0,129) KHÔNG phản ánh "số hóa làm giảm
năng suất" mà là \textbf{artifact của Tier-1 saturation} ở các nước tiên
tiến, xem Mục 4.4.5 ICT exclusion test.

\textbf{(vii) Pattern phi tuyến FDI ≥10\%}, dạng chữ U với cực tiểu
Emerging (4,7\%): Advanced 11,1\% dẫn đến Upper-middle 8,4\% dẫn đến
Emerging 4,7\% dẫn đến Frontier 5,9\% dẫn đến SIDS 23,5\%. Hai mô hình
FDI: (a) Advanced. MNE hub + Vùng Vịnh hạn chế; (b) SIDS, du lịch + viễn
thông.

\textbf{(viii) Đợt khảo sát 2025, mẫu đơn năm lớn nhất} \emph{(NEW D2)}.
\textbf{14 nước × 16.979 doanh nghiệp}: (a) IND FSTS sụt 5 đpt (schema +
Atmanirbhar Bharat); (b) Fiji website 74,8\% \textgreater{} Singapore
66,1\%; \textbf{Kiribati đối lập với website 18,7\%, FSTS 1,03\%}; (c)
Vùng Vịnh+Brunei resource-driven Advanced confirmed; (d) R\&D
schema-induced overestimation cảnh báo.

\textbf{(ix) Khung phân tích cấp ngành (industry-level)} \emph{(NEW D3,
Mục 4.8)}. 9 ngành ISIC Rev. 4: Manufacturing-only subsample, ICT
exclusion test, Construction subsample test Vùng Vịnh, Tourism/Hotels
separation cho SIDS. 5 giả thuyết I1-I5.

\textbf{(x) Pipeline tái lập được cho 4 thế hệ schema WBES}. Pipeline 5
bước Python; tất cả Bảng 4.1--4.10, 5.1, 6.1, Phụ lục A đều tái tạo
được. \textbf{Gói tái lập (replication package)} mở.

\textbf{Hàm ý cho CĐ2 và luận án}:

(a) \textbf{Hệ giả thuyết H1--H6}: H1 phi tuyến; H2 TCI điều tiết; H3
DAI có điều kiện; H4 phân nhóm con thể chế (5+8 phân nhóm con); H5 cụm
tài nguyên; H6 chi phí buộc phải quốc tế hóa (\textbf{7 SIDS gồm
Kiribati extreme}).

(b) \textbf{8 phân nhóm con với hiệu ứng cố định} cho CĐ2.

(c) \textbf{Hai đặc tả mô hình kiểm định vững}: Đặc tả 1, phạm vi đầy đủ
(n=101.185); Đặc tả 2, 2018-2025 (n≈50.000, DAI/TCI 5 thành phần).

(d) \textbf{Industry FE + 5 kiểm định vững mẫu con} \emph{(NEW D3)}.

\textbf{Phạm vi không trình bày}: hồi quy đa biến (CĐ2); panel/IV
identification (CĐ2); industry-level descriptive tables (Phase 2
7/2026); năng suất USD PPP (Phase 1 6/2026); resource dependence trực
tiếp (Phase 1).

\hypertarget{475-phuxe2n-biux1ec7t-3-thuux1eadt-ngux1eef-tham-nhux169ng-trong-wbes--bribery-vs-graft}{%
\paragraph{4.7.5 Phân biệt 3 thuật ngữ tham nhũng trong WBES --- Bribery
vs
Graft}\label{475-phuxe2n-biux1ec7t-3-thuux1eadt-ngux1eef-tham-nhux169ng-trong-wbes--bribery-vs-graft}}

\begin{quote}
\textbf{Nguồn chuẩn}: World Bank (n.d.). \emph{Tài liệu tóm tắt: Các chỉ
số khảo sát doanh nghiệp Enterprise Surveys} (file 04 v2.5 Section L).
Cross-cite glossary v1.2 Nhóm 6 \#53-54 (commit e9a73ae) + Nhóm 7
\#64-66 (commit beb47d3).
\end{quote}

WBES đo lường tham nhũng qua \textbf{3 chỉ số riêng biệt}, phân biệt rõ
giữa \textbf{hành vi chủ thể} (subject behavior, hối lộ vs tham ô) và
\textbf{mức độ phổ biến} (incidence vs depth). Sự phân biệt này quan
trọng cho phân tích chính sách: hối lộ và tham ô có nguyên nhân + giải
pháp khác nhau, dù đôi khi gộp chung dưới label "tham nhũng" trong văn
liệu chính sách Việt Nam.

\textbf{Bảng 4.7.5}. \emph{Phân biệt 3 chỉ số tham nhũng WBES --- định
nghĩa chính thức + cách đo + ngưỡng cảnh báo.}

\begin{longtable}[]{@{}llllll@{}}
\toprule\noalign{}
\# & Chỉ số (Tiếng Anh) & Tiếng Việt CHUẨN WB & Định nghĩa & Cách đo
(công thức) & Ngưỡng cảnh báo \\
\midrule\noalign{}
\endhead
\bottomrule\noalign{}
\endlastfoot
(a) & \textbf{Bribery} \emph{(khái niệm gốc)} & \textbf{Hối lộ} & Hành
vi đưa/nhận tiền/quà cho lợi thế trong giao dịch công (cấp phép, thanh
tra, hoàn thuế, kết nối điện thoại/điện) & (đo qua 3 chỉ số dưới đây) &
--- \\
(b) & \textbf{Graft} \emph{(khái niệm gốc)} & \textbf{Tham ô}
\emph{(rộng hơn hối lộ)} & Biển thủ công quỹ + lạm dụng chức vụ trục lợi
cá nhân & (đo qua chỉ số (b-1) dưới đây) & --- \\
(a-1) & \textbf{Bribery Incidence} & \textbf{Tỷ lệ hối lộ} & \% DN gặp
ít nhất 1 yêu cầu hối lộ trong 6 giao dịch công khảo sát & Tử số: \# DN
có ≥1 yêu cầu hối lộ; Mẫu: \# DN trong khảo sát có giao dịch trong 6
dạng & \textgreater20\% = báo động đỏ \\
(a-2) & \textbf{Bribery Depth} & \textbf{Độ sâu hối lộ} & \% giao dịch
công có hối lộ trên tổng số giao dịch công của DN & Tử số: \# giao dịch
có hối lộ; Mẫu: tổng \# giao dịch công & \textgreater15\% = báo động
đỏ \\
(b-1) & \textbf{Graft Index} & \textbf{Chỉ số tham ô} & Tỷ lệ hối lộ
trong 6 giao dịch (gồm kết nối điện thoại/điện) NHƯNG \textbf{loại trừ
thanh tra thuế} & Tử số: như Bribery Incidence; Mẫu: 5 giao dịch (loại
bỏ tax inspection vì thiên lệch khai thuế) & \textgreater10\% = báo động
đỏ \\
\end{longtable}

\textbf{Tại sao loại trừ thanh tra thuế khỏi Graft Index}: Doanh nghiệp
khai báo "phải hối lộ thanh tra thuế" có thể bị thiên lệch (selection
bias) vì những DN trốn thuế dễ thừa nhận (do đã có hành vi vi phạm),
trong khi DN tuân thủ thuế ít gặp tình huống thanh tra cụ thể. Loại trừ
tax inspection giúp Graft Index phản ánh \textbf{tham ô hệ thống}
(systemic graft) thay vì hành vi cá biệt liên quan đến trốn thuế.

\textbf{Hàm ý cho CĐ2 và luận án}:

\begin{itemize}
\tightlist
\item
  \textbf{(1) CĐ1 v3.x} sẽ dùng "Hối lộ" (Bribery) và "Tham ô" (Graft)
  thay vì "tham nhũng" tổng quát khi cite số liệu cụ thể từ WBES. Khi
  nói chung về vấn đề thể chế hoặc chính sách quốc gia, vẫn có thể dùng
  "tham nhũng" như một tổng thể.
\item
  \textbf{(2) CĐ2 specification} sẽ có \textbf{3 biến phân biệt}:
  \texttt{bribery\_incidence\_pct} (a-1), \texttt{bribery\_depth\_pct}
  (a-2), \texttt{graft\_index\_pct} (b-1) từ WBES question j7a-j7f sang
  kiểm tra hệ số khác nhau xuyên 8 phân nhóm con thể chế. Giả thuyết H8
  (mới): tác động tham nhũng lên hiệu quả doanh nghiệp khác biệt giữa
  "extensive corruption" (Bribery Incidence cao, Depth thấp, tham nhũng
  phổ biến nhưng mỗi giao dịch ít) và "intensive corruption" (Incidence
  thấp, Depth cao, chỉ một số ít DN nhưng giao dịch nào cũng tham
  nhũng).
\item
  \textbf{(3) Liên kết với Mục 4.5.5 rào cản \#3} (cạnh tranh phi chính
  thức): DN phi chính thức tránh thuế nhưng cũng có thể \textbf{né tránh
  hối lộ (bypass bribery)} do hoạt động ngoài hệ thống chính thức, tạo
  ra \textbf{tình thế tiến thoái lưỡng nan chính sách (policy dilemma)}:
  \textbf{chính thức hóa (formalization)} có thể tăng \textbf{mức độ
  tuân thủ (compliance)} nhưng cũng tăng \textbf{mức phơi nhiễm hối lộ
  (exposure to bribery)}.
\item
  \textbf{(4) Liên kết với khung phân tích cấp ngành Mục 4.8}: ngành Xây
  dựng (cường độ hợp đồng cao, high contract intensity) thường có
  Bribery Depth cao do từng giao dịch công lớn; ngành Khai khoáng/Du
  lịch (tỷ trọng công ty đa quốc gia, high MNE share) có Bribery
  Incidence cao nhưng Depth thấp do \textbf{áp lực tuân thủ MNE (MNE
  compliance pressure)}.
\end{itemize}

\hypertarget{48-khung-phuxe2n-tuxedch-cux1ea5p-nguxe0nh-industry-level-kux1ebf-houx1ea1ch-cux1112}{%
\subsubsection{4.8 Khung phân tích cấp ngành (industry-level), kế hoạch
CĐ2}\label{48-khung-phuxe2n-tuxedch-cux1ea5p-nguxe0nh-industry-level-kux1ebf-houx1ea1ch-cux1112}}

\begin{quote}
\textbf{Mới ở v2.9 (D3)}: Khung dữ liệu (schema) WBES phân loại theo
\texttt{a3a} (mã ngành) và \texttt{a4a} (mã ISIC 4 chữ số).

\textbf{Mở rộng v3.4 (07/05/2026)}: Bảng 4.8.1 thêm cột "Dynamism
profile (Kafouros 2023)"; thêm đoạn lý giải sau bảng về industry ×
institutional quality interaction; mở rộng 5 dẫn đến 6 hàm ý phương pháp
luận với hàm ý (f) test interaction term.
\end{quote}

\textbf{Bảng 4.8.1}. \emph{Khung phân loại 9 ngành ISIC Rev. 4 cho nhóm
dữ liệu WBES.}

\begin{longtable}[]{@{}lllll@{}}
\toprule\noalign{}
Ngành & Mã ISIC & Đặc điểm & Pattern dự kiến & \textbf{Dynamism profile
(Kafouros 2023)} \\
\midrule\noalign{}
\endhead
\bottomrule\noalign{}
\endlastfoot
Manufacturing, high-tech & C (20--21, 26--30) \emph{(pharma, computers,
electrical, motor vehicles)} & Thâm dụng vốn, định hướng xuất khẩu &
FSTS cao; R\&D cao; ISO cao & \textbf{High technological dynamism}
\emph{(Schumpeter Mark II, deepening pattern)} \\
Manufacturing, low-tech & C (10--18, 22--25, 31--33) \emph{(food,
beverages, textiles, leather, furniture, basic metals)} & Thâm dụng lao
động & FSTS trung bình; R\&D thấp & \textbf{Low technological dynamism}
\emph{(widening pattern; longer technology life cycle)} \\
Wholesale/retail & G (45--47) & Thâm dụng lao động, định hướng trong
nước & FSTS thấp; website cao; FDI thấp & \textbf{High market dynamism}
\emph{(volatile consumer demand, short product cycles)} \\
Tourism/hotels & I (55--56) & Dịch vụ phụ thuộc du lịch & FDI cao ở SIDS
& \textbf{High market dynamism} \emph{(volatile demand, pandemic,
geopolitical, seasonal)} \\
Transport & H (49--53) & Kinh tế mạng lưới & FDI cao; ISO cao & Trung
gian (low-medium dynamism) \\
ICT & J (58--63) & Thâm dụng tri thức & R\&D cao; website
\textasciitilde100\% & \textbf{High technological dynamism} \emph{(rapid
tech evolution, IPR-sensitive)} \\
Construction & F (41--43) & Theo dự án, định hướng trong nước & FSTS
thấp; FDI ở GCC & \textbf{Low technological dynamism + low market
dynamism} \emph{(stable demand, traditional methods)} \\
Mining & B (05--09) & Tài nguyên dẫn dắt & FDI cao; ISO ở MNE-led &
\textbf{Low technological dynamism} \emph{(commodity-driven,
institutional rents)} \\
Finance & K (64--66) & Có quy chế, xuyên biên giới & FDI cao; website
\textasciitilde100\% & \textbf{High market dynamism} \emph{(demand
shocks, AI disruption Wang/Huang/Hong 2024)} \\
Khác & A, D, L & Đa dạng & Hỗn hợp & Hỗn hợp \\
\end{longtable}

\textbf{5 giả thuyết cấp ngành cho CĐ2}: I1 Manufacturing FSTS
dominance; I2 ICT digital-native; I3 Tourism drive FDI ở SIDS; I4 Mining
drive resource cluster; I5 Construction dominate Vùng Vịnh.

\textbf{Lý giải lý thuyết. Kafouros et al. (2023) industry ×
institutional quality interaction}:

Khung 9 ngành trên không chỉ phân loại theo logic kinh tế mà còn theo
logic lý thuyết của \textbf{Kafouros et al. (2023 --- \emph{Global
Strategy Journal}, Vol 14, 56--83)}. Kafouros và cộng sự (n=12.888 firms
· 16 CEE economies · 2004-2011 · 72.082 obs) chứng minh chất lượng thể
chế (đo bằng Rule of Law từ World Governance Indicators) tương tác hai
chiều với động học ngành: \textbf{(H1)} ở các ngành có
\textbf{technological dynamism} cao (Schumpeter Mark II, pharma,
computers, electrical, ICT, motor vehicles), chất lượng thể chế
\textbf{khuếch đại tích cực} hiệu quả doanh nghiệp do ba cơ chế
(partnership identification, IPR enforcement, interfirm market
exchange); \textbf{(H2)} ngược lại, ở các ngành có \textbf{market
dynamism} cao (volatility cầu, tourism, retail, finance), chất lượng thể
chế \textbf{suy yếu} vì doanh nghiệp phải nội bộ hóa (internalize) các
chức năng để giảm phụ thuộc thị trường ngoài. \textbf{Hàm ý cho CĐ1 và
CĐ2}: (a) hệ số âm DAI ở Advanced phân nhóm innovation-driven (-0,129
trong Spec 1, xem Mục 4.4 và Phase 2 HĐ1) có thể bị nhiễu bởi ICT firms
saturated; (b) Bảng 4.6 cho thấy Δ R\&D Emerging giảm -42,1 đpt, nếu
phân tách high-tech vs low-tech Manufacturing, có thể thấy pattern khác
biệt; (c) \textbf{CĐ2 cần thêm giả thuyết H7}: tương tác
\texttt{tech\_dynamism\ ×\ institutional\_quality} có hệ số dương; tương
tác \texttt{market\_dynamism\ ×\ institutional\_quality} có hệ số âm.
Đây là test trực tiếp cho replication của Kafouros (2023) trong bối cảnh
châu Á --- đóng góp lý thuyết quan trọng vì Kafouros chỉ test trong bối
cảnh CEE.

\textbf{6 hàm ý phương pháp luận} \emph{)}:

(a) Industry FE 9 ngành theo Bảng 4.8.1. (b) Manufacturing-only
subsample, tách cho high-tech vs low-tech theo Kafouros (2023)
classification. (c) Tourism separation cho SIDS Pacific (boundary case
extension Mục 1.3). (d) Construction subsample test Vùng Vịnh
resource-driven Advanced. (e) ICT exclusion test cho DAI --- đặc biệt
liên quan Phase 2 HĐ1 Tier-1 Digital Presence rebrand. \textbf{(f) Test
interaction term \texttt{tech\_dynamism\ ×\ institutional\_quality} và
\texttt{market\_dynamism\ ×\ institutional\_quality}}, robustness check
thứ 7 cho CĐ2 H7 mới (sau Manufacturing-only, ICT-excluded,
Tourism-separated SIDS, Construction-tested Gulf, Mining-excluded
resource, Anchor model BREADY 2025). Đây là test trực tiếp cho
\textbf{replication Kafouros et al. (2023) trong bối cảnh châu Á}, bảng
Industry × Institutional Quality interaction cần thiết để
confirm/disconfirm H1-H2 của Kafouros với 41 nước châu Á (vs 16 nước CEE
của Kafouros). Nếu interaction tech\_dynamism × institutional\_quality
dương và significant ở mức p\textless.05, đây là evidence châu
Á-specific cho industry-bounded institutional advantage --- đóng góp lý
thuyết của luận án.

\hypertarget{49-phuxe2n-nhuxf3m-con-emerging-phuxe1t-hiux1ec7n-dux1ecb-biux1ec7t-nux1ed9i-bux1ed9}{%
\subsubsection{4.9 Phân nhóm con Emerging, phát hiện dị biệt nội
bộ}\label{49-phuxe2n-nhuxf3m-con-emerging-phuxe1t-hiux1ec7n-dux1ecb-biux1ec7t-nux1ed9i-bux1ed9}}

\begin{quote}
\textbf{Mới ở v2.7 (D1)}: Phát hiện phương pháp luận thứ hai cho CĐ2.
\end{quote}

\textbf{Bảng 4.9}. \emph{Phân nhóm con Emerging, số liệu tổng hợp
2009--2025 (n=47.803).}

\begin{longtable}[]{@{}llllllllll@{}}
\toprule\noalign{}
Phân nhóm con & Quốc gia & n\_firms & FSTS (\%) & DN xuất khẩu (\%) &
FDI ≥10\% (\%) & Website (\%) & R\&D (\%) & ISO (\%) & sd log \\
\midrule\noalign{}
\endhead
\bottomrule\noalign{}
\endlastfoot
Emerging. FDI dẫn dắt SEA & VNM, IDN, PHL & 13.779 & \textbf{13,2} &
22,1 & 11,4 & 47,3 & 4,8 & 18,8 & 1,53 \\
Emerging, dân số lớn & IND, LKA, JOR & 32.119 & \textbf{7,2} & 13,8 &
1,9 & 49,5 & 19,8 & 27,6 & 1,16 \\
Emerging, tài nguyên & MNG & 1.905 & \textbf{5,0} & 9,7 & 4,7 & 50,1 &
20,8 & 15,4 & 1,16 \\
\textbf{Tổng Emerging} & 7 nước & \textbf{47.803} & 8,6 & 15,5 & 4,7 &
49,2 & 16,4 & 24,9 & 1,24 \\
\end{longtable}

Năm phát hiện: (1) FSTS phân tầng 5,0\%-7,2\%-13,2\% bị che giấu; (2)
FDI chênh lệch 6× ASEAN-3 vs Nam Á+Tây Á; (3) TCI ngược dấu; (4)
Mongolia boundary case; (5) sd log 1,53 vs 1,16, two-tier (Hsieh \&
Klenow, 2009).

\hypertarget{410-phuxe2n-tuxedch-suxe2u-ux111ux1ee3t-khux1ea3o-suxe1t-2025}{%
\subsubsection{4.10 Phân tích sâu đợt khảo sát
2025}\label{410-phuxe2n-tuxedch-suxe2u-ux111ux1ee3t-khux1ea3o-suxe1t-2025}}

\begin{quote}
\textbf{Mới ở v2.8 (D2)}: Đợt 2025 chiếm 16.979 doanh nghiệp (16,8\%
pool, v3.1 với Kiribati 2025).
\end{quote}

\textbf{Bảng 4.10}. \emph{14 quốc gia trong đợt 2025 (n=16.979, v3.1 với
Kiribati 2025).}

\begin{longtable}[]{@{}lllllllll@{}}
\toprule\noalign{}
Quốc gia & ISO3 & Phân nhóm con & n\_firms & FSTS (\%) & FDI (\%) & R\&D
(\%) & Website (\%) & sd log \\
\midrule\noalign{}
\endhead
\bottomrule\noalign{}
\endlastfoot
Ấn Độ & IND & Emerging & 10.479 & 2,7 & 1,9 & 2,2 & 41,8 & 0,83 \\
Nepal¹ & NPL & Frontier & 1.740 & n/a & n/a & n/a & n/a & n/a \\
Saudi Arabia & SAU & Advanced & 1.002 & 2,7 & 9,5 & 1,7 & 30,2 &
\textbf{0,47} \\
Thái Lan & THA & Upper-middle & 813 & 9,3 & 6,3 & 8,7 & 61,9 & 1,38 \\
Sri Lanka & LKA & Emerging & 607 & 16,1 & 1,8 & 4,1 & 48,8 & 1,00 \\
Mongolia & MNG & Emerging & 601 & 5,9 & 3,2 & 20,8 & 64,7 & 1,15 \\
Qatar & QAT & Advanced & 480 & 2,3 & 19,4 & 0,6 & 63,3 &
\textbf{0,31} \\
Afghanistan & AFG & Frontier & 480 & 6,4 & 1,0 & 25,8 & 45,8 & 1,42 \\
Maldives & MDV & Frontier & 154 & 5,8 & 4,5 & 22,4 & \textbf{73,4} &
1,36 \\
Fiji & FJI & SIDS & 151 & 12,5 & 9,9 & 18,8 & \textbf{74,8} & 1,09 \\
Solomon Is. & SLB & SIDS & 150 & 4,8 & 20,0 & 11,3 & 53,3 & 1,25 \\
Brunei & BRN & Advanced & 150 & 5,1 & 26,2 & 18,9 & \textbf{80,7} &
1,10 \\
Kuwait & KWT & Advanced & 150 & \textbf{0,4} & 0,0 & 20,7 & 69,3 &
1,15 \\
\textbf{Kiribati²} & \textbf{KIR} & \textbf{SIDS} & \textbf{150} &
\textbf{1,0} & \textbf{0,7} & \textbf{14,0} & \textbf{18,7} &
\textbf{1,48} \\
\textbf{Tổng 2025} & , & , & \textbf{16.979} & \textbf{3,8} &
\textbf{2,9} & \textbf{5,4} & \textbf{44,2} & \textbf{0,90} \\
\end{longtable}

\emph{Ghi chú: ¹Nepal 2025, schema BREADY chưa thống nhất. ²Kiribati
2025: Lower Middle Income; viện trợ AUS/NZ \textasciitilde30\% GDP.}

\textbf{Bảy phát hiện}: (1) IND FSTS sụt 5 đpt (schema effect); (2) THA
"digital up, exports down"; (3) Fiji website 74,8\% \textgreater{}
Singapore (digital leapfrog); (4) Vùng Vịnh + Brunei resource-driven
Advanced confirmed; (5) R\&D schema-induced overestimation; (6) 2025
wave validation sample CĐ2; \textbf{(7) Kiribati extreme, FSTS 1,03\%,
FDI 0,7\%, website 18,7\%, đối lập với Fiji digital leapfrog. Bằng chứng
dị biệt SIDS rõ rệt: "high-digital" (FJI, MDV) vs "isolated rural"
(KIR). CĐ2 cần tách 2 phân nhóm con SIDS.}

\textbf{Bảy hàm ý cho CĐ2}: (a) 2025 validation test bed; \textbf{(b)
Schema FE PostBREADY2024 enhanced với anchor model} \emph{(chi tiết Mục
4.11 v3.7)}; (c) Advanced sub-grouping test 11 quốc gia; (d) SIDS
digital leapfrog evidence cho H6, phân biệt \textbf{Tier-1} (website
saturation, đã bão hòa) vs \textbf{Tier-2} (transformation, đang phát
triển); (e) Two-wave panel cho 6 nước; (f) tách phân nhóm con SIDS,
"high-digital" vs "isolated", mở rộng 8 dẫn đến 9 phân nhóm con; (g)
Tách 2025 thành \textbf{"Panel hậu đại dịch độc lập"} \emph{(chi tiết
Mục 4.11 v3.7)}.

\hypertarget{411-cuxf4-lux1eadp-ux111ux1ee9t-guxe3y-schema-bready-2025-3-ux111ux1ec1-xuux1ea5t-phux1b0ux1a1ng-phuxe1p-luux1eadn}{%
\subsubsection{4.11 Cô lập đứt gãy schema BREADY 2025, 3 đề xuất phương
pháp
luận}\label{411-cuxf4-lux1eadp-ux111ux1ee9t-guxe3y-schema-bready-2025-3-ux111ux1ec1-xuux1ea5t-phux1b0ux1a1ng-phuxe1p-luux1eadn}}

Đợt 2025 (n=16.979 firms, \textbf{16,8\% pool}) sử dụng schema mới
BREADY khác bản hỏi cũ, gây cú sốc phi lý: IND FSTS rơi 7,7\%→2,7\% (xem
Mục 4.10), R\&D nhiều nước tăng vọt do \textbf{hiệu ứng bảng hỏi
(questionnaire effect)} chứ không phải thực tế. Theo : \emph{"Đứt gãy
schema BREADY 2025 đe dọa trực tiếp đến tính liên tục của phân tích
chuỗi thời gian. Mô hình kiểm tra điểm uốn sẽ cho ra kết luận sai lệch
hoàn toàn nếu không xử lý."} NCS đã có cảnh báo trong Mục 4.10 (mục 5:
R\&D schema-induced overestimation) nhưng cần làm sâu hơn với 3 đề xuất
phương pháp luận.

\hypertarget{4111-ux111ux1ec1-xuux1ea5t-3a-biux1ebfn-giux1ea3-post_bready_2024-trong-mux1ecdi-spec-tux1ed5ng-gux1ed9p}{%
\paragraph{\texorpdfstring{4.11.1 Đề xuất 3a. Biến giả
\texttt{Post\_BREADY\_2024} trong mọi spec tổng
gộp}{4.11.1 Đề xuất 3a. Biến giả Post\_BREADY\_2024 trong mọi spec tổng gộp}}\label{4111-ux111ux1ec1-xuux1ea5t-3a-biux1ebfn-giux1ea3-post_bready_2024-trong-mux1ecdi-spec-tux1ed5ng-gux1ed9p}}

Khai báo biến giả \texttt{Post\_BREADY\_2024} = 1 cho mọi observation
thuộc đợt khảo sát 2025+ (BREADY schema), = 0 cho 2009--2024 (PICS3 +
Standardized + early BREADY 2018-2024). Biến giả này:

\begin{itemize}
\tightlist
\item
  \textbf{Hấp thụ schema effect tĩnh}: lấy out mức trung bình của các
  thay đổi đo lường giữa 2 generation bảng hỏi (e.g., IND FSTS
  7,7\%→2,7\% có thể một phần do BREADY questions hỏi rõ hơn về "doanh
  thu xuất khẩu trực tiếp" vs PICS3 hỏi tổng quát hơn).
\item
  \textbf{Cảnh báo phương pháp luận}: \emph{"Biến giả đâu phải cây đũa
  thần. Nó chỉ như miếng bọt biển hút tác động tĩnh, không sửa được
  phương sai."} Vì vậy cần đi cùng đề xuất 3b (anchor model) để kiểm tra
  robustness.
\item
  \textbf{Triển khai}: thêm \texttt{Post\_BREADY\_2024} vào Spec 1 (full
  pool, 101.185 firms) và Spec 2 (2018-2025, \textasciitilde50.000
  firms) cho CĐ2.
\end{itemize}

\hypertarget{4112-ux111ux1ec1-xuux1ea5t-3b-muxf4-huxecnh-neo-anchor-model}{%
\paragraph{4.11.2 Đề xuất 3b. Mô hình neo (anchor
model)}\label{4112-ux111ux1ec1-xuux1ea5t-3b-muxf4-huxecnh-neo-anchor-model}}

Quy trình 3 bước:

(i) Chạy hồi quy với data chỉ đến 2024 (n≈84.000 firms, 13 đợt khảo sát
2009-2024) dẫn đến có \textbf{bộ hệ số "neo"} (anchor coefficients) cho
FSTS, FSTS², TCI, DAI và các tương tác.

(ii) \textbf{Khóa hệ số} từ bước (i). Trong R:
\texttt{glmer(...\ offset\ =\ X1\ *\ beta1\_anchor\ +\ X2\ *\ beta2\_anchor,\ data=full\_pool\_with\_2025)}.

(iii) Re-run hồi quy với data đầy đủ 2009-2025 (n=101.185), so sánh:

\begin{itemize}
\tightlist
\item
  Coefficient stability test (Wald test, Chow test): hệ số FSTS, FSTS²,
  TCI, DAI có thay đổi significantly khi thêm 2025 không?
\item
  Likelihood ratio test giữa 2 mô hình: Δ AIC, Δ BIC.
\item
  Coefficient ratio test: nếu
  \texttt{beta\_2025\ /\ beta\_anchor\ \textgreater{}\ 1.5x} hoặc
  \texttt{\textless{}\ 0.67x} dẫn đến schema effect significant.
\end{itemize}

\textbf{Mục đích}: kiểm tra "cấu trúc bảng hỏi mới có làm đảo chiều ý
nghĩa thống kê cốt lõi hay không". Đây là \textbf{robustness check thứ 6
cho CĐ2} (sau Manufacturing-only, ICT-excluded, Tourism-separated SIDS,
Construction-tested Gulf, Mining-excluded resource).

\hypertarget{4113-ux111ux1ec1-xuux1ea5t-3c-tuxe1ch-2025-thuxe0nh-panel-hux1eadu-ux111ux1ea1i-dux1ecbch-ux111ux1ed9c-lux1eadp}{%
\paragraph{4.11.3 Đề xuất 3c. Tách 2025 thành "Panel hậu đại dịch độc
lập"}\label{4113-ux111ux1ec1-xuux1ea5t-3c-tuxe1ch-2025-thuxe0nh-panel-hux1eadu-ux111ux1ea1i-dux1ecbch-ux111ux1ed9c-lux1eadp}}

Thay vì cố ép 2025 vào đường xu hướng lịch sử, dành riêng một
sub-section trong CĐ2 (đề xuất Mục X.Y trong Mục 6 mô hình) để xử lý
2025 như \textbf{tập mẫu xác thực cho kỷ nguyên hậu COVID + AI}. Câu hỏi
nghiên cứu mới: \emph{"Quy luật thể chế và U-curve giai đoạn 2009-2024
có còn hiệu lực trong \textquotesingle trạng thái bình thường
mới\textquotesingle{} (new normal) hậu COVID + AI bùng nổ?"}

Tiến trình:

\begin{itemize}
\tightlist
\item
  \textbf{Bước 1}: Chạy mô hình chính trên data 2009-2024, báo cáo H1-H6
  + H7 + Lin-Mehlum tests.
\item
  \textbf{Bước 2}: Test cùng mô hình trên data 2025 only (n=16.979) như
  \textbf{out-of-sample validation panel}. Ghi rõ "post-pandemic + AI
  era panel" (post-COVID 2020-2022 + AI surge 2023-2025).
\item
  \textbf{Bước 3}: Báo cáo:

  \begin{itemize}
  \tightlist
  \item
    Hệ số nào giữ stability (validation passed)
  \item
    Hệ số nào thay đổi (validation failed dẫn đến discuss as evidence
    của paradigm shift hậu COVID + AI)
  \item
    Effect size differences và policy implications.
  \end{itemize}
\end{itemize}

\textbf{Nhận xét nội bộ}: \emph{"Biến rủi ro dữ liệu thành chương phân
tích hướng tới tương lai cực kỳ giá trị."}

\hypertarget{4114-liuxean-kux1ebft-vux1edbi-mux1ee5c-6-yux1ebfu-tux1ed1-giux1ea3i-thuxedch-vuxe0-mux1ee5c-73-huxe0m-uxfd-cux1112}{%
\paragraph{4.11.4 Liên kết với Mục 6 (yếu tố giải thích) và Mục 7.3 (hàm
ý
CĐ2)}\label{4114-liuxean-kux1ebft-vux1edbi-mux1ee5c-6-yux1ebfu-tux1ed1-giux1ea3i-thuxedch-vuxe0-mux1ee5c-73-huxe0m-uxfd-cux1112}}

3 đề xuất 3a-3b-3c kết hợp tạo thành \textbf{hệ thống cô lập
triple-defense} chống đứt gãy schema 2025:

\begin{longtable}[]{@{}llll@{}}
\toprule\noalign{}
Đề xuất & Tác dụng & Robustness check \# & Liên kết \\
\midrule\noalign{}
\endhead
\bottomrule\noalign{}
\endlastfoot
3a Biến giả Post\_BREADY\_2024 & Hấp thụ schema effect tĩnh &
(parametric correction) & Spec 1 + Spec 2 \\
3b Anchor model & Test stability hệ số & \#6 (sau
Mfg/ICT/Tourism/Gulf/Mining) & CĐ2 robustness section \\
3c Panel hậu đại dịch độc lập & Out-of-sample validation hậu COVID+AI &
\#7 panel test & CĐ2 Mục X.Y dedicated section \\
\end{longtable}

Cả 3 đều là tiền đề cho file 16 Mục 7.3.4 hàm ý cho CĐ2, robustness
check section.

\begin{center}\rule{0.5\linewidth}{0.5pt}\end{center}

\hypertarget{chuyuxean-ux111ux1ec1-tiux1ebfn-sux129-sux1ed1-1--bux1ea3n-nhuxe1p-ux111ux1ea7y-ux111ux1ee7-phux1ea7n-3-chux1b0ux1a1ng-57--tuxe0i-liux1ec7u-tham-khux1ea3o--phux1ee5-lux1ee5c}{%
\section{CHUYÊN ĐỀ TIẾN SĨ SỐ 1 --- BẢN NHÁP ĐẦY ĐỦ (PHẦN 3: CHƯƠNG 5--7
+ TÀI LIỆU THAM KHẢO + PHỤ
LỤC)}\label{chuyuxean-ux111ux1ec1-tiux1ebfn-sux129-sux1ed1-1--bux1ea3n-nhuxe1p-ux111ux1ea7y-ux111ux1ee7-phux1ea7n-3-chux1b0ux1a1ng-57--tuxe0i-liux1ec7u-tham-khux1ea3o--phux1ee5-lux1ee5c}}

\begin{quote}
Tiếp nối \texttt{thesis/14\_cd1\_part1\_intro\_theory\_vi.md} và
\texttt{thesis/15\_cd1\_part2\_findings\_vi.md}. Bảng thuật ngữ
Anh-Việt: \texttt{thesis/09b\_vn\_term\_glossary.md}. Tổng hợp Asia
context: \texttt{thesis/\_asia\_context\_synthesis\_2026.md}. Hình minh
họa: \texttt{thesis/figures/} (11 hình; chạy
\texttt{python3\ generate\_figures.py}).
\end{quote}

\begin{center}\rule{0.5\linewidth}{0.5pt}\end{center}

\hypertarget{chux1b0ux1a1ng-5--bux1ea3y-tiux1ec3u-cux1ea3nh-ux111iux1ec3n-huxecnh}{%
\subsection{CHƯƠNG 5 --- BẢY TIỂU CẢNH ĐIỂN
HÌNH}\label{chux1b0ux1a1ng-5--bux1ea3y-tiux1ec3u-cux1ea3nh-ux111iux1ec3n-huxecnh}}

\hypertarget{51-singapore-advanced-ux111ux1ed5i-mux1edbi-suxe1ng-tux1ea1o-dux1eabn-dux1eaft-innovation-driven-n623-ux111ux1ee3t-2023}{%
\subsubsection{5.1 Singapore (Advanced đổi mới sáng tạo dẫn dắt,
innovation-driven, n=623, đợt
2023)}\label{51-singapore-advanced-ux111ux1ed5i-mux1edbi-suxe1ng-tux1ea1o-dux1eabn-dux1eaft-innovation-driven-n623-ux111ux1ee3t-2023}}

FSTS 7,1\%; doanh nghiệp xuất khẩu 17,8\%; website 66,1\%; ISO 23,3\%;
R\&D 7,5\%; FDI 31,5\%. Độ lệch chuẩn log năng suất (sd log) 1,03.
\textbf{Biên trên (upper boundary)} của nhóm Advanced đổi mới sáng tạo
dẫn dắt. Theo bài báo P3 (Singapore. MIR): mô hình M8 cho R² hiệu chỉnh
= 0,196; tương tác FSTS² × DAI = 3,119; điểm uốn (turning point) ở mức
\textasciitilde85\%.

\hypertarget{52-saudi-arabia-qatar-vuxe0-kuwait-advanced-tuxe0i-nguyuxean-dux1eabn-dux1eaft-resource-driven-n1632-3-ux111ux1ee3t-nux103m-2025}{%
\subsubsection{5.2 Saudi Arabia, Qatar và Kuwait (Advanced tài nguyên
dẫn dắt, resource-driven, n=1.632, 3 đợt năm
2025)}\label{52-saudi-arabia-qatar-vuxe0-kuwait-advanced-tuxe0i-nguyuxean-dux1eabn-dux1eaft-resource-driven-n1632-3-ux111ux1ee3t-nux103m-2025}}

\begin{longtable}[]{@{}lllll@{}}
\toprule\noalign{}
Chỉ số & Saudi Arabia & Qatar & Kuwait & Singapore \\
\midrule\noalign{}
\endhead
\bottomrule\noalign{}
\endlastfoot
FSTS (\%) & 2,7 & 2,3 & \textbf{0,4} & 7,1 \\
FDI ≥10\% (\%) & 9,5 & 19,4 & 0,0 & \textbf{31,5} \\
R\&D dương (\%) & 1,7 & 0,6 & \textbf{20,7} & 7,5 \\
sd log năng suất & \textbf{0,47} & \textbf{0,31} & 1,15 & 1,03 \\
\end{longtable}

Năm phát hiện: (1) \textbf{nhà nước tô (rentier state)}, Beblawi (1987);
Hertog (2010); Hvidt (2013); (2) \textbf{phân bổ sai nguồn lực đảo
chiều} so với Hsieh \& Klenow (2009); (3) Kuwait Vision 2035 với R\&D
20,7\%; (4) \textbf{thiên lệch của DAI đơn thành phần}; (5) \textbf{phân
nhóm con (sub-grouping) Advanced}, innovation-driven so với
resource-driven là một dạng biến thể chế kiểu Varieties of Capitalism
(Hall \& Soskice, 2001). \textbf{Bối cảnh xung đột Trung Đông 2026}: IMF
(2026, April) điều chỉnh giảm tăng trưởng Saudi Arabia mạnh từ
\textasciitilde4,5\% xuống 3,1\% do giảm sản lượng dầu; ADB (2026,
April), báo cáo \emph{Asian Development Outlook April 2026: The Middle
East Conflict Challenges Resilience in Asia and the Pacific}, dự báo
Pacific 3,4\%, Đông Nam Á 4,7\%, Việt Nam 7,0\%, củng cố luận điểm về dị
biệt Advanced innovation vs resource.

\hypertarget{53-viux1ec7t-nam-n3077-3-ux111ux1ee3t-khux1ea3o-suxe1t-cux1eadp-nhux1eadt-v34-vux1edbi-adb-vietnam-outlook-2026}{%
\subsubsection{5.3 Việt Nam (n=3.077, 3 đợt khảo sát), cập nhật v3.4 với
ADB Vietnam Outlook
2026}\label{53-viux1ec7t-nam-n3077-3-ux111ux1ee3t-khux1ea3o-suxe1t-cux1eadp-nhux1eadt-v34-vux1edbi-adb-vietnam-outlook-2026}}

FSTS 23,2\%, sau đó 17,9\%, sau đó 16,1\% (suy giảm); doanh nghiệp xuất
khẩu 37,1\%, sau đó 23,8\%; ISO 17--23\%; R\&D 6,1\% (đợt 2023). Pattern
\textbf{kinh tế hai tầng (two-tier economy)}, doanh nghiệp FDI hướng
xuất khẩu hiệu quả cao đan xen với doanh nghiệp nội địa năng suất thấp
(CIEM, 2023; Tran \& Pham, 2024).

\textbf{Cập nhật bối cảnh 2026, ADB Vietnam Economic Outlook 2026}: Theo
ADB (2026, April, \emph{Vietnam Economic Outlook 2026: Navigating the
Crosscurrents}), Việt Nam ghi nhận \textbf{GDP lịch sử 2025 đạt 8,0\%}
(cao nhất khu vực Đông Nam Á) và dự phóng \textbf{7,2\% (2026), sau đó
7,0\% (2027)}, vượt mức bình quân ASEAN 4,6\% và Đông Nam Á 4,7\%. Bốn
động lực chính: (a) FDI commitment \$2,4 tỷ qua 18 dự án chiến lược; (b)
Năng suất lao động tăng 5,1\% giai đoạn 2026-2027; (c) Hai-tầng kinh tế
tiếp tục phân kỳ + ADB Policy-Based Loans \$2,4B target SME nội địa; (d)
GVC re-positioning từ assembly downstream sang mid-stream design/R\&D
theo 3 trụ cột "Resilience + Environmental Sustainability +
Inclusiveness".

\textbf{Cập nhật macro deep, 4 mảng bổ sung}:

(i) \textbf{Lạm phát}: ADB dự báo \textbf{4,0\% (2026), sau đó 3,8\%
(2027)}, nằm trong vùng mục tiêu của Ngân hàng Nhà nước (≤4,5\%). Áp lực
từ giá năng lượng + xung đột Trung Đông kéo dài + đứt gãy chuỗi cung ứng
vận tải biển. Lạm phát kiểm soát được tạo dư địa cho chính sách tiền tệ
nới lỏng.

(ii) \textbf{Cấu trúc ngành} (đóng góp tăng trưởng 2026): \textbf{công
nghiệp + xây dựng 7,7\%} (FDI manufacturing + smart manufacturing);
\textbf{dịch vụ 7,5\%} (du lịch phục hồi sau COVID, retail e-commerce,
fintech); \textbf{nông nghiệp 3,6\%} (climate-resilient agriculture,
công nghệ cao). Cơ cấu này phản ánh chuyển dịch tiếp tục từ nông nghiệp
sang công nghiệp dẫn đến dịch vụ chất lượng cao.

(iii) \textbf{Bốn trụ cột tăng trưởng}: (a) \textbf{Đầu tư công}, gói
800.000 tỷ VND giải ngân ưu tiên hạ tầng giao thông + năng lượng sạch +
chuyển đổi số chính phủ; (b) \textbf{Chính sách tiền tệ nới lỏng}, Ngân
hàng Nhà nước duy trì lãi suất tái cấp vốn thấp + mở rộng tín dụng có
chọn lọc cho SME + manufacturing; (c) \textbf{FDI inflows}, \$2,4 tỷ
commitment qua 18 dự án chiến lược (fintech, năng lượng sạch, logistics,
semiconductor packaging); (d) \textbf{Xuất khẩu}, duy trì
\textgreater30\% kim ngạch sang Mỹ (đối tác lớn nhất) + đa dạng hóa qua
CPTPP + RCEP + ACFTA.

(iv) \textbf{Năm rủi ro chính}: (a) \textbf{Gián đoạn chuỗi cung ứng
toàn cầu} do Trung Đông + chip shortage + vận tải biển; (b)
\textbf{Chính sách thuế quan Mỹ 2026} --- thuế 10\% tạm thời, nguy cơ
tăng 15\% (xem Mục 1.1 file 14 v3.11 Lớp 7); (c) \textbf{Thanh khoản nội
địa} --- nợ doanh nghiệp/GDP cao (\textasciitilde120\%); (d) \textbf{Nợ
xấu hệ thống ngân hàng} --- đặc biệt từ bất động sản; (e) \textbf{Trái
phiếu doanh nghiệp yếu} --- sau khủng hoảng Tân Hoàng Minh + Vạn Thịnh
Phát 2022-2024 --- cần củng cố niềm tin nhà đầu tư trái phiếu trong
nước.

\textbf{Hàm ý cho CĐ2}: bổ sung \textbf{biến
\texttt{Vietnam\_HighValue\_FDI\_2026}} đo bước chuyển từ assembly sang
higher-value của FDI manufacturing như test cho \textbf{H7 industry ×
institutional dynamism interaction} (Kafouros et al., 2023). Bổ sung
thêm \textbf{biến \texttt{Vietnam\_Sector\_Mix\_2026}} (3 cấp: công
nghiệp / dịch vụ / nông nghiệp) làm control variable trong robustness
check \#6 (sector composition adjustment).

\hypertarget{54-trung-quux1ed1c-n4889}{%
\subsubsection{5.4 Trung Quốc
(n=4.889)}\label{54-trung-quux1ed1c-n4889}}

FSTS 10,9\%, sau đó 8,8\%; FDI ≥10\% chiếm 6,0\%. Quan hệ FSTS -- năng
suất là đường cong bậc hai (inverted-U) với điểm uốn 49,4\% (2012) /
47,2\% (2024) --- bền vững theo kiểm định Paternoster (Đỗ \& Phan, 2026
--- \emph{IJOEM}). \textbf{Cập nhật v3.1c}: Wang, Huang và Hong (2024
--- \emph{IRFA}) phân tích 80\% bank risk models ở PRC phụ thuộc
dominant tech providers --- \textbf{concentration risk} trong digital
ecosystem PRC.

\textbf{Khung Zombie firms PRC --- lý giải bổ sung cho điểm uốn bậc hai
(inverted-U) \textasciitilde49\%}: Theo các phân tích nội bộ và văn liệu
Caballero, Hoshi và Kashyap (2008 --- \emph{American Economic Review}),
khái niệm \textbf{"doanh nghiệp xác sống" (zombie firms)} chỉ những
doanh nghiệp tồn tại nhờ tín dụng ưu đãi hoặc trợ cấp ngầm, không tạo
lợi nhuận thực + không khấu hao bình thường. Trong bối cảnh PRC sau 2015
(chiến dịch khử rủi ro tài chính + supply-side structural reform), tỷ lệ
zombie firms ở các SOE + ngành thừa năng lực sản xuất (steel, coal, real
estate, shipbuilding) vẫn cao bất chấp các đợt thắt chặt.

\textbf{Ba cơ chế zombie firms làm méo TFP tổng thể PRC}:

(i) \textbf{Chậm đào thải} (slow exit): Doanh nghiệp xác sống không phá
sản dẫn đến vẫn chiếm dụng vốn vay ngân hàng + đất đai + lao động ---
kéo dài chu kỳ phân bổ sai nguồn lực (Hsieh \& Klenow, 2009 channel).
Trong WBES PRC pool, một số firms zombie có FSTS cao (cố xuất khẩu để
duy trì doanh thu) nhưng năng suất thực thấp dẫn đến góp phần kéo điểm
uốn bậc hai về phía \textasciitilde49\%.

(ii) \textbf{Cản trở entry} (entry barrier): SME mới không thể chiếm thị
phần vì zombie firms vẫn occupy thị trường nội địa với giá dumping. Liên
kết Mục 5.5 Emerging Asia: Việt Nam + Indonesia + India tránh được
zombie effect nhờ market discipline mạnh hơn dẫn đến entry rate cao +
creative destruction (Schumpeter, 1942) hoạt động tốt hơn.

(iii) \textbf{Concentration risk + zombie hybrid}: Wang/Huang/Hong
(2024) finding (80\% bank risk models PRC dependent on dominant tech
providers) \textbf{kết hợp} với zombie firms structure tạo ra rủi ro hệ
thống --- nếu fintech provider lock-in fails + zombie firms massively
default đồng thời dẫn đến systemic banking crisis. Đây là rủi ro mới cho
China-exposed Vietnamese exporters.

\textbf{Hàm ý cho CĐ2}: bổ sung \textbf{biến
\texttt{Zombie\_Firm\_Indicator\_PRC}} (proxy: firms với ROA \textless0
trong 3 năm liên tiếp + tăng nợ vay) để test xem điểm uốn bậc hai
\textasciitilde49\% có shifted khi loại zombie firms khỏi PRC subsample.
Robustness check \#7 (zombie-excluded re-estimation) --- bổ sung sau
Manufacturing-only / ICT-excluded / Tourism-separated SIDS /
Construction-tested Gulf / Mining-excluded resource / Sector composition
adjustment.

\hypertarget{55-tux1ed5ng-hux1ee3p-emerging-asia-n42278--bux1ed5-sung-india-case-v32}{%
\subsubsection{5.5 Tổng hợp Emerging Asia (n=42.278) --- bổ sung India
case
(v3.2)}\label{55-tux1ed5ng-hux1ee3p-emerging-asia-n42278--bux1ed5-sung-india-case-v32}}

FSTS 7,6\%; FDI ≥10\% chiếm 4,4\%; tỷ lệ doanh nghiệp R\&D dương 16\%.
Phân tán nội bộ lớn (sd log 2,18) --- phản ánh dị biệt rộng xuyên 7 nước
Emerging.

\textbf{Trường hợp Ấn Độ --- chứng cứ lan tỏa công nghệ (technology
spillovers) (NEW v3.2)}: Theo Sikdar và Mukhopadhyay (2026 ---
\emph{Asian Development Review, 43}(1), pp.37-75), nghiên cứu panel
\textbf{4.293 doanh nghiệp Ấn Độ giai đoạn 2006-2019}. Phát hiện: FDI là
kênh \textbf{lan tỏa ngang (horizontal spillover)} quan trọng nhất;
doanh nghiệp lớn và hiệu quả cao thu nhận lan tỏa (capture spillovers)
nhiều nhất; \textbf{dòng vốn FDI vào (FDI inflows)} tăng từ \$1,705M
(1990-2000) dẫn đến \$42,396M (2012-2022); chi tiêu nghiên cứu-phát
triển (R\&D) trì trệ ở 0,7\% GDP; \textbf{đơn đăng ký sáng chế của cư
dân (resident patent filing)} tăng +11,9\%/năm 2012-2022.

\hypertarget{56-mongolia--tux1eeb-lux1eddi-nguyux1ec1n-tuxe0i-nguyuxean-ux111ux1ebfn-chuyux1ec3n-ux111ux1ed5i-critical-minerals-n1905-4-ux111ux1ee3t-20092025}{%
\subsubsection{5.6 Mongolia --- từ lời nguyền tài nguyên đến chuyển đổi
critical minerals (n=1.905, 4 đợt
2009--2025)}\label{56-mongolia--tux1eeb-lux1eddi-nguyux1ec1n-tuxe0i-nguyuxean-ux111ux1ebfn-chuyux1ec3n-ux111ux1ed5i-critical-minerals-n1905-4-ux111ux1ee3t-20092025}}

FSTS đình trệ 4-6\%; FDI ≥10\% giảm rõ \textbf{7,2\%, sau đó 3,2\%}; tỷ
lệ website tăng 39\%, sau đó 65\%. Pattern \textbf{lời nguyền tài
nguyên} (Auty, 1993; Sachs \& Warner, 2001).

\textbf{Bước ngoặt 2026 --- chuyển đổi critical minerals (ADB, 2026,
May)}: Mongolia ở giao điểm chiến lược của \textbf{11 critical minerals}
(Cu, Li, Mo, Mn, Co, Ni, W, REE, etc.). Cu exports 1,69 Mt (2024) =
22,1\% tổng xuất khẩu. Oyu Tolgoi target 500.000 tấn/năm 2028-2036. PRC
processing dominance: 100\% graphite, 90\% Mn, 70\% Co. CĐ2 thêm biến
\texttt{Critical\_Minerals\_Exposure}.

\begin{quote}
\emph{Hình 5.6.1, Mongolia: Evolution 2009--2025 (4 đợt khảo sát WBES).
Hình sẽ được sinh từ
\texttt{generate\_figures.py::fig\_5\_6\_mongolia()} trong bản nộp cuối;
PENDING.}
\end{quote}

\hypertarget{57-sids-thuxe1i-buxecnh-dux1b0ux1a1ng-7-nux1b0ux1edbc--cux1eadp-nhux1eadt-v31-n1371}{%
\subsubsection{5.7 SIDS Thái Bình Dương (7 nước --- cập nhật v3.1,
n=1.371)}\label{57-sids-thuxe1i-buxecnh-dux1b0ux1a1ng-7-nux1b0ux1edbc--cux1eadp-nhux1eadt-v31-n1371}}

\textbf{Bảy nước}: Fiji (FJI), Papua New Guinea (PNG), Solomon Islands
(SLB), Tonga (TON), Vanuatu (VUT), Samoa (WSM), \textbf{Kiribati (KIR
--- bổ sung 2025)}. Tổng n=1.371. FSTS 6,3\%; FDI ≥10\% \textbf{23,5\%};
đổi mới sản phẩm \textbf{41,5\%}; website \textbf{58,9\%}. Pattern
\textbf{thích nghi trong điều kiện ràng buộc (adaptation under
constraint)}.

\textbf{Kiribati 2025 --- trường hợp biên CỰC ĐOAN nhất}: FSTS 1,03\%,
FDI 0,7\%, website 18,7\%, ISO 1,3\%, sd log 1,48, SME 93,3\%. GDP/đầu
người \textasciitilde1.700 USD, viện trợ AUS/NZ \textasciitilde30\% GDP.

\hypertarget{575-bux1ed1i-cux1ea3nh-kinh-tux1ebf-vux129-muxf4-pics-pacific}{%
\paragraph{5.7.5 Bối cảnh kinh tế vĩ mô PICs
Pacific}\label{575-bux1ed1i-cux1ea3nh-kinh-tux1ebf-vux129-muxf4-pics-pacific}}

PICs có 4 trụ cột MIRAB (Bertram, 2006): (1) khu vực công chủ đạo
(30-50\% việc làm chính thức); (2) phụ thuộc 4 nguồn ngoại sinh --- nông
nghiệp tự cấp + ODA 30\% GDP + remittances 25-40\% GDP + khai thác tài
nguyên; (3) thanh niên thất nghiệp 15-30\% + dân số \textless25 tuổi
\textgreater50\%; (4) migration labor PLMS Australia + RSE New Zealand
như channel sinh tồn (Vanuatu RSE 15-20\% GDP). Liên kết: Fiji website
74,8\% \textgreater{} Singapore (PLMS+remittance corridors); Kiribati
18,7\% (outer islands no internet); SIDS innovation 41,5\% (bricolage
Baker \& Nelson 2005). CĐ2: 3 control variables PICs-specific
(\texttt{aid\_share\_GDP\_pct}, \texttt{remittance\_share\_GDP\_pct},
\texttt{seasonal\_worker\_outflow\_pct}).

\hypertarget{58-so-suxe1nh-tux1ed5ng-hux1ee3p}{%
\subsubsection{5.8 So sánh tổng
hợp}\label{58-so-suxe1nh-tux1ed5ng-hux1ee3p}}

\textbf{Bảng 5.1}. \emph{Ma trận so sánh đa chiều bảy tiểu cảnh điển
hình.}

\begin{longtable}[]{@{}llllllll@{}}
\toprule\noalign{}
Chỉ số & Singapore & Việt Nam & Trung Quốc & Em Asia & Mongolia & SIDS
Thái Bình Dương¹ & Saudi+Qatar+Kuwait \\
\midrule\noalign{}
\endhead
\bottomrule\noalign{}
\endlastfoot
n\_firms & 623 & 3.077 & 4.889 & 42.278 & 1.905 & \textbf{1.371 (v3.1)}
& 1.632 \\
FSTS (\%) & 7,1 & 19,1 & 9,9 & 7,6 & \textbf{5,0} & 6,3 &
\textbf{2,4} \\
FDI ≥10\% (\%) & \textbf{31,5} & 11,4 & 6,0 & 4,4 & \textbf{4,7} & 23,5
& \textbf{11,5} \\
R\&D dương (\%) & 7,5 & \textbf{3,1} & \textbf{39,4} & 16 &
\textbf{20,8} & 11,8 & \textbf{3,1} \\
Website (\%) & 66,1 & \textbf{47,0} & \textbf{58,7} & \textbf{48,0} &
\textbf{50,1} & 58,9 & \textbf{43,6} \\
sd log năng suất & 1,03 & 1,38 & 1,20 & 2,18 & \textbf{1,16} & 1,32 &
\textbf{0,49} \\
\end{longtable}

\begin{center}\rule{0.5\linewidth}{0.5pt}\end{center}

\hypertarget{chux1b0ux1a1ng-6--cuxe1c-yux1ebfu-tux1ed1-giux1ea3i-thuxedch-sux1a1-bux1ed9}{%
\subsection{CHƯƠNG 6 --- CÁC YẾU TỐ GIẢI THÍCH SƠ
BỘ}\label{chux1b0ux1a1ng-6--cuxe1c-yux1ebfu-tux1ed1-giux1ea3i-thuxedch-sux1a1-bux1ed9}}

\textbf{Bảng 6.1}. \emph{Hệ số tương quan Pearson (n=101.185 --- v3.1).}

\begin{longtable}[]{@{}llllll@{}}
\toprule\noalign{}
Yếu tố & Advanced & Upper-middle & Emerging & Frontier & SIDS \\
\midrule\noalign{}
\endhead
\bottomrule\noalign{}
\endlastfoot
FDI ≥10\% & −0,113 & −0,023 & +0,113 & +0,068 & \textbf{+0,222} \\
FSTS & \textbf{+0,113} & −0,086 & +0,062 & +0,014 & −0,050 \\
TCI (R\&D + ISO) & \textbf{+0,128} & +0,016 & −0,029 & +0,019 &
\textbf{+0,155} \\
DAI (website) & \textbf{−0,129} & +0,012 & +0,016 & +0,070 & −0,049 \\
\end{longtable}

\textbf{Sáu phát hiện chính}: (1) FDI dương mạnh SIDS (+0,222); (2) FDI
âm Advanced; (3) FSTS dương Advanced, Uppsala+OLI; (4) TCI dương
SIDS+Advanced; (5) DAI âm Advanced, Tier-1 saturation; (6) Đảo dấu
cross-regime --- institutional moderation.

\textbf{Cảnh báo schema BREADY 2025 (v3.3)}: bảng 6.1 có thể bị nhiễu
schema effect --- kiểm chứng qua "triple-defense" Mục 4.11 file 15.

\textbf{Lý giải dị biệt cross-regime qua Xu (2024) de jure-de facto
(v3.3)}: gap formal-enforcement explains FDI sign reversal Emerging vs
Advanced.

\textbf{Lý giải bổ sung --- Institutional Transaction Costs xuyên quốc
gia (TCE Coase-Williamson, mở rộng Banalieva \& Dhanaraj 2019)}: Đối
chiếu hai bản thảo P3 Singapore (Mar et al., 2026 --- \emph{MIR}) + P4
Việt Nam (Đỗ \& Phan, 2026 --- \emph{APJM}) cho thấy dị biệt
cross-regime trong Bảng 6.1 không chỉ phản ánh chất lượng thể chế formal
mà còn \textbf{chi phí giao dịch thể chế (institutional transaction
costs)} đo lường hạ tầng thương mại + logistics + năng lực hấp thụ chi
phí điều phối xuyên biên giới (cross-border coordination cost
absorption). Khung CDCM (Mục 2.7 file 14 v3.13) hợp nhất 3 chiều xuyên
quốc gia:

\textbf{Bảng 6.1.2}. \emph{Đối chiếu I-P curve + cơ chế DAI/TCI giữa
Việt Nam vs Singapore --- bằng chứng cho khung CDCM.}

\begin{longtable}[]{@{}lll@{}}
\toprule\noalign{}
Chiều & \textbf{Việt Nam (Emerging, n=2.958, 3 wave)} &
\textbf{Singapore (Advanced innovation, n=623, 1 wave)} \\
\midrule\noalign{}
\endhead
\bottomrule\noalign{}
\endlastfoot
\textbf{Hình dáng I-P curve} & \textbf{Inverted-U rõ nét} --- điểm uốn
(TP) ≈ 39-46\% FSTS xuyên 3 wave 2009-2023; Lind-Mehlum p\textless0,013
& \textbf{Tuyến tính dương + bậc hai nhẹ} --- TP ≈ 82\% FSTS (vùng dữ
liệu thưa, 3,2\% firms vượt 70\% FSTS); Lind-Mehlum p=0,303 KHÔNG xác
nhận \\
\textbf{Cơ chế DAI} & \textbf{Stage-contingent} (Tier 1 = website only):
mạnh 2009 → null 2015 → âm 2023 (β=-0,912, p=,043) ở DAI×FSTS →
\textbf{điểm nghẽn ở xuất khẩu cao} & \textbf{Conditional scaling} (Tier
1+2 = website + e-payment): DAI×FSTS² β=3,119 p=,005 dương mạnh →
\textbf{đòn bẩy ở xuất khẩu cao} \\
\textbf{2SLS DAI robustness} & β co về 0,02, p=,94 → \textbf{không xác
lập nhân quả} (Đỗ \& Phan, 2026 --- \emph{APJM}) & (chưa kiểm định
IV) \\
\textbf{Vai trò TCI} & \textbf{Scarce advantage} --- TCIz IV β=1,639
p\textless,001 robust 2SLS; điều tiết làm phẳng I-P curve (3/4 mô hình
significant) & \textbf{Hygiene factor} --- direct β=0,168 p\textless,001
nhưng KHÔNG điều tiết I-P curve significantly \\
\end{longtable}

\textbf{Lý giải Institutional Transaction Costs (TCE)}: Ở Việt Nam, hạ
tầng thương mại số yếu (Tier 1 only) khiến chi phí điều phối xuyên biên
giới (cross-border coordination cost) tăng vọt theo siêu tuyến tính
(super-linearly --- Brynjolfsson \& McAfee, 2014) ở mức xuất khẩu trung
bình dẫn đến đẩy điểm uốn (TP) về vùng FSTS thấp 39-46\%. Ở Singapore,
hạ tầng số (Tier 1+2: website + e-payment + integrated logistics) đã hấp
thụ phần lớn chi phí này dẫn đến đẩy ranh giới quá tải (overload
boundary) về vùng FSTS cực cao \textasciitilde82\%. Vì vậy \textbf{bằng
chứng "đường cong U ngược nghẹt thở" của Lu \& Beamish (2004) chỉ vận
hành ở các nước có institutional transaction costs trung bình-cao} ---
nơi chi phí điều phối phi tuyến đủ tạo turnaround point trong vùng FSTS
observed; ở các nước có hạ tầng số trưởng thành, U-curve được "kéo dãn"
về phía cực đoan và có thể không quan sát được trong dữ liệu thực tế.

\textbf{Hàm ý cho khung CDCM (Context-Contingent Digital Capability
Model)}: Khung này dự đoán trong CĐ2 + luận án 5 chương: (a) \textbf{41
nước châu Á} sẽ có một continuum institutional transaction costs từ
Singapore (cực thấp) dẫn đến Việt Nam/Indonesia (trung bình) dẫn đến
Mongolia/Pakistan/Bangladesh (cao); (b) Điểm uốn TP của I-P curve sẽ
correlate ngược với mức trưởng thành hạ tầng số quốc gia (đo qua DAI
Tier 1+2 trung bình quốc gia); (c) Vai trò TCI sẽ chuyển từ scarce
advantage (ở các nước transaction cost cao) sang hygiene factor (ở các
nước transaction cost thấp) --- confirm giả thuyết H7 (industry ×
institutional dynamism interaction) của Kafouros et al. (2023). Đây là
\textbf{đóng góp lý thuyết MỚI} cho luận án --- replication CDCM xuyên
41 nước châu Á.

\textbf{Tham chiếu}: Coase (1937); Williamson (1985); Brynjolfsson \&
McAfee (2014); Banalieva \& Dhanaraj (2019); Kafouros et al. (2023); Mar
et al. (2026, P3 Singapore manuscript, \emph{MIR} under review); Đỗ \&
Phan (2026, P4 Việt Nam --- \emph{APJM}); Mục 2.7 file 14 v3.13 (Khung
4-Tier Verhoef + CDCM).

\textbf{Hiệu ứng đẩy nhanh giao hàng + bất định chính sách thuế quan
trong khung ICRV (front-loading effect + tariff uncertainty)}: Theo phân
tích nội bộ, bất định chính sách thuế quan Mỹ 2026 (xem Mục 1.1 file 14
v3.11 Lớp 7) tạo \textbf{hiệu ứng đẩy nhanh giao hàng (front-loading
effect)} --- doanh nghiệp xuất khẩu Việt Nam đẩy nhanh giao hàng ngắn
hạn nhằm tận dụng cửa sổ thuế thấp 10\% trước khi tăng lên 15\% trong
nửa cuối 2026. Hậu quả: doanh thu xuất khẩu tăng ngắn hạn 2025 (đóng góp
vào GDP 8,0\%) nhưng làm giảm đà tăng trưởng 2026-2027 do hợp đồng tương
lai bị "ăn trước". Trong khung ICRV, đẩy nhanh giao hàng + bất định
chính sách hoạt động như \textbf{cú sốc ngoại sinh (exogenous shock)}
đảo dấu hệ số FSTS giữa hai nhóm: (a) \textbf{Emerging năng lực FDI cao}
(Việt Nam doanh nghiệp FDI sản xuất chế biến) --- có khả năng đẩy nhanh
+ phòng ngừa rủi ro tốt dẫn đến FSTS-năng suất tương quan dương mạnh
ngắn hạn 2025 dẫn đến đảo chiều âm 2026-2027 do \textbf{cạn kiệt đơn
hàng tương lai (exhaustion of forward orders)}; (b) \textbf{Emerging
năng lực FDI thấp} (doanh nghiệp nội địa SME không có \textbf{tồn kho dự
phòng (buffer inventory)}) --- không thể đẩy nhanh dẫn đến FSTS-năng
suất tương quan yếu hoặc âm cả 2 giai đoạn. \textbf{Phân hóa lớn vs
SME}: tập đoàn lớn hấp thụ chi phí thuế quan qua đa thị trường + phòng
ngừa tài chính (financial hedging); SME phải lựa chọn hấp thụ chi phí
(giảm biên lợi nhuận) vs chuyển giá khách hàng (mất đơn hàng). Hàm ý
CĐ2: bổ sung \textbf{biến \texttt{Tariff\_Front\_Loading\_2025\_2026}}
(proxy: tỷ trọng xuất khẩu sang Mỹ × thời điểm giao hàng quý 1-2 2026)
vào kiểm định độ vững (robustness check) \#8 (đẩy nhanh giao hàng được
hiệu chỉnh --- front-loading-adjusted re-estimation).

\hypertarget{61-pattern-giux1edbi-tuxednh-trong-quux1ea3n-luxfd--sux1edf-hux1eefu--bux1eb1ng-chux1ee9ng-eap-chuxe2u-uxe1-dux1eabn-ux111ux1ea7u}{%
\subsubsection{6.1 Pattern giới tính trong quản lý + sở hữu --- bằng
chứng EAP châu Á dẫn
đầu}\label{61-pattern-giux1edbi-tuxednh-trong-quux1ea3n-luxfd--sux1edf-hux1eefu--bux1eb1ng-chux1ee9ng-eap-chuxe2u-uxe1-dux1eabn-ux111ux1ea7u}}

\textbf{Bảng 6.1.1}. \emph{Pattern giới tính trong TMT + ownership ---
phân tầng theo khu vực địa lý.}

\begin{longtable}[]{@{}llll@{}}
\toprule\noalign{}
Khu vực & \% nữ TMT (cấp cao) & \% nữ ownership majority & Khoảng cách
(TMT -- Owner) \\
\midrule\noalign{}
\endhead
\bottomrule\noalign{}
\endlastfoot
\textbf{East Asia \& Pacific (EAP)} & \textbf{33,4\%} \emph{(dẫn đầu
globally)} & \textbf{25,7\%} \emph{(dẫn đầu globally)} & +7,7 đpt \\
Latin America \& Caribbean & 27,8\% & 18,3\% & +9,5 đpt \\
Sub-Saharan Africa & 21,3\% & 14,9\% & +6,4 đpt \\
Europe \& Central Asia & 19,8\% & 12,1\% & +7,7 đpt \\
South Asia & 9,2\% & 8,5\% & +0,7 đpt \\
\textbf{Middle East \& North Africa (MENA)} & \textbf{3,3\%} \emph{(thấp
nhất globally)} & \textbf{1,7\%} \emph{(thấp nhất globally)} & +1,6
đpt \\
\end{longtable}

Sáu quan sát: (a) \textbf{EAP dẫn đầu, Asian gender paradox}: EAP 33,4\%
nữ TMT cao nhất toàn cầu trong khi South Asia và MENA thấp nhất, phản
ánh dị biệt thể chế-văn hóa xuyên "châu Á" không đồng nhất; (b)
\textbf{MENA thấp nhất, structural barriers}: MENA 3,3\% nữ TMT + 1,7\%
nữ ownership, thấp nhất toàn cầu, phản ánh rào cản thể chế-luật pháp đặc
thù (World Bank, 2024 --- \emph{Women, Business and the Law}); (c)
\textbf{Ownership ceiling phổ biến}: 43/50 economies có ownership
ceiling cao hơn management representation --- gợi ý barrier đặc thù ở
cấp sở hữu vốn (equity ownership) so với vị trí quản lý; (d)
\textbf{Boundary references TMT moderation}: Mardones-Ibáñez (2025 ---
\emph{SAGE Open}) Chile context + Al-Najjar et al. (2025, \emph{IJHRM})
S\&P 500 U-shaped TMT gender effect, không suy rộng trực tiếp sang ICRV
Emerging/Frontier; (e) \textbf{Technology complement hypothesis}
(Carboni et al., WBES 192K obs, 158 nước, 2006--2023): female top
manager (\texttt{b7a}) tương tác với technology adoption --- tác động
ESG-performance mạnh hơn khi doanh nghiệp có DAI cao; (f) \textbf{Firm
performance heterogeneity} (World Bank, 2024 --- \emph{Unlocking Global
Growth: Closing the Gender Gap in Business}): doanh nghiệp có nữ sở hữu
đa số có doanh thu/nhân viên thấp hơn trung bình nhưng tốc độ tăng
trưởng việc làm cao hơn ở nhóm Frontier/Emerging --- pattern phù hợp Mục
4.3 CAGR việc làm SIDS cao nhất (5,77\%).

\textbf{Liên kết với Mục 4.5.6}: Chi tiết phân tầng female top manager
theo ICRV sub-regime và hàm ý biến \texttt{b7a} trong CĐ2 specification
xem Mục 4.5.6 file 15.

\textbf{Nguồn}: Carboni et al. (2024/2025 --- WBES ESG indicators);
World Bank (2024) \emph{Unlocking Global Growth: Closing the Gender Gap
in Business} (WBES 167 nền kinh tế); Mardones-Ibáñez (2025 ---
\emph{SAGE Open}); Al-Najjar et al. (2025 --- \emph{IJHRM}).

\hypertarget{62-ux111ux1b0ux1eddng-cong-u-cux1ee7a-innovation-digital-tech}{%
\subsubsection{6.2 Đường cong U của innovation digital
tech}\label{62-ux111ux1b0ux1eddng-cong-u-cux1ee7a-innovation-digital-tech}}

\textbf{Giai đoạn 1 (chi phí đầu tư cao)}: 1-3 năm đầu, giảm hiệu quả
tài chính ngắn hạn --- phù hợp DAI âm -0,129 Advanced + "digital
theatre" (Verhoef et al., 2021). \textbf{Giai đoạn 2 (productivity
payoff)}: vượt threshold \textasciitilde3\% revenue/năm × 3-5 năm
sustained dẫn đến tăng năng suất + cạnh tranh quốc tế + profit margin.
Liên kết Tier-1/Tier-2 Mục 4.4.5 file 15. Bằng chứng: Trung Quốc R\&D
39,4\% late-stage Tier-2; Vùng Vịnh transition phase; Việt Nam R\&D
6,1\% early Tier-2.

\begin{center}\rule{0.5\linewidth}{0.5pt}\end{center}

\hypertarget{chux1b0ux1a1ng-7--khoux1ea3ng-trux1ed1ng-thux1ef1c-tiux1ec5n-vuxe0-kux1ebft-luux1eadn}{%
\subsection{CHƯƠNG 7 --- KHOẢNG TRỐNG THỰC TIỄN VÀ KẾT
LUẬN}\label{chux1b0ux1a1ng-7--khoux1ea3ng-trux1ed1ng-thux1ef1c-tiux1ec5n-vuxe0-kux1ebft-luux1eadn}}

\hypertarget{71-khoux1ea3ng-trux1ed1ng-nghiuxean-cux1ee9u-thux1ef1c-tiux1ec5n}{%
\subsubsection{7.1 Khoảng trống nghiên cứu thực
tiễn}\label{71-khoux1ea3ng-trux1ed1ng-nghiuxean-cux1ee9u-thux1ef1c-tiux1ec5n}}

Năm khoảng trống chưa được lấp đầy trong văn liệu:

\textbf{KT1 --- Three-way moderation TCI×DAI×FSTS chưa được kiểm định
đồng thời trên mẫu liên quốc gia}: P3/P4/P5 kiểm định từng cặp điều tiết
riêng biệt. Mô hình M7 CĐ2 sẽ lấp khoảng trống này với 108 cặp quốc
gia--năm.

\textbf{KT2 --- Phân nhóm con Advanced: innovation-driven vs
resource-driven chưa được hệ thống hóa}: Nghiên cứu hiện hành (Bausch \&
Krist, 2007; Kirca et al., 2012) gộp tất cả OECD vào "developed". Bằng
chứng Singapore vs Saudi/Qatar/Kuwait cho thấy hai sub-group có pattern
hoàn toàn khác nhau.

\textbf{KT3 --- Xác nhận nhân quả Digital Saturation Paradox qua IV
estimation}: Hệ số DAI âm (−0,129) có thể phản ánh: (a) true saturation;
(b) reverse causality; (c) omitted variable bias. IV instrument hợp lệ
(broadband coverage) cần thiết để phân biệt.

\textbf{KT4 --- Temporal heterogeneity ba giai đoạn 2009--2025 chưa được
meta-analyzed}: Bằng chứng P3 (Paternoster z = 3,353, p \textless{}
0,001) và P5 (TCI Paternoster p = 0,011) cho thấy parameter instability
quan trọng.

\textbf{KT5 --- SIDS Pacific như trường hợp biên lý thuyết của thuyết
Uppsala}: Forced internationalization penalty tại 7 SIDS không khớp với
bất kỳ dạng hàm I--P trong văn liệu. Cần khung kết hợp MIRAB + Briguglio
(1995) + RBV.

\hypertarget{72-kux1ebft-luux1eadn-chuxednh}{%
\subsubsection{7.2 Kết luận
chính}\label{72-kux1ebft-luux1eadn-chuxednh}}

Tám kết luận từ 101.185 doanh nghiệp, 47 nền kinh tế, 108 cặp quốc
gia--năm, 14 sóng WBES 2009--2025:

\textbf{KL1} --- Phi tuyến xuyên chế độ thể chế là quy luật: Inverted-U
tại Emerging (Việt Nam TP 39--46\%; Trung Quốc TP 47--49\%); gần tuyến
tính tại Advanced (Singapore TP 88,6\%); forced penalty tại SIDS.

\textbf{KL2} --- TCI là level-shifter phổ quát, DAI là
context-contingent amplifier: TCI dương ở mọi nhóm; DAI đổi dấu theo
ICRV regime --- đây là nội hàm cốt lõi của CDCM.

\textbf{KL3} --- Digital Saturation Paradox là cơ chế đo lường, không
phải thất bại thị trường: DAI âm tại Advanced là artifact của Tier-1
measurement trong bối cảnh Tier-2/3 đã được áp dụng.

\textbf{KL4} --- Thể chế là moderator, không phải mediator: de jure--de
facto gap (Xu, 2024) giải thích đảo dấu FDI cross-regime. Chính sách
nâng cao thể chế hiệu quả nhất khi song hành với nâng cấp TCI và DAI.

\textbf{KL5} --- Temporal heterogeneity ba giai đoạn là thực tế cấu
trúc: DAI effect mạnh hơn trong BREADY (sau COVID-19 digital
acceleration). Triple-defense system đảm bảo kết luận không bị nhiễu
schema.

\textbf{KL6} --- EAP dẫn đầu toàn cầu về nữ quản lý cấp cao (33,4\%)
nhưng khoảng cách sở hữu dai dẳng: Glass ceiling ở lớp sở hữu chưa được
phá vỡ dù management ceiling đang thu hẹp.

\textbf{KL7} --- SIDS Pacific là trường hợp biên lý thuyết với forced
internationalization penalty: FDI (+0,222) là nguồn năng suất chủ đạo;
FSTS âm. MIRAB model phù hợp hơn Uppsala.

\textbf{KL8} --- Phân nhóm con Advanced là cần thiết cho phân tích IB
châu Á: Singapore và Saudi/Qatar/Kuwait cùng nhãn "Advanced" nhưng
pattern năng suất doanh nghiệp hoàn toàn khác nhau.

\hypertarget{73-huxe0m-uxfd-cho-luux1eadn-uxe1n-chuxednh-suxe1ch-vuxe0-chuyuxean-ux111ux1ec1-2}{%
\subsubsection{7.3 Hàm ý cho luận án, chính sách và Chuyên đề
2}\label{73-huxe0m-uxfd-cho-luux1eadn-uxe1n-chuxednh-suxe1ch-vuxe0-chuyuxean-ux111ux1ec1-2}}

\hypertarget{731-huxe0m-uxfd-luxfd-thuyux1ebft}{%
\paragraph{7.3.1 Hàm ý lý
thuyết}\label{731-huxe0m-uxfd-luxfd-thuyux1ebft}}

\textbf{(1) Uppsala có điều kiện --- không phổ quát}: FSTS dương tại
Advanced, không có ý nghĩa tại Frontier xác nhận boundary condition: lý
thuyết Uppsala vận hành chỉ trong môi trường thể chế đủ mạnh
(Williamson, 1985).

\textbf{(2) CDCM như khung lý thuyết tích hợp mới cho IB digital era}:
Phần bù năng suất của digital adoption là hàm của (i) mức độ phát triển
thể chế, (ii) vị trí trên đường cong I--P, (iii) tier của digital
capability được đo lường. Ba điều kiện tương tác phi tuyến --- advance
so với "digital premium" đồng nhất (Verhoef et al., 2021).

\textbf{(3) Advanced Innovation-Performance Interface (AIPI)}: FSTS²×DAI
= +3,119 tại Singapore gợi ý cơ chế mới: tại FSTS cao, doanh nghiệp cần
DAI như "phối hợp kỹ thuật số xuyên biên giới" để quản lý độ phức tạp
GVC. Cần kiểm định tại Advanced khác (Hàn Quốc, Đài Loan).

\textbf{(4) Digital theatre là thực tế đo lường ở cả hai chiều}: Tại
Advanced, DAI Tier-1 là hygiene factor bão hòa (Verhoef et al., 2021).
Tại Emerging, DAI Tier-1 có thể là proxy obsolescence --- đo lường không
còn chụp được cơ chế kinh tế thực.

\textbf{(5) Institutional context là contingency layer, không phải
background factor}: Đảo dấu cross-regime không phải noise mà là tín hiệu
lý thuyết. Nghiên cứu IB tương lai nên kiểm định moderation effects theo
ICRV regime trước khi báo cáo average effect size.

\hypertarget{732-huxe0m-uxfd-phux1b0ux1a1ng-phuxe1p-luux1eadn}{%
\paragraph{7.3.2 Hàm ý phương pháp
luận}\label{732-huxe0m-uxfd-phux1b0ux1a1ng-phuxe1p-luux1eadn}}

\textbf{(1) Pooled cross-national WBES như data platform cho IB emerging
markets}: 101.185 doanh nghiệp từ 47 nền kinh tế là mẫu lớn nhất từ
trước đến nay cho I--P analysis toàn diện tại châu Á--Thái Bình Dương.

\textbf{(2) Triple-defense system cho mixed-generation panel là đóng góp
phương pháp luận}: Ba tầng kiểm chứng, (3a) Post\_BREADY\_2024 dummy,
(3b) anchor model stability test, (3c) independent post-pandemic panel,
là giao thức có thể chuẩn hóa cho bất kỳ nghiên cứu dùng WBES với nhiều
thế hệ schema.

\textbf{(3) Winsorization within country-year clusters tốt hơn overall
winsorization}: Overall winsorization bóp méo phân phối tại Frontier (sd
log 2,18+). Within-cluster bảo tồn dị biệt cross-regime quan trọng về lý
thuyết.

\textbf{(4) HC1 robust SE như chuẩn mực cho WBES multi-country
regression}: HC3 cần thiết cho subsample nhỏ (SIDS n \textless{} 200).
HC1 là cân bằng tốt giữa hiệu suất ước lượng và độ bền vững cho full
sample.

\textbf{(5) Anchor model strategy bảo vệ inference khi schema chuyển
đổi}: 6 biến bất biến cross-schema (ln\_employees, FSTS, ISO,
innovate\_product, website, ln\_LP) cho phép so sánh hệ số cấu trúc giữa
PICS3 và BREADY.

\hypertarget{733-huxe0m-uxfd-chuxednh-suxe1ch-cho-doanh-nghiux1ec7p-viux1ec7t-nam-vuxe0-khu-vux1ef1c-chuxe2u-uxe1}{%
\paragraph{7.3.3 Hàm ý chính sách cho doanh nghiệp Việt Nam và khu vực
châu
Á}\label{733-huxe0m-uxfd-chuxednh-suxe1ch-cho-doanh-nghiux1ec7p-viux1ec7t-nam-vuxe0-khu-vux1ef1c-chuxe2u-uxe1}}

\textbf{(1) Capability-first sequencing --- nâng cấp TCI trước khi mở
rộng FSTS}: β TCI = 0,179 (IV-causal, P3 Việt Nam) xác nhận: xuất khẩu
chỉ tạo năng suất khi doanh nghiệp có năng lực công nghệ nền tảng.
Chương trình hỗ trợ xuất khẩu cần song hành với nâng cấp TCI.

\textbf{(2) Ưu tiên Tier-2 digital cho SME Việt Nam}: Website 47\% là
điểm đầu vào; VietQR + TikTok Shop + Shopee Asia Pacific là
Tier-1.5/Tier-2 thực tế hơn. AI Adoption Voucher (\textasciitilde50
triệu VND/firm) giúp vượt qua giai đoạn đầu tư ban đầu.

\textbf{(3) Phân biệt FDI chiến lược với FDI lắp ráp}: FSTS suy giảm
23,2\%, sau đó 16,1\% phản ánh assembly FDI tăng nhưng giá trị gia tăng
nội địa thấp. KPI cần theo dõi: tỷ lệ R\&D/revenue, tỷ lệ nhà cung cấp
nội địa Tier-1, patent co-filed.

\textbf{(4) GVC repositioning: assembly downstream sang mid-stream
design/R\&D}: Việt Nam ở turning point (TP 39--46\%). Trung Quốc R\&D
39,4\% và TP cao hơn (47--49\%) cho thấy: TCI cao dẫn đến TP cao dẫn đến
chịu đựng được mức độ QTH cao hơn.

\textbf{(5) ICRV-based Policy Learning Platform cho ASEAN}: Trao đổi
best practices giữa nền kinh tế cùng nhóm IV Emerging (Việt Nam,
Indonesia, Philippines, Ấn Độ) thay vì học từ Singapore/Hàn Quốc (nhóm I
--- context khác biệt quá lớn).

\textbf{(6) Chính sách số hóa phân cấp theo vùng và quy mô}: DAI effect
phụ thuộc context --- đầu tư website khác nhau tại TP.HCM (Tier-1.5) vs
nông thôn đồng bằng sông Cửu Long (Tier-1 chưa bão hòa).

\textbf{(7) Migration labor partnership Pacific cho Việt Nam --- bài học
PLMS/RSE}

(a) \textbf{PLMS Australia}: Học mô hình 30.000 lao động PICs/năm dẫn
đến triển khai Vietnamese Skilled Worker Mobility Scheme (nông nghiệp
công nghệ cao + aged-care); (b) \textbf{RSE New Zealand}: Aquaculture +
Coffee + Cashew Skills Exchange --- tận dụng chuyên môn đặc thù Việt
Nam; (c) \textbf{Remittance corridor digital}: Hợp tác Wise/MoMo giảm
phí từ \textasciitilde7\% xuống \textless{} 3\%; (d) \textbf{Skills
upgrading theo Korea EPS}: Đào tạo kỹ năng trước khi xuất khẩu lao động
+ return migration channel.

\textbf{(8) Chính sách số hóa Asia cho SME Việt Nam --- 4 lộ trình tham
khảo}

(a) Singapore/Hàn Quốc AI full-stack (Tier-3) --- không phù hợp để copy
trực tiếp; (b) Nhật Bản 60\% DN AI back-office (Tier-2) ---
\textbf{realistic target 2026--2030}; (c) Trung Quốc smart manufacturing
(Tier-2/3) --- học positive, tránh concentration risk (Wang/Huang/Hong,
2024); (d) \textbf{Mô hình ưu tiên Việt Nam Tier-1.5}: VietQR (100\%
miễn phí) dẫn đến Social commerce dẫn đến Cross-border Shopee/Amazon dẫn
đến Cloud + AI basic. \textbf{Hàm ý CĐ2}: biến
\texttt{Digital\_Tier\_Vietnam} (4 cấp) test productivity payoff.

\hypertarget{734-huxe0m-uxfd-cho-chuyuxean-ux111ux1ec1-2--8-robustness-checks-bux1eaft-buux1ed9c}{%
\paragraph{7.3.4 Hàm ý cho Chuyên đề 2 --- 8 robustness checks bắt
buộc}\label{734-huxe0m-uxfd-cho-chuyuxean-ux111ux1ec1-2--8-robustness-checks-bux1eaft-buux1ed9c}}

\textbf{(1) Thay DV}: ln(LP) dẫn đến ROS, revenue growth, employment
growth. Kết quả chính cần bền vững ≥ 2/3 DV.

\textbf{(2) Lọc mẫu}: Loại micro firms (\textless{} 5 lao động) + SOE.
Kết quả cần bền vững trong SME tư nhân (n ≈ 80.000+).

\textbf{(3) Thay phương pháp}: OLS dẫn đến quantile regression
(p25/p50/p75); bootstrap SE 500 replications; HC3.

\textbf{(4) Functional form test}: LOWESS phi tham số; Lind--Mehlum
(2010) test bắt buộc xác nhận inverted-U; AIC/BIC so sánh linear vs
quadratic vs cubic.

\textbf{(5) Placebo tests}: Thay FSTS + TCI bằng noise variables. Kết
quả chính không được có ý nghĩa thống kê.

\textbf{(6) Nhạy cảm ICRV cutoffs}: ±0,1 WGI unit; PPP GDP/capita thay
GNI; jackknife loại từng quốc gia (47 lần).

\textbf{(7) Temporal robustness}: Chow test parameter stability;
Paternoster z-test so sánh turning points giữa ba giai đoạn.

\textbf{(8) Schema boundary test}: Triple-defense Mục 4.11 --- điều kiện
tiên quyết trước khi báo cáo BREADY 2025 như bằng chứng nhân quả.

\hypertarget{74-hux1ea1n-chux1ebf-cux1ee7a-chuyuxean-ux111ux1ec1}{%
\subsubsection{7.4 Hạn chế của chuyên
đề}\label{74-hux1ea1n-chux1ebf-cux1ee7a-chuyuxean-ux111ux1ec1}}

Bảy hạn chế cấu trúc cần được thừa nhận minh bạch:

\textbf{(1) Cross-sectional không cho phép suy luận nhân quả mạnh}: WBES
là repeated cross-sections, không phải panel cân bằng. Kết luận nhân quả
dựa trên IV từ P3 Việt Nam và không thể khái quát trực tiếp sang 47 nền
kinh tế.

\textbf{(2) DAI bị giới hạn Tier-1 và Tier-2}: PICS3/Standardized chỉ có
website binary (Tier-1). BREADY có Tier-2 nhưng không đo Tier-3/4.
Chuyên đề chỉ phát biểu về phần bù Tier-1/2, không phải toàn bộ digital
adoption spectrum.

\textbf{(3) TCI composite với trọng số đồng đều}: Trung bình đơn giản 4
thành phần (ISO, R\&D, đổi mới sản phẩm, công nghệ nước ngoài) chưa được
xác nhận bằng factor analysis --- trọng số tối ưu có thể khác nhau giữa
các nhóm ICRV.

\textbf{(4) Manager items thiếu trong một số sóng}:
\texttt{exp\_manager}, \texttt{gender} không có trong tất cả 14 sóng
WBES. H4 CĐ2 ước lượng trên subsample, tạo selection bias tiềm tàng.

\textbf{(5) Không tách biệt COVID-19 effect khỏi digital adoption
effect}: BREADY 2018--2025 bao gồm COVID disruption (2020--2022) + AI
acceleration (2023+). Kết quả BREADY 2025 nên được giải thích như
"post-pandemic + digital acceleration combined".

\textbf{(6) SIDS coverage chưa đầy đủ}: 7/\textasciitilde52 SIDS theo
UNCTAD. Tuvalu, Nauru, Marshall Islands, Micronesia chưa có WBES. Kết
luận về forced penalty có thể chưa đại diện toàn bộ SIDS Pacific.

\textbf{(7) Không có dữ liệu sector-specific digital intensity}: WBES
không phân biệt DAI theo ngành. Doanh nghiệp sản xuất (MES) và dịch vụ
(CRM) có hệ số DAI Tier-1 giống nhau nhưng digital intensity thực tế
khác nhau. Industry FE kiểm soát một phần nhưng không giải quyết triệt
để.

\hypertarget{75-kux1ebf-houx1ea1ch-houxe0n-thiux1ec7n}{%
\subsubsection{7.5 Kế hoạch hoàn
thiện}\label{75-kux1ebf-houx1ea1ch-houxe0n-thiux1ec7n}}

Bốn giai đoạn từ tháng 6 đến tháng 9 năm 2026:

\textbf{Giai đoạn 1 --- Tháng 6/2026}: Chạy M0--M7 trên R với dữ liệu
WBES harmonized đầy đủ. Lind--Mehlum test. Paternoster z-test temporal.
Xuất 8 bảng kết quả (Bảng 4.3--4.10) định dạng APA 7th. Cập nhật Bảng
5.1 với số liệu thực từ regression. Nộp bản nháp CĐ1 hoàn thiện cho
PGS.TS. Phan Anh Tú.

\textbf{Giai đoạn 2 --- Tháng 7/2026}: Nộp P3 (APJM) và P4 (MIR). Chuẩn
bị dữ liệu WBES cho CĐ2 (101.185 obs, 47 nền kinh tế, 3 thế hệ schema
harmonized). Phác thảo M0--M7 specification CĐ2.

\textbf{Giai đoạn 3 --- Tháng 8/2026}: Ước lượng M0--M7 CĐ2 với HC1
robust SE + country/year FE. Kiểm định H1--H6. Chạy 8 robustness checks.
Viết Chương 4 (Kết quả) và Chương 5 (Thảo luận) CĐ2.

\textbf{Giai đoạn 4 --- Tháng 9/2026}: Tích hợp CĐ1 + CĐ2 + P3/P4/P5 vào
khung luận án tổng thể. Viết Chương 1 (Tổng quan) và Chương 6 (Kết luận)
luận án. Nộp bản thảo luận án hoàn chỉnh lần 1 cho hội đồng hướng dẫn.

\begin{center}\rule{0.5\linewidth}{0.5pt}\end{center}

\hypertarget{tuxe0i-liux1ec7u-tham-khux1ea3o}{%
\subsection{TÀI LIỆU THAM KHẢO}\label{tuxe0i-liux1ec7u-tham-khux1ea3o}}

\emph{Trích dẫn theo APA 7th edition. Danh mục đầy đủ tại}
\texttt{thesis/04\_references\_apa7.md} \emph{(v2.5).}

ADB. (2024). \emph{Asian development outlook 2024: Strengthening climate
resilience}. Asian Development Bank.

ADB. (2026, April). \emph{Asian development outlook April 2026: The
Middle East conflict challenges resilience in Asia and the Pacific}.
Asian Development Bank.

ADB. (2026). \emph{Vietnam economic outlook 2026: Navigating the
crosscurrents}. Asian Development Bank.

Aydemir, M. F., Ozsahin, M., \& Algur, E. (2024). The
internationalization-firm performance relationship: A meta-analysis.
\emph{International Journal of Management Economics and Business},
\emph{20}(2), 287--312.

Bausch, A., \& Krist, M. (2007). The effect of context-related
moderators on the internationalization--performance relationship:
Evidence from meta-analysis. \emph{Management International Review,
47}(3), 319--347.

Beblawi, H. (1987). The rentier state in the Arab world. In H. Beblawi
\& G. Luciani (Eds.), \emph{The rentier state} (pp. 49--62). Croom Helm.

Bertram, G. (2006). Introduction: The MIRAB model in the twenty-first
century. \emph{Asia Pacific Viewpoint, 47}(1), 1--13.

Briguglio, L. (1995). Small island developing states and their economic
vulnerabilities. \emph{World Development, 23}(9), 1615--1632.

Brynjolfsson, E., \& McAfee, A. (2014). \emph{The second machine age:
Work, progress, and prosperity in a time of brilliant technologies}. W.
W. Norton \& Company.

Coase, R. H. (1937). The nature of the firm. \emph{Economica, 4}(16),
386--405.

Cohen, W. M., \& Levinthal, D. A. (1990). Absorptive capacity: A new
perspective on learning and innovation. \emph{Administrative Science
Quarterly, 35}(1), 128--152.

Contractor, F. J., Kundu, S. K., \& Hsu, C. C. (2003). A three-stage
theory of international expansion. \emph{Journal of International
Business Studies, 34}(1), 5--18.

David, P. A. (1990). The dynamo and the computer: An historical
perspective on the modern productivity paradox. \emph{American Economic
Review, 80}(2), 355--361.

Tantawy, A., Brown, T. E., \& Eshima, Y. (2024). Institutional voids and
new venture performance: The moderating role of
founders\textquotesingle{} political ties. \emph{European Management
Review}, \emph{21}(2), 312--328.
\url{https://doi.org/10.1111/emre.12601}

Wagner, J. (2025). Productivity and international firm activities: What
do we know? \emph{Working Paper Series in Economics, University of
Lüneburg}, No. 425.

Đỗ, T. H., \& Phan, A. T. (2024). \emph{Meta-analysis of
internationalization--performance relationship in Asia-Pacific SMEs}.
Proceedings of ICBEF 2025. Business Innovation Academy.

Đỗ, T. H., \& Phan, A. T. (2026). Internationalization and firm
performance in Vietnam. \emph{Asia Pacific Journal of Management} (under
review).

Hall, P. A., \& Soskice, D. (Eds.). (2001). \emph{Varieties of
capitalism: The institutional foundations of comparative advantage}.
Oxford University Press.

Hertog, S. (2010). \emph{Princes, brokers, and bureaucrats: Oil and the
state in Saudi Arabia}. Cornell University Press.

Hsieh, C.-T., \& Klenow, P. J. (2009). Misallocation and manufacturing
TFP in China and India. \emph{Quarterly Journal of Economics, 124}(4),
1403--1448.

Hvidt, M. (2013). Economic diversification in GCC countries: Past record
and future trends. \emph{LSE Kuwait Programme Research Papers, 27}.

Johanson, J., \& Vahlne, J.-E. (1977). The internationalization process
of the firm. \emph{Journal of International Business Studies, 8}(1),
23--32.

Kafouros, M., Wang, C., Piperopoulos, P., \& Zhang, M. (2023). Academic
publishing in business and management. \emph{Global Strategy Journal,
13}(1), 3--28.

Kirca, A. H., Hult, G. T. M., Deligonul, S., Perryy, M. Z., \& Cavusgil,
S. T. (2012). A multilevel examination of the drivers of firm
multinationality. \emph{Journal of Management, 38}(2), 502--530.

Lu, J. W., \& Beamish, P. W. (2004). International diversification and
firm performance: The S-curve hypothesis. \emph{Academy of Management
Journal, 47}(4), 598--609.

Marano, V., Arregle, J.-L., Hitt, M. A., Spadafora, E., \& van Essen, M.
(2016). Home country institutions and the
internationalization-performance relationship. \emph{Journal of
Management, 42}(5), 1075--1110.

Mar, J., Đỗ, T. H., \& Phan, A. T. (2026). Digital adoption and
internationalization in Singapore. \emph{Management International
Review} (under review).

Meyer, K. E., Neck, H. M., \& Meeks, M. D. (2017). The
entrepreneurship-strategic management interface. In M. A. Hitt et al.
(Eds.), \emph{Strategic entrepreneurship} (pp. 19--44). Blackwell.

Sachs, J. D., \& Warner, A. M. (2001). The curse of natural resources.
\emph{European Economic Review, 45}(4--6), 827--838.

Sikdar, A., \& Mukhopadhyay, A. (2026). Technology spillovers and firm
performance in India. \emph{Asian Development Review, 43}(1), 37--75.

Stallkamp, M., \& Schotter, A. P. J. (2021). Platforms without borders?
\emph{Global Strategy Journal, 11}(1), 58--80.

UNCTAD. (2023). \emph{World investment report 2023}. United Nations
Conference on Trade and Development.

Verhoef, P. C., Broekhuizen, T., Bart, Y., Bhattacharya, A., Dong, J.
Q., Fabian, N., \& Haenlein, M. (2021). Digital transformation: A
multidisciplinary reflection and research agenda. \emph{Journal of
Business Research, 122}, 889--901.

Wang, J., Huang, Y., \& Hong, H. (2024). Technology concentration and
bank risk: Evidence from Chinese commercial banks. \emph{International
Review of Financial Analysis, 91}, 102968.

Williamson, O. E. (1985). \emph{The economic institutions of
capitalism}. Free Press.

World Bank. (2025). \emph{World Bank Enterprise Surveys 2025}. World
Bank Group.

Wu, J., Li, S., \& Selover, D. D. (2022). Foreign direct investment vs
foreign trade: Substitutes or complements? \emph{Journal of
International Business Studies, 53}(9), 1726--1743.

Xu, D. (2024). De jure versus de facto institutional distance.
\emph{Global Strategy Journal, 14}(2), 187--215.

Zaheer, S. (1995). Overcoming the liability of foreignness.
\emph{Academy of Management Journal, 38}(2), 341--363.

\begin{center}\rule{0.5\linewidth}{0.5pt}\end{center}

\hypertarget{phux1ee5-lux1ee5c}{%
\subsection{PHỤ LỤC}\label{phux1ee5-lux1ee5c}}

\hypertarget{phux1ee5-lux1ee5c-a--phux1ea1m-vi-thux1ef1c-tux1ebf-cux1ee7a-nhuxf3m-dux1eef-liux1ec7u-wbes}{%
\subsubsection{Phụ lục A --- Phạm vi thực tế của nhóm dữ liệu
WBES}\label{phux1ee5-lux1ee5c-a--phux1ea1m-vi-thux1ef1c-tux1ebf-cux1ee7a-nhuxf3m-dux1eef-liux1ec7u-wbes}}

\textbf{Bảng A.1}. \emph{Phân bổ 47 nền kinh tế theo nhóm ICRV ---
101.185 doanh nghiệp, 108 cặp quốc gia--năm.}

\begin{longtable}[]{@{}llll@{}}
\toprule\noalign{}
Nhóm ICRV & Nền kinh tế & n (xấp xỉ) & Sóng khảo sát \\
\midrule\noalign{}
\endhead
\bottomrule\noalign{}
\endlastfoot
\textbf{Nhóm I --- Advanced đổi mới sáng tạo} & Singapore, Hong Kong
SAR, Hàn Quốc, Đài Loan, Israel & \textasciitilde8.200 & 2009--2025 \\
\textbf{Nhóm II --- Advanced tài nguyên} & Saudi Arabia, Qatar, Kuwait,
Bahrain, Brunei & \textasciitilde4.100 & 2013--2025 \\
\textbf{Nhóm III --- Upper-middle chuyển đổi} & Trung Quốc, Malaysia,
Thái Lan, Kazakhstan, Armenia, Georgia & \textasciitilde12.300 &
2009--2023 \\
\textbf{Nhóm IV --- Emerging chế tạo} & Việt Nam, Indonesia,
Philippines, Ấn Độ, Sri Lanka, Jordan, Mông Cổ & \textasciitilde21.000 &
2009--2025 \\
\textbf{Nhóm V --- Frontier cận biên} & Bangladesh, Pakistan, Lào,
Campuchia, Myanmar, Nepal, Bhutan, Maldives, Uzbekistan, Tajikistan,
Kyrgyzstan, Turkmenistan, Afghanistan, Timor-Leste, Iraq, Lebanon, Yemen
& \textasciitilde53.200 & 2009--2025 \\
\textbf{Nhóm VI --- SIDS Pacific} & Fiji, Papua New Guinea, Solomon
Islands, Tonga, Vanuatu, Samoa, Kiribati & \textasciitilde2.385 &
2009--2025 \\
\textbf{Tổng} & \textbf{47 nền kinh tế} &
\textbf{\textasciitilde101.185} & \textbf{14 sóng} \\
\end{longtable}

\emph{Ghi chú:} n xấp xỉ do harmonization loại quan sát thiếu biến phụ
thuộc. Phân bổ chính xác tại Bảng 4.1.

\textbf{Ba thế hệ schema WBES}:

\begin{longtable}[]{@{}llll@{}}
\toprule\noalign{}
Schema & Năm & n & DAI coverage \\
\midrule\noalign{}
\endhead
\bottomrule\noalign{}
\endlastfoot
PICS3 & 2009--2012 & 14.171 & Tier-1 only (website binary) \\
Standardized & 2013--2017 & 24.564 & Tier-1 only \\
BREADY/BEE & 2018--2025 & 62.450 & Tier-1 + Tier-2 (e-payment
k33/k38) \\
\end{longtable}

Triple-defense system (Mục 4.11) đảm bảo kết quả từ BREADY 2025 không bị
nhiễu schema effect.

\hypertarget{phux1ee5-lux1ee5c-bg}{%
\subsubsection{Phụ lục B--G}\label{phux1ee5-lux1ee5c-bg}}

Sẽ mở rộng trong giai đoạn hoàn thiện tháng 6--8/2026:

\begin{itemize}
\tightlist
\item
  B: Quy trình harmonization ba thế hệ schema chi tiết
\item
  C: Bảng phân loại 47 nước vào ICRV 6 nhóm (đầy đủ)
\item
  D: Kết quả M0--M7 đầy đủ (khi phân tích hoàn tất)
\item
  E: Kiểm định robustness 8 loại (kết quả chi tiết)
\item
  F: PRISMA-style literature search protocol
\item
  G: Biểu đồ turning point theo nhóm ICRV
\end{itemize}

\begin{center}\rule{0.5\linewidth}{0.5pt}\end{center}

\emph{Chuyên đề tiến sĩ số 1 --- Phiên bản nộp v1.0 (12/05/2026)}
\emph{Nghiên cứu sinh: Đỗ Thùy Hương --- Trường Đại học Kỹ thuật Công
nghệ Vĩnh Long (VLUTE)} \emph{Người hướng dẫn: PGS.TS. Phan Anh Tú ---
Trường Kinh tế, Đại học Cần Thơ (CTU)}

\end{document}
